% Leçon 38 — Expression écrite : origine du caractère 电 et radicaux ; caractères des réseaux sociaux — Chinois 1ère A — Cours signé OnBuch+
% Programme officiel MINESEC. Compiler avec : tectonic ZHA38-ecrit-origine-du-caractere-et-radicaux-caracteres.tex
\newcommand{\DOCMATIERE}{Chinois}
\newcommand{\DOCNIVEAU}{1ère A}
\newcommand{\DOCMODULE}{Module 5 : Mass médias et télécommunication}
\newcommand{\DOCLECON}{ZHA38}
\newcommand{\DOCTITRE}{Leçon 38 — Expression écrite : origine du caractère 电 et radicaux ; caractères des réseaux sociaux}
% OnBuch+ — préambule « cours signés » ENRICHI (compilé avec tectonic / XeLaTeX)
% Système visuel « Pop » d'OnBuch+ : fond crème (jamais blanc), bordures encre
% épaisses, ombres « dures » décalées, titres Archivo Black en capitales.
% Version MATHÉMATIQUES : graphes (pgfplots/tikz), boîtes pédagogiques
% (définition, propriété, méthode, attention, le savais-tu, exercice résolu,
% exercice, TP, bilan), page de garde et signature OnBuch+ ancrée en bas.
%
% Macros attendues dans main.tex AVANT \input{preamble} :
%   \DOCMATIERE  \DOCNIVEAU  \DOCMODULE  \DOCLECON (numéro)  \DOCTITRE
\documentclass[11pt]{article}

\usepackage{fontspec}
\usepackage[french]{babel} % français = langue principale ; le chinois via \zh{..} / textebox (xeCJK)
\usepackage{xeCJK}
\usepackage{pifont} % \ding{51} \ding{55} utilisés par les modèles
\usepackage[a4paper,margin=2cm,top=3.4cm,bottom=2.8cm,headheight=1.4cm,footskip=1.3cm]{geometry}
\usepackage{amsmath,amssymb,amsfonts}
\usepackage{mathtools} % \xmapsto, \coloneqq... souvent utilisés par les modèles
\usepackage{cancel} % simplifications barrées (bilans de forces)
% \abs{z} (module, valeur absolue) et \norm{u} (norme), utilisés par les modèles
\providecommand{\abs}{}\let\abs\relax\DeclarePairedDelimiter{\abs}{\lvert}{\rvert}
\providecommand{\norm}{}\let\norm\relax\DeclarePairedDelimiter{\norm}{\lVert}{\rVert}
\usepackage[table]{xcolor} % [table] : fournit aussi \rowcolors (alternance de lignes)
\usepackage{pagecolor}
\usepackage{tikz}
\usetikzlibrary{arrows.meta,positioning,calc,shapes.geometric,shapes.misc,shapes.arrows,decorations.pathmorphing,decorations.pathreplacing,decorations.markings,patterns,shadows,mindmap,trees,fit,backgrounds,matrix,intersections,angles,quotes}
\usepackage{pgfplots}
\usepackage[version=4]{mhchem} % équations nucléaires \ce{^{235}_{92}U}
\usepackage[european,siunitx]{circuitikz} % schémas électriques
\pgfplotsset{compat=1.18}
\usepgfplotslibrary{fillbetween,groupplots} % aires entre courbes (calcul intégral)
\usepackage{enumitem}
\usepackage{fancyhdr}
% [most] charge toutes les bibliothèques tcolorbox (listings, minted, raster,
% documentation...) et gonfle inutilement la mémoire TeX sur un document long
% (~300 boîtes) : on ne charge que ce qui est réellement utilisé.
\usepackage[skins,breakable]{tcolorbox}
\usepackage{graphicx}
\usepackage{colortbl}
\usepackage{booktabs}
\usepackage{tabularx}
\usepackage{diagbox} % case d'en-tête diagonale des tableaux à double entrée
% Colonnes extensibles centrée / gauche / droite (C, L, R), souvent utilisées par les modèles
\newcolumntype{C}{>{\centering\arraybackslash}X}
\newcolumntype{L}{>{\raggedright\arraybackslash}X}
\newcolumntype{R}{>{\raggedleft\arraybackslash}X}
\usepackage{array}
\usepackage{multicol}
\usepackage{multirow}
\usepackage{siunitx}
% parse-numbers=false : accepte aussi \SI{2,5 \times 10^{-3}}{...}
\sisetup{locale=FR,output-decimal-marker={,},group-minimum-digits=4,parse-numbers=false}
\DeclareSIUnit\ohm{\ensuremath{\symup{\Omega}}} % U+2126 absent de la police
\DeclareSIUnit\division{div} % réglages d'oscilloscope (ms/div, V/div)
\DeclareSIUnit\tr{tr} % tours (vitesse de rotation, ex. \SI{1500}{\tr\per\minute})
\DeclareSIUnit\molar{M} % concentration molaire (ex. \SI{3}{\molar}, \SI{1}{\micro\molar})
\DeclareSIUnit\an{an} % année (le modèle confond parfois avec \year, un compteur TeX) — ex. \SI{5}{\kilo\meter\per\an}
\DeclareSIUnit\FCFA{FCFA} % franc CFA (ex. \SI{500}{\FCFA\per\kilo\gram})
\DeclareSIUnit\degreeBrix{\textdegree Brix} % teneur en sucre d'un sirop/jus (réfractométrie)
\DeclareSIUnit\month{mois} % le modèle utilise parfois \month, un compteur TeX, comme unité de durée
\usepackage{qrcode}
% \degree : commande fréquemment utilisée par les modèles pour un angle ou une
% température en dehors de \SI{}{} (non fournie par les paquets chargés).
\providecommand{\degree}{^{\circ}}
% \qed (fin de démonstration), utilisé par les modèles sans amsthm
\providecommand{\qed}{\hfill$\square$}
% \barre : double barre des valeurs interdites dans un tableau de variations (tkz-tab)
\providecommand{\barre}{\ensuremath{\|}}
% environnement proof (amsthm non chargé) utilisé par les modèles
\ifdefined\proof\else\newenvironment{proof}[1][Démonstration]{\par\noindent\textit{#1.}\ }{\qed\par}\fi
\usepackage{lastpage}
\usepackage{hyperref}
\hypersetup{hidelinks}

% ============================================================
% Palette « Pop » OnBuch+ — mêmes hex que la classe P (plus_ui.dart)
% ============================================================
\definecolor{poporange}{HTML}{F59321}
\definecolor{popdark}{HTML}{E07A0C}
\definecolor{popcream}{HTML}{FFF6E8}
\definecolor{popink}{HTML}{1C1714}
\definecolor{poppurple}{HTML}{7A5AE0}
\definecolor{popgreen}{HTML}{2E9E6A}
\definecolor{poppink}{HTML}{E0457B}
\definecolor{popblue}{HTML}{2F7DE1}
\definecolor{popgold}{HTML}{C98A12}
\definecolor{popmuted}{HTML}{8A7F76}
\definecolor{popline}{HTML}{EADFD2}
\definecolor{popgray}{HTML}{8A7F76}
\colorlet{popgrey}{popgray}
% Alias tolérants (le modèle nomme parfois les couleurs différemment)
\colorlet{popwhite}{white}
\colorlet{popblack}{popink}
\colorlet{popred}{poppink}
\colorlet{popyellow}{popgold}
% Teintes claires (fonds de tableaux/encadrés) — le modèle nomme parfois
% les couleurs avec un suffixe L (ex. popblueL) sans les avoir définies.
\colorlet{poporangeL}{poporange!20}
\colorlet{popdarkL}{popdark!20}
\colorlet{poppurpleL}{poppurple!20}
\colorlet{popgreenL}{popgreen!20}
\colorlet{poppinkL}{poppink!20}
\colorlet{popblueL}{popblue!20}
\colorlet{popgoldL}{popgold!20}
\colorlet{popredL}{popred!20}
\colorlet{popgrayL}{popgray!20}
% Teintes claires pour fonds de tableaux / zones
\colorlet{popblueL}{popblue!10!white}
\colorlet{poporangeL}{poporange!14!white}
\colorlet{popgreenL}{popgreen!12!white}
\colorlet{poppurpleL}{poppurple!12!white}
\colorlet{poppinkL}{poppink!12!white}
\colorlet{popgoldL}{popgold!14!white}

\pagecolor{popcream}

% ============================================================
% Typographie
% ============================================================
\defaultfontfeatures{Ligatures=TeX}
\setmainfont{PlusJakartaSans-Medium.ttf}[Path=../../../fonts/,BoldFont=PlusJakartaSans-Bold.ttf]
% Symboles, API et emojis absents de la police principale : repli sur DejaVu Sans (caractères actifs)
\newfontface\symfont{DejaVuSans.ttf}[Path=../../../fonts/]
\makeatletter
\newcommand\symactive[1]{\catcode#1=13 \begingroup\lccode`\~=#1 \lowercase{\endgroup\def~}{{\symfont\char#1}}}
\newcommand\symblank[1]{\catcode#1=13 \begingroup\lccode`\~=#1 \lowercase{\endgroup\def~}{}}
\newcommand\symhyph[1]{\catcode#1=13 \begingroup\lccode`\~=#1 \lowercase{\endgroup\def~}{\mbox{-}}}
\newcount\symc
\newcommand\symrange[2]{\symc=#1 \@whilenum\symc<#2 \do{\expandafter\symactive\expandafter{\the\symc}\advance\symc by 1 }}
\newcommand\symblankrange[2]{\symc=#1 \@whilenum\symc<#2 \do{\expandafter\symblank\expandafter{\the\symc}\advance\symc by 1 }}
\makeatother
\symrange{"0250}{"02FF}\symactive{"2713}\symactive{"2714}\symactive{"2717}\symactive{"2718}\symactive{"2610}\symactive{"2611}\symactive{"2612}
\symrange{"2600}{"27BF}
\catcode"2705=13 \begingroup\lccode`\~="2705 \lowercase{\endgroup\def~}{{\symfont\char"2713}}\symblank{"26BD}\symblank{"26A1}
\symrange{"2460}{"257F}\symactive{"223C}\symactive{"2248}\symactive{"2260}
\symactive{"01F8}\symactive{"01F9}\symactive{"1E3E}\symactive{"1E3F}
\symhyph{"2011}\symhyph{"2E1A}\symblankrange{"1F300}{"1FAFF}\symblank{"FE0F}
% Caractères chinois : WenQuanYi Zen Hei (sous-ensemble GB2312 + pinyin), gras/italique simulés
\setCJKmainfont{WenQuanYiZenHei-subset.ttf}[Path=../../../fonts/,AutoFakeBold=3,AutoFakeSlant=0.2]
\xeCJKsetup{CJKspace=false}
% Pinyin : voyelles à caron (ǎ ǐ ǒ ǔ ǖ ǘ ǚ ǜ) absentes de la police principale -> caractères actifs
\newfontface\pyfont{WenQuanYiZenHei-subset.ttf}[Path=../../../fonts/]
\newcommand\pyactive[1]{\catcode#1=13 \begingroup\lccode`\~=#1 \lowercase{\endgroup\def~}{{\pyfont\char#1}}}
\pyactive{"01CD}\pyactive{"01CE}\pyactive{"01CF}\pyactive{"01D0}\pyactive{"01D1}\pyactive{"01D2}\pyactive{"01D3}\pyactive{"01D4}\pyactive{"01D5}\pyactive{"01D6}\pyactive{"01D7}\pyactive{"01D8}\pyactive{"01D9}\pyactive{"01DA}\pyactive{"01DB}\pyactive{"01DC}
% Maths en sans-serif (Fira Math) avec chiffres et lettres droites de Plus
% Jakarta, pour que les formules se fondent dans le texte.
\usepackage[math-style=ISO,bold-style=ISO]{unicode-math}
\setmathfont{FiraMath-Regular.otf}[Path=../../../fonts/]
\setmathfont{PlusJakartaSans-Medium.ttf}[Path=../../../fonts/,range={up/num,up/{Latin,latin},\mathup}]
\usepackage{icomma} % virgule décimale sans espace en mode math
% Caractères mathématiques Unicode absents des polices de texte (Plus Jakarta,
% Archivo Black) : rendus par leur équivalent mathématique, en texte comme en
% formule — sinon ils s'affichent en carré vide (ex. « Arithmétique dans ℤ »).
\usepackage{newunicodechar}
\newunicodechar{ℝ}{\ensuremath{\mathbb{R}}}
\newunicodechar{ℤ}{\ensuremath{\mathbb{Z}}}
\newunicodechar{ℂ}{\ensuremath{\mathbb{C}}}
\newunicodechar{ℕ}{\ensuremath{\mathbb{N}}}
\newunicodechar{ℚ}{\ensuremath{\mathbb{Q}}}
\newunicodechar{𝔻}{\ensuremath{\mathbb{D}}}
\newunicodechar{α}{\ensuremath{\alpha}}
\newunicodechar{β}{\ensuremath{\beta}}
\newunicodechar{γ}{\ensuremath{\gamma}}
\newunicodechar{δ}{\ensuremath{\delta}}
\newunicodechar{ε}{\ensuremath{\varepsilon}}
\newunicodechar{θ}{\ensuremath{\theta}}
\newunicodechar{λ}{\ensuremath{\lambda}}
\newunicodechar{μ}{\ensuremath{\mu}}
\newunicodechar{σ}{\ensuremath{\sigma}}
\newunicodechar{τ}{\ensuremath{\tau}}
\newunicodechar{φ}{\ensuremath{\varphi}}
\newunicodechar{ψ}{\ensuremath{\psi}}
\newunicodechar{ω}{\ensuremath{\omega}}
\newunicodechar{Σ}{\ensuremath{\Sigma}}
% majuscules produites par \MakeUppercase dans les titres (α-aminés -> Α)
\newunicodechar{Α}{\ensuremath{\alpha}}
\newunicodechar{Β}{\ensuremath{\beta}}
\newunicodechar{Γ}{\ensuremath{\Gamma}}
\newunicodechar{Δ}{\ensuremath{\Delta}}
\newunicodechar{Ε}{\ensuremath{\varepsilon}}
\newunicodechar{Θ}{\ensuremath{\Theta}}
\newunicodechar{Λ}{\ensuremath{\Lambda}}
\newunicodechar{Μ}{\ensuremath{\mu}}
\newunicodechar{Τ}{\ensuremath{\tau}}
\newunicodechar{Φ}{\ensuremath{\Phi}}
\newunicodechar{Ψ}{\ensuremath{\Psi}}
\newunicodechar{Ω}{\ensuremath{\Omega}}
\newunicodechar{↦}{\ensuremath{\mapsto}}
\newunicodechar{∈}{\ensuremath{\in}}
\newunicodechar{∉}{\ensuremath{\notin}}
\newunicodechar{∧}{\ensuremath{\wedge}}
\newunicodechar{⇔}{\ensuremath{\Leftrightarrow}}
\newunicodechar{⟺}{\ensuremath{\Longleftrightarrow}}
\newunicodechar{⇒}{\ensuremath{\Rightarrow}}
\newunicodechar{⟹}{\ensuremath{\Longrightarrow}}
\newunicodechar{∘}{\ensuremath{\circ}}
\newunicodechar{∪}{\ensuremath{\cup}}
\newunicodechar{∩}{\ensuremath{\cap}}
\newunicodechar{≡}{\ensuremath{\equiv}}
\newunicodechar{⊂}{\ensuremath{\subset}}
\newunicodechar{⊥}{\ensuremath{\perp}}
\newunicodechar{∃}{\ensuremath{\exists}}
\newunicodechar{∀}{\ensuremath{\forall}}
\newunicodechar{∛}{\ensuremath{\sqrt[3]{\,}}}
\newunicodechar{∜}{\ensuremath{\sqrt[4]{\,}}}
\newunicodechar{⃗}{\ensuremath{{}^{\to}}}
\newunicodechar{ⁿ}{\ifmmode^{n}\else\textsuperscript{\ensuremath{n}}\fi}
\newunicodechar{ˣ}{\ifmmode^{x}\else\textsuperscript{\ensuremath{x}}\fi}
\newunicodechar{ᵇ}{\ifmmode^{b}\else\textsuperscript{\ensuremath{b}}\fi}
\newunicodechar{ʳ}{\ifmmode^{r}\else\textsuperscript{\ensuremath{r}}\fi}
\newunicodechar{ᵘ}{\ifmmode^{u}\else\textsuperscript{\ensuremath{u}}\fi}
\newunicodechar{ʸ}{\ifmmode^{y}\else\textsuperscript{\ensuremath{y}}\fi}
\newunicodechar{ᵃ}{\ifmmode^{a}\else\textsuperscript{\ensuremath{a}}\fi}
\newunicodechar{ᵉ}{\ifmmode^{e}\else\textsuperscript{\ensuremath{e}}\fi}
\newunicodechar{ᵏ}{\ifmmode^{k}\else\textsuperscript{\ensuremath{k}}\fi}
\newunicodechar{ᵖ}{\ifmmode^{p}\else\textsuperscript{\ensuremath{p}}\fi}
\newunicodechar{ᴺ}{\ifmmode^{N}\else\textsuperscript{\ensuremath{N}}\fi}
\newunicodechar{ᶜ}{\ifmmode^{c}\else\textsuperscript{\ensuremath{c}}\fi}
\newunicodechar{ʰ}{\ifmmode^{h}\else\textsuperscript{\ensuremath{h}}\fi}
\newunicodechar{ᵠ}{\ifmmode^{q}\else\textsuperscript{\ensuremath{q}}\fi}
\newunicodechar{ₙ}{\ifmmode_{n}\else\textsubscript{\ensuremath{n}}\fi}
\newunicodechar{ₐ}{\ifmmode_{a}\else\textsubscript{\ensuremath{a}}\fi}
\newunicodechar{ᵢ}{\ifmmode_{i}\else\textsubscript{\ensuremath{i}}\fi}
\newunicodechar{ᵣ}{\ifmmode_{r}\else\textsubscript{\ensuremath{r}}\fi}
\newunicodechar{ₖ}{\ifmmode_{k}\else\textsubscript{\ensuremath{k}}\fi}
\newunicodechar{ᵦ}{\ifmmode_{\beta}\else\textsubscript{\ensuremath{\beta}}\fi}
\newunicodechar{ₓ}{\ifmmode_{x}\else\textsubscript{\ensuremath{x}}\fi}
\newfontfamily\popdisplay{ArchivoBlack-Regular.ttf}[Path=../../../fonts/]
\newfontfamily\poplogo{SpaceGrotesk-Bold.ttf}[Path=../../../fonts/]
\newfontfamily\popsemi{PlusJakartaSans-SemiBold.ttf}[Path=../../../fonts/]
\newfontfamily\popxbold{PlusJakartaSans-ExtraBold.ttf}[Path=../../../fonts/]
\newfontfamily\popxboldls{PlusJakartaSans-ExtraBold.ttf}[Path=../../../fonts/,LetterSpace=3.0]

\setlist[enumerate]{leftmargin=1.7em,itemsep=5pt,topsep=4pt}
\setlist[itemize]{leftmargin=1.7em,itemsep=4pt,topsep=4pt,label={\color{poporange}\textbullet}}
\setlength{\parindent}{0pt}
\setlength{\parskip}{5pt}
\renewcommand{\arraystretch}{1.25}

% pgfplots : style Pop par défaut
\pgfplotsset{
  every axis/.append style={axis line style={popink,line width=1pt},tick style={popink},
    grid style={popline,line width=0.5pt},label style={font=\small},tick label style={font=\footnotesize}},
  popcurve/.style={poporange,line width=1.8pt,no markers,smooth},
}
% Même style disponible dans un \draw TikZ simple (hors pgfplots)
\tikzset{popcurve/.style={poporange,line width=1.8pt}}
\tikzset{popgrid/.style={popline,very thin}}
\colorlet{popcurve}{poporange} % le nom du style est parfois utilisé comme couleur

% ============================================================
% Pill / squiggle
% ============================================================
\newcommand{\poppill}[3]{%
  \tikz[baseline=-0.55ex]\node[draw=popink,line width=1.2pt,rounded corners=2.6pt,%
    fill=#2,inner xsep=6.5pt,inner ysep=3pt]{%
    \popxboldls\fontsize{7}{7}\selectfont\color{#3}\MakeUppercase{#1}};%
}
\newcommand{\popsquiggle}[1]{%
  \tikz[baseline=-0.4ex]{
    \draw[#1,line width=2.6pt,line cap=round]
      plot [smooth] coordinates {(0,0.09) (0.34,0.01) (0.68,0.21) (1.02,0.01) (1.36,0.10)};
  }%
}
% Mot-clé surligné (marqueur orange sous le texte)
\newcommand{\cle}[1]{{\popxbold\color{popdark}#1}}

% ============================================================
% En-tête / pied
% ============================================================
\pagestyle{fancy}
\fancyhf{}
\renewcommand{\headrulewidth}{0pt}
\renewcommand{\footrulewidth}{0pt}
\lhead{\raisebox{-0.15ex}{\poplogo\fontsize{13}{13}\selectfont\color{popink}OnBuch\color{poporange}+}}
\rhead{\poppill{\DOCMATIERE\ \textbullet\ \DOCNIVEAU\ \textbullet\ Leçon \DOCLECON}{white}{popink}}
\lfoot{\popxbold\fontsize{7.6}{7.6}\selectfont\color{popmuted}\MakeUppercase{Cours signé OnBuch+}\ \textbullet\ {\popsemi\DOCTITRE}}
\rfoot{\tikz[baseline=-0.6ex]\node[rounded corners=8.5pt,fill=popink,minimum height=17pt,minimum width=17pt,inner xsep=5pt,inner sep=0]{\popxbold\fontsize{8}{8}\selectfont\color{popcream}\thepage\,/\,\pageref*{LastPage}};}

% ============================================================
% Titres
% ============================================================
\newcommand{\chaptitre}[2]{%
  \begin{tcolorbox}[enhanced,colback=poporange,colframe=popink,boxrule=2.6pt,arc=11pt,
      drop shadow=popink,left=15pt,right=15pt,top=12pt,bottom=11pt,before skip=3pt,after skip=18pt]
    \poppill{#2}{popink}{popcream}\par\vspace{9pt}
    {\raggedright\hyphenpenalty=10000\exhyphenpenalty=10000\popdisplay\fontsize{22}{24}\selectfont\color{popcream}\MakeUppercase{#1}\par}
  \end{tcolorbox}
}
% Section numérotée : pastille orange avec le numéro + titre Archivo
\newcounter{popsec}
\newcommand{\coursec}[1]{%
  \stepcounter{popsec}\setcounter{popsub}{0}%
  \par\vspace{16pt}\noindent
  \begin{minipage}{\linewidth}
  \tikz[baseline=-0.7ex]\node[draw=popink,line width=1.6pt,fill=poporange,rounded corners=4pt,minimum size=22pt,inner sep=0]{\popdisplay\fontsize{12}{12}\selectfont\color{popcream}\thepopsec};%
  \hspace{8pt}{\popdisplay\fontsize{15}{17}\selectfont\color{popink}\MakeUppercase{#1}}\par
  \vspace{4pt}\noindent{\color{poporange}\rule{36pt}{3.6pt}}
  \end{minipage}\par\nopagebreak\vspace{8pt}
}
\newcounter{popsub}
\newcommand{\courssub}[1]{%
  \stepcounter{popsub}%
  \par\vspace{9pt}\noindent{\popxbold\fontsize{11.5}{13}\selectfont\color{popink}{\color{poporange}\thepopsec.\thepopsub}\hspace{6pt}#1}\par\nopagebreak\vspace{4pt}
}

% ============================================================
% Boîtes pédagogiques — même recette : fond blanc, bordure épaisse colorée,
% ombre dure encre, badge plein. Titre optionnel [#1].
% ============================================================
\tcbset{popbase/.style={breakable,enhanced,colback=white,boxrule=2.3pt,arc=9pt,
  shadow={3pt}{-3pt}{0pt}{popink},left=14pt,right=14pt,top=11pt,bottom=11pt,before skip=11pt,after skip=15pt,
  fontupper=\popsemi\fontsize{10.6}{14.5}\selectfont\color{popink},
  fontlower=\popsemi\fontsize{10.6}{14.5}\selectfont\color{popink}}}
% Coche dessinée (le glyphe ✓ n'existe pas dans les polices du document)
\newcommand{\popcheck}[1]{\tikz[baseline=-0.5ex]\draw[#1,line width=1.6pt,line cap=round,line join=round](-0.13,0.0)--(-0.04,-0.09)--(0.13,0.1);}
\providecommand{\coche}{\popcheck{popgreen}}
\providecommand{\resultat}[1]{\par\smallskip\noindent\fcolorbox{popgreen}{popgreenL}{\parbox{\dimexpr\linewidth-2\fboxsep-2\fboxrule\relax}{\textbf{Résultat.} #1}}\par\smallskip}
\newcommand{\popboxhead}[3]{\poppill{#1}{#2}{white}\ifx\relax#3\relax\else\hspace{6pt}{\popxbold\fontsize{10.6}{12}\selectfont\color{popink}#3}\fi\par\vspace{7pt}}

\newtcolorbox{aretenir}[1][]{popbase,colframe=popgreen,before upper={\popboxhead{À retenir}{popgreen}{#1}}}
% Texte en chinois (document de lecture, dialogue, extrait) : cadre
% d'étude. Un titre optionnel (source, auteur) s'affiche dans l'en-tête.
\newtcolorbox{textebox}[1][]{popbase,colframe=popblue,before upper={\popboxhead{文本 · Texte}{popblue}{#1}},after upper={}}
% Marques ✓ / ✗ (forme correcte / fautive) souvent utilisées par les modèles
\providecommand{\cross}{\textcolor{poppink}{\ensuremath{\times}}}
\providecommand{\xmark}{\cross}\providecommand{\crossmark}{\cross}\providecommand{\cmark}{\ensuremath{\checkmark}}
% Ligne de réponse / justification à compléter (inventée par les modèles)
\providecommand{\justification}[1]{\par\noindent\textit{Justification~:} \dotfill\par}
% \textipa (package tipa non chargé) : la police ne couvre pas l'API -> simple passage
\providecommand{\textipa}[1]{#1}
% Soulignement (ulem non chargé) : \uline, \ul
\providecommand{\uline}[1]{\underline{#1}}\providecommand{\ul}[1]{\underline{#1}}
% \glossaire{\item ...} inventée par les modèles : liste de glossaire
\providecommand{\glossaire}[1]{\begin{itemize}#1\end{itemize}}
% \alert (beamer) inventée par les modèles : texte mis en évidence
\providecommand{\alert}[1]{\textbf{\textcolor{poppink}{#1}}}
% variantes ASCII d'ordinaux écrites en macro
\providecommand{\iere}{\textsuperscript{re}}\providecommand{\ieme}{\textsuperscript{e}}\providecommand{\ere}{\textsuperscript{re}}\providecommand{\eme}{\textsuperscript{e}}\providecommand{\er}{\textsuperscript{er}}\providecommand{\ie}{\textsuperscript{e}}
% Ordinaux français écrits en macro par les modèles : 1\ière, 3\ième
\newcommand{\ière}{\textsuperscript{re}}\newcommand{\ième}{\textsuperscript{e}}
% \exemple (inventée par les modèles) : étiquette « Exemple : »
\providecommand{\exemple}{\textbf{Exemple~:}\ }
% \square utilisable aussi hors mode math (cases à cocher)
\AtBeginDocument{\let\squareMath\square\renewcommand{\square}{\ensuremath{\squareMath}}}
% Mot ou phrase en chinois à l'intérieur d'un paragraphe français
\newcommand{\zh}[1]{\ifmmode\text{#1}\else{#1}\fi}
\let\eng\zh \let\esp\zh \let\all\zh % alias tolérés
% flèches d'intonation inventées par les modèles
\providecommand{\textsearrow}{}\renewcommand{\textsearrow}{\ensuremath{\searrow}}\providecommand{\textnearrow}{}\renewcommand{\textnearrow}{\ensuremath{\nearrow}}
% \fr{..} : mot français (inventé par les modèles)
\providecommand{\fr}[1]{\textit{#1}}
% zhbox : encadré de texte chinois inventé par les modèles
\newenvironment{zhbox}{\begin{quote}}{\end{quote}}
% \courssubsub : titre de niveau 3 inventé par les modèles
\providecommand{\courssubsub}[1]{\par\medskip\noindent\textbf{#1}\par\smallskip}
% \zhbf : texte chinois en gras inventé par les modèles
\providecommand{\zhbf}[1]{\textbf{#1}}
% \vs : « versus » inventé par les modèles
\providecommand{\vs}{\textit{vs}\ }
% \blank : case à compléter inventée par les modèles
\providecommand{\blank}[1][]{\underline{\hspace{1.6em}}}
\newtcolorbox{exemplebox}[1][]{popbase,colframe=poporange,before upper={\popboxhead{Exemple}{poporange}{#1}}}
\newtcolorbox{definition}[1][]{popbase,colframe=popblue,before upper={\popboxhead{Définition}{popblue}{#1}}}
\newtcolorbox{propriete}[1][]{popbase,colframe=poppurple,before upper={\popboxhead{Propriété}{poppurple}{#1}}}
\newtcolorbox{methode}[1][]{popbase,colframe=popgold,before upper={\popboxhead{Méthode}{popgold}{#1}}}
\newtcolorbox{attention}[1][]{popbase,colframe=poppink,before upper={\popboxhead{Attention}{poppink}{#1}}}
\newtcolorbox{experience}[1][]{popbase,colframe=popblue,colback=popblueL,before upper={\popboxhead{Expérience}{popblue}{#1}}}
% « Le savais-tu ? » : carte violette pleine (contexte, culture, Cameroun)
\newtcolorbox{savaistu}[1][]{popbase,colback=poppurple,colframe=popink,
  fontupper=\popsemi\fontsize{10.4}{14.2}\selectfont\color{white},
  before upper={\poppill{Le savais-tu\,?}{white}{poppurple}\ifx\relax#1\relax\else\hspace{6pt}{\popxbold\color{white}#1}\fi\par\vspace{7pt}}}
% Exercice résolu : énoncé en haut, \tcblower puis la solution sur fond crème
\newcounter{popexo}\newcounter{popexr}
\newtcolorbox{exoresolu}[1][]{popbase,colframe=poporange,colbacklower=popcream,
  segmentation style={popink,line width=1.2pt,dashed},
  before upper={\stepcounter{popexr}\popboxhead{Exercice résolu \thepopexr}{poporange}{#1}},
  before lower={\poppill{Solution}{popink}{popcream}\par\vspace{6pt}}}
% Exercice (non résolu en place) — la correction est dans la partie Corrigés
\newtcolorbox{exercice}[1][]{popbase,colframe=popink,
  before upper={\stepcounter{popexo}\popboxhead{Exercice \thepopexo}{popink}{#1}}}
% Corrigé
\newtcolorbox{corrige}[1][]{popbase,colframe=popgreen,colback=popgreenL,
  before upper={\popboxhead{Corrigé}{popgreen}{#1}}}
% Objectifs / prérequis / bilan
\newtcolorbox{objectifs}{popbase,colframe=popink,colback=poporangeL,
  before upper={\popboxhead{Objectifs de la leçon}{popink}{}}}
\newtcolorbox{prerequis}{popbase,colframe=popink,colback=popblueL,
  before upper={\popboxhead{Prérequis}{popblue}{}}}
\newtcolorbox{bilan}{popbase,colframe=popink,colback=white,boxrule=2.8pt,
  before upper={\popboxhead{Fiche bilan}{poporange}{L'essentiel à connaître pour l'examen}}}
% Figure encadrée avec légende
\newtcolorbox{popfigure}[1][]{enhanced,breakable=false,colback=white,colframe=popline,boxrule=1.4pt,arc=8pt,
  left=10pt,right=10pt,top=10pt,bottom=8pt,before skip=10pt,after skip=12pt,halign=center,
  lower separated=false,halign lower=center,
  fontlower=\popsemi\fontsize{9}{11.5}\selectfont\color{popmuted},
  #1}
\newcommand{\legende}[1]{\par\vspace{6pt}{\popsemi\fontsize{9}{11.5}\selectfont\color{popmuted}#1}}

% ============================================================
% Signature OnBuch+
% ============================================================
\newcommand{\poplogoinline}[1]{{\poplogo\fontsize{#1}{#1}\selectfont\color{popink}OnBuch\color{poporange}+}}
\newcommand{\signatureonbuch}{%
  \begin{tcolorbox}[enhanced,colback=popink,colframe=popink,boxrule=2.2pt,arc=10pt,
      drop shadow=poporange,left=14pt,right=14pt,top=10pt,bottom=10pt,before skip=0pt,after skip=0pt]
    \tikz[baseline=-0.7ex]\node[circle,fill=popgreen,draw=popcream,line width=1.3pt,minimum size=22pt,inner sep=0]{\popcheck{white}};%
    \hspace{9pt}\begin{minipage}[c]{0.62\linewidth}
      {\popxbold\fontsize{12}{13}\selectfont\color{popcream}Signé\ \textbullet\ {\poplogo OnBuch\color{poporange}+}}\par\vspace{2pt}
      {\raggedright\popsemi\fontsize{8}{10.5}\selectfont\color{popline}\MakeUppercase{Équipe pédagogique OnBuch+}\\\MakeUppercase{Conforme au programme officiel MINESEC}\par}
    \end{minipage}\hfill
    {\poplogo\fontsize{22}{22}\selectfont\color{popcream}OnBuch\color{poporange}+}
  \end{tcolorbox}%
}

% ============================================================
% Page de garde
% #1 titre · #2 matière · #3 niveau · #4 module · #5 n° leçon · #6 durée ·
% #7 description / objectifs (texte ou itemize)
% ============================================================
\newcommand{\pagegarde}[7]{%
  \thispagestyle{empty}
  \newgeometry{a4paper,margin=1.7cm,top=1.6cm,bottom=1.4cm}%
  \begin{tikzpicture}[remember picture,overlay]
    \fill[poporange!18] ([xshift=-3.2cm,yshift=-3.6cm]current page.north east) circle (4.6cm);
    \fill[poppurple!14] ([xshift=2.4cm,yshift=7.6cm]current page.south west) circle (3.8cm);
  \end{tikzpicture}%
  {\poplogo\fontsize{30}{30}\selectfont\color{popink}OnBuch\color{poporange}+}\hfill
  \raisebox{0.6ex}{\poppill{Cours signé \textbullet\ Premium}{popink}{popcream}}\par
  \vspace{1pt}\popsquiggle{poppurple}\par
  \vspace{8pt}
  \begin{tcolorbox}[enhanced,colback=poporange,colframe=popink,boxrule=2.8pt,arc=13pt,
      drop shadow=popink,left=18pt,right=18pt,top=12pt,bottom=13pt,after skip=10pt]
    \poppill{#2 \textbullet\ #3}{popink}{popcream}\hspace{5pt}\poppill{#4}{white}{popink}\par\vspace{8pt}
    {\popdisplay\fontsize{14}{14}\selectfont\color{popink}LEÇON #5}\par\vspace{4pt}
    {\raggedright\hyphenpenalty=10000\exhyphenpenalty=10000\popdisplay\fontsize{25}{27}\selectfont\color{popcream}\MakeUppercase{#1}\par}
    \vspace{8pt}
    {\popxbold\fontsize{9.5}{9.5}\selectfont\color{popink}\MakeUppercase{Durée conseillée : #6}}
  \end{tcolorbox}
  \begin{tcolorbox}[enhanced,colback=white,colframe=popink,boxrule=2.2pt,arc=10pt,
      drop shadow=popink,left=16pt,right=16pt,top=10pt,bottom=10pt,after skip=8pt]
    \poppill{Dans cette leçon}{poppurple}{white}\par\vspace{5pt}
    \popsemi\fontsize{9.3}{12.6}\selectfont\color{popink}#7
  \end{tcolorbox}
  \noindent\begin{minipage}[t]{0.487\linewidth}
    \begin{tcolorbox}[enhanced,colback=popgreen,colframe=popink,boxrule=2.2pt,arc=10pt,
        drop shadow=popink,left=11pt,right=11pt,top=10pt,bottom=10pt,height=3.1cm,valign=center]
      \tcbox[colback=white,colframe=popink,boxrule=1.3pt,arc=5pt,boxsep=0pt,nobeforeafter,
        left=3pt,right=3pt,top=3pt,bottom=3pt,valign=center]{\qrcode[height=1.9cm]{https://play.google.com/store/apps/details?id=com.luvvix.onbuch&hl=fr}}%
      \hspace{8pt}\begin{minipage}[c]{0.5\linewidth}
        {\popdisplay\fontsize{11}{12.5}\selectfont\color{white}\MakeUppercase{Télécharge OnBuch}}\par\vspace{4pt}
        {\raggedright\popsemi\fontsize{8.4}{11}\selectfont\color{white}Gratuit \textbullet\ Google Play\\Tuteur IA, annales, concours, résultats\par}
      \end{minipage}
    \end{tcolorbox}
  \end{minipage}\hfill
  \begin{minipage}[t]{0.487\linewidth}
    \begin{tcolorbox}[enhanced,colback=poppurple,colframe=popink,boxrule=2.2pt,arc=10pt,
        drop shadow=popink,left=11pt,right=11pt,top=10pt,bottom=10pt,height=3.1cm,valign=center]
      \tikz[baseline=-0.7ex]\node[circle,fill=white,draw=popink,line width=1.3pt,minimum size=1.6cm,inner sep=0]{\popdisplay\fontsize{24}{24}\selectfont\color{poppurple}+};%
      \hspace{8pt}\begin{minipage}[c]{0.6\linewidth}
        {\popdisplay\fontsize{11}{12.5}\selectfont\color{white}\MakeUppercase{Passe à OnBuch+}}\par\vspace{4pt}
        {\raggedright\popsemi\fontsize{8.4}{11}\selectfont\color{white}Tous les cours signés, Tuteur IA illimité, fiches PDF — onglet OnBuch+ de l'app\par}
      \end{minipage}
    \end{tcolorbox}
  \end{minipage}
  \vspace{8pt}
  \popcontenu
  \vfill
  \signatureonbuch
  \restoregeometry
  \newpage
}

% Rangée « Ce cours contient » (page de garde)
\newcommand{\poptile}[3]{% #1 couleur #2 gros texte #3 légende
  \begin{tcolorbox}[enhanced,colback=#1,colframe=popink,boxrule=2pt,arc=9pt,shadow={2.5pt}{-2.5pt}{0pt}{popink},
      left=6pt,right=6pt,top=9pt,bottom=9pt,halign=center,valign=center,height=1.75cm,width=\linewidth,nobeforeafter]
    {\popdisplay\fontsize{12.5}{13}\selectfont\color{white}\MakeUppercase{#2}}\par\vspace{3pt}
    {\centering\popsemi\fontsize{7.8}{9.5}\selectfont\color{white}#3\par}
  \end{tcolorbox}}
\newcommand{\popcontenu}{%
  \par\noindent{\popxboldls\fontsize{7.5}{7.5}\selectfont\color{popmuted}CE COURS SIGNÉ CONTIENT}\par\vspace{7pt}
  \noindent\begin{minipage}[t]{0.235\linewidth}\poptile{popblue}{Cours}{complet, illustré, pas à pas}\end{minipage}\hfill
  \begin{minipage}[t]{0.235\linewidth}\poptile{poporange}{Méthodes}{et exercices résolus}\end{minipage}\hfill
  \begin{minipage}[t]{0.235\linewidth}\poptile{poppink}{Exercices}{type examen + corrigés}\end{minipage}\hfill
  \begin{minipage}[t]{0.235\linewidth}\poptile{popgreen}{Fiche bilan}{l'essentiel à retenir}\end{minipage}\par}

% Sommaire
\newtcolorbox{sommaire}{enhanced,colback=white,colframe=popink,boxrule=2.2pt,arc=10pt,
  drop shadow=popink,left=18pt,right=18pt,top=13pt,bottom=13pt,before skip=6pt,after skip=20pt,
  before upper={\poppill{Sommaire}{poppurple}{white}\par\vspace{9pt}\popsemi\fontsize{10.6}{14.5}\selectfont\color{popink}}}

% Signature de fin de cours — poussée tout en bas de la dernière page
\newcommand{\findecours}{%
  \par\vfill\nopagebreak
  \begin{center}\popsquiggle{poporange}\end{center}\vspace{4pt}
  \signatureonbuch
}
\AtBeginDocument{\def\llbracket{\mathopen{[\mkern-3.5mu[}}\def\rrbracket{\mathclose{]\mkern-3.5mu]}}} % intervalles d'entiers (glyphe ⟦ absent de la police)
% Glyphes absents de Fira Math : redéfinis à partir de glyphes disponibles
\AtBeginDocument{%
  \def\setminus{\mathbin{\backslash}}\let\smallsetminus\setminus
  \def\vdots{\mathord{\vbox{\baselineskip3.2pt\lineskiplimit0pt\kern4pt\hbox{.}\hbox{.}\hbox{.}}}}%
  \def\ddots{\mathinner{\mkern1mu\raise6.5pt\hbox{.}\mkern2mu\raise3.6pt\hbox{.}\mkern2mu\raise0.7pt\hbox{.}\mkern1mu}}%
  \def\triangle{\mathord{\scalebox{0.9}{$\symup{\Delta}$}}}%
  \def\nabla{\mathord{\raisebox{\depth}{\scalebox{1}[-1]{$\symup{\Delta}$}}}}%
  \def\checkmark{\popcheck{popdark}}%
  \def\star{\mathbin{\ast}}\def\bigstar{\mathord{\ast}}%
  \def\diamond{\mathbin{\scalebox{0.6}{$\Diamond$}}}%
  \def\bigcup{\mathop{\mathchoice{\raisebox{-0.3em}{\scalebox{1.7}{$\cup$}}}{\scalebox{1.25}{$\cup$}}{\cup}{\cup}}\displaylimits}%
  \def\bigcap{\mathop{\mathchoice{\raisebox{-0.3em}{\scalebox{1.7}{$\cap$}}}{\scalebox{1.25}{$\cap$}}{\cap}{\cap}}\displaylimits}%
  \def\lesseqqgtr{\mathrel{\lessgtr}}\def\bigoplus{\mathop{\oplus}\displaylimits}%
}
\providecommand{\ssi}{si et seulement si} % abréviation fréquente
\AtBeginDocument{\providecommand{\wideparen}{\overparen}} % arc de cercle

%% FIGFIX-BEGIN v5 — correctifs globaux des figures (docs/onbuch-generation/tools/figfix)
\usepackage{adjustbox}\usepackage{etoolbox}\tcbuselibrary{raster}
% toute figure TikZ/pgfplots plus large que la ligne est réduite à la largeur disponible
\BeforeBeginEnvironment{tikzpicture}{\begin{adjustbox}{max width=\linewidth}}
\AfterEndEnvironment{tikzpicture}{\end{adjustbox}}
% légendes pgfplots lisibles par-dessus les courbes
\pgfplotsset{every axis/.append style={legend style={fill=white,fill opacity=0.92,text opacity=1,draw=black!15,inner sep=3pt}}}
% étiquettes d'arêtes (syntaxe edge["..."]) avec fond blanc
\tikzset{every edge quotes/.append style={fill=white,inner sep=1.2pt}}
%% FIGFIX-END
\begin{document}

\pagegarde{Leçon 38 — Expression écrite : origine du caractère 电 et radicaux ; caractères des réseaux sociaux}{Chinois}{1ère A}{Module 5 : Mass médias et télécommunication}{ZHA38}{2 h}{Cette leçon d'expression écrite explore l'origine pictographique du caractère 电 et ses radicaux, puis présente les caractères essentiels du vocabulaire des réseaux sociaux. L'élève apprend à écrire correctement ces caractères, à rédiger un texte simple sur les médias sociaux et à traduire des phrases dans les deux sens (thème et version).
\vspace{4pt}
\begin{multicols}{2}\small
\begin{itemize}
\item À la fin de cette leçon, je sais raconter l'origine pictographique du caractère 电 (éclair → électricité)
\item À la fin de cette leçon, je sais identifier les radicaux 雨 (pluie), 田 (champ) et 讠(parole) dans les caractères du thème
\item À la fin de cette leçon, je sais écrire correctement, dans l'ordre des traits, les caractères 电, 话, 网, 信, 友, 发
\item À la fin de cette leçon, je sais distinguer les caractères proches (电/申/由, 话/活, 网/冈)
\item À la fin de cette leçon, je sais utiliser les mots-clés des réseaux sociaux : 微信, 上网, 朋友圈, 发消息, 点赞
\item À la fin de cette leçon, je sais rédiger un texte simple de 5 à 8 phrases sur mon usage des réseaux sociaux avec des connecteurs (因为, 所以, 但是, 也)
\end{itemize}\end{multicols}}

\begin{sommaire}
\begin{multicols}{2}
\begin{enumerate}
\item Origine et évolution du caractère 电
\item Radicaux et écriture des caractères du thème
\item Vocabulaire des réseaux sociaux
\item Modèle de texte commenté et connecteurs
\item Thème et version guidés
\item Entraînement à l'examen et grille d'évaluation
\item Activité d'intégration
\item Méthodes et exercices résolus
\item Exercices
\item Corrigés des exercices
\item Fiche bilan
\end{enumerate}
\end{multicols}
\end{sommaire}

\begin{objectifs}
\begin{itemize}
\item À la fin de cette leçon, je sais raconter l'origine pictographique du caractère 电 (éclair → électricité)
\item À la fin de cette leçon, je sais identifier les radicaux 雨 (pluie), 田 (champ) et 讠(parole) dans les caractères du thème
\item À la fin de cette leçon, je sais écrire correctement, dans l'ordre des traits, les caractères 电, 话, 网, 信, 友, 发
\item À la fin de cette leçon, je sais distinguer les caractères proches (电/申/由, 话/活, 网/冈)
\item À la fin de cette leçon, je sais utiliser les mots-clés des réseaux sociaux : 微信, 上网, 朋友圈, 发消息, 点赞
\item À la fin de cette leçon, je sais rédiger un texte simple de 5 à 8 phrases sur mon usage des réseaux sociaux avec des connecteurs (因为, 所以, 但是, 也)
\item À la fin de cette leçon, je sais traduire des phrases simples du français vers le chinois (thème) et du chinois vers le français (version)
\item À la fin de cette leçon, je sais m'auto-évaluer à l'aide d'une grille d'évaluation de production écrite
\end{itemize}
\end{objectifs}

\begin{prerequis}
\begin{itemize}
\item Radicaux de base vus en Seconde : 人, 口, 日, 木, 氵
\item Règles générales de l'ordre des traits (haut avant bas, gauche avant droite, horizontal avant vertical)
\item Vocabulaire des médias des leçons précédentes du Module 5 : 手机, 电脑, 电视, 电话
\item Structure de base Sujet + Verbe + Complément et particule 的
\item Connecteurs simples déjà rencontrés : 和, 也, 很
\end{itemize}
\end{prerequis}

\newpage

\coursec{Origine et évolution du caractère 电}

\courssub{Situation d'accroche : un éclair dans votre téléphone}

À Yaoundé comme à Douala, presque tous les élèves de 1re A passent leur soirée sur WeChat, TikTok ou Facebook. Mais savez-vous comment un Chinois écrit « téléphone » ? Le caractère \zh{电} (diàn, « électricité ») dessinait autrefois un éclair dans le ciel ! Aujourd'hui, il se cache dans \zh{电话} (téléphone), \zh{电视} (télévision) et \zh{电脑} (ordinateur). Dans cette leçon, nous allons remonter à l'origine de ce caractère, maîtriser ses radicaux, puis apprendre à écrire et traduire des textes simples sur les réseaux sociaux — un sujet que vous connaissez mieux que vos professeurs !

\textbf{Problématique.} Comment un dessin d'éclair vieux de plus de trois mille ans a-t-il pu devenir le signe de l'électricité, puis la clé de tous nos appareils modernes ? Comprendre l'histoire de \zh{电}, c'est comprendre la logique de l'écriture chinoise : chaque caractère raconte une histoire, et cette histoire nous aide à mémoriser, à écrire sans faute et à traduire correctement.

\courssub{Observation : un petit texte pour entrer dans la matière}

Lisons d'abord ce court texte. Ne cherchez pas à tout comprendre : repérez simplement le caractère \zh{电} et les mots qui le contiennent.

\begin{textebox}[L'origine de 电]
电原来是闪电的样子。古人看到天上的闪电，就画出了这个字。\\
\emph{Diàn yuánlái shì shǎndiàn de yàngzi. Gǔrén kàndào tiānshàng de shǎndiàn, jiù huà chūle zhège zì.}\\
« À l'origine, 电 représentait la foudre. Les anciens, en voyant l'éclair dans le ciel, ont dessiné ce caractère. »\\[4pt]
现在，电是“电”的意思：电话、电视、电脑都有这个字。\\
\emph{Xiànzài, diàn shì “diàn” de yìsi : diànhuà, diànshì, diànnǎo dōu yǒu zhège zì.}\\
« Aujourd'hui, 电 signifie “électricité” : on retrouve ce caractère dans téléphone, télévision et ordinateur. »
\end{textebox}

\textbf{Questions de compréhension (oral).}
\begin{enumerate}
\item Que représentait \zh{电} à l'origine ? (\emph{Un éclair, la foudre.})
\item Que signifie \zh{电} aujourd'hui ? (\emph{L'électricité.})
\item Citez deux mots modernes qui contiennent \zh{电}. (\emph{电话, 电视, 电脑…})
\end{enumerate}

\begin{definition}[Pictogramme]
Un \cle{pictogramme} (en chinois \zh{象形字}, xiàngxíngzì, « caractère qui imite la forme ») est un caractère dont la forme ancienne \textbf{dessine l'objet qu'il désigne}. Exemples déjà rencontrés : \zh{人} (rén, personne, silhouette debout), \zh{山} (shān, montagne, trois sommets), \zh{木} (mù, arbre, tronc et branches). Le caractère \zh{电} est lui aussi né d'un dessin : celui d'un éclair.
\end{definition}

\courssub{De l'éclair au caractère : les quatre étapes de l'évolution}

\textbf{Étape 1 — Le pictogramme.} Sur les os et carapaces de tortue (甲骨文, jiǎgǔwén, écriture divinatoire des Shang), le caractère représentait un \textbf{éclair zébrant le ciel} : un trait central \emph{sinueux}, en zigzag, comme la foudre qui descend du ciel en se tordant. Les anciens Chinois, comme les anciens Camerounais d'ailleurs, voyaient dans la foudre un phénomène puissant et mystérieux : il fallait un signe pour le nommer.

\textbf{Étape 2 — La forme sur bronze.} Sur les vases rituels en bronze (金文, jīnwén), le dessin s'épaissit : le zigzag devient plus régulier, on ajoute parfois des gouttes de pluie autour, car l'éclair vient avec l'orage.

\textbf{Étape 3 — La forme traditionnelle 電.} Avec le temps, le caractère se complexifie et devient \zh{電} (forme traditionnelle, encore utilisée à Taïwan et Hong Kong). On peut le décomposer :
\begin{itemize}
\item en haut : \zh{雨} (yǔ, la \textbf{pluie}) — la clé sémantique, car l'éclair vient du ciel et de l'orage ;
\item en bas : \zh{申} (shēn) — qui représente à l'origine la foudre qui \textbf{s'étend} dans le ciel, et qui donne aussi le son.
\end{itemize}
C'est donc un caractère \emph{sémantico-phonétique} : une partie donne le sens (雨, le ciel pluvieux), l'autre suggère le son (申).

\textbf{Étape 4 — La forme simplifiée 电.} En 1956, la République populaire de Chine simplifie l'écriture pour faciliter l'alphabétisation : \zh{電} perd sa clé de la pluie et redevient presque le pictogramme ancien : \zh{电}. C'est la forme que nous apprenons et que nous utilisons dans ce cours.

\begin{popfigure}
\begin{tikzpicture}[node distance=0.6cm and 1.1cm]
\node[draw=popblue, fill=popblueL, rounded corners, align=center, minimum width=2.6cm, minimum height=1.4cm] (picto) {\textbf{Pictogramme}\\ 甲骨文\\ \emph{dessin d'un éclair}};
\node[draw=popgreen, fill=popgreenL, rounded corners, align=center, minimum width=2.6cm, minimum height=1.4cm, right=of picto] (bronze) {\textbf{Bronze}\\ 金文\\ \emph{zigzag régulier}};
\node[draw=poporange, fill=poporangeL, rounded corners, align=center, minimum width=2.6cm, minimum height=1.4cm, right=of bronze] (trad) {\textbf{Traditionnel}\\ 電\\ 雨 + 申};
\node[draw=poppurple, fill=poppurpleL, rounded corners, align=center, minimum width=2.6cm, minimum height=1.4cm, right=of trad] (simp) {\textbf{Simplifié}\\ 电\\ \emph{5 traits}};
\draw[-{Stealth[length=3mm]}, very thick, popdark] (picto) -- (bronze);
\draw[-{Stealth[length=3mm]}, very thick, popdark] (bronze) -- (trad);
\draw[-{Stealth[length=3mm]}, very thick, popdark] (trad) -- (simp);
\end{tikzpicture}
\legende{Frise d'évolution du caractère 电 : du dessin d'éclair au caractère simplifié moderne.}
\end{popfigure}

\begin{savaistu}[Pourquoi dit-on « frapper le téléphone » ?]
En chinois, « passer un appel » se dit \zh{打电话} (dǎ diànhuà), littéralement « \textbf{frapper} le téléphone ». Pourquoi \zh{打} (dǎ, frapper, actionner) ? Parce que les premiers téléphones, au début du XXe siècle, se manipulaient avec une \textbf{manivelle} qu'il fallait tourner énergiquement pour appeler l'opératrice ! Le verbe est resté dans la langue même si la manivelle a disparu. De même, on dit \zh{发邮件} (fā yóujiàn, « envoyer un courriel ») avec \zh{发} (fā, émettre, expédier).
\end{savaistu}

\courssub{Glissement de sens : de l'éclair à l'électricité}

Comment passe-t-on de « éclair » à « électricité » ? Le raisonnement des savants chinois du XIXe siècle est très naturel : quand l'électricité arrive en Chine, on cherche un mot pour la nommer. Or, le phénomène électrique le plus visible dans la nature, c'est \textbf{l'éclair}. Le vieux caractère \zh{电} est donc « recyclé » : il prend le sens moderne d'\emph{électricité}, puis, par extension, désigne \textbf{tout ce qui fonctionne à l'électricité}.

Ce phénomène existe aussi en français : le mot « téléphone » vient du grec \emph{têle} (loin) et \emph{phônê} (voix). Le chinois procède de la même façon, mais avec ses propres briques : au lieu de racines grecques, il combine des caractères.

\begin{propriete}[Règle de formation : 电 + nom d'organe ou de fonction]
Pour nommer un appareil moderne, le chinois combine \cle{电} (électricité) avec un caractère indiquant un \textbf{organe} ou une \textbf{fonction} :
\begin{center}
电 (électricité) + organe/fonction = nom d'un appareil électrique
\end{center}
Ainsi : 电 + 话 (parole) = 电话 « parole électrique » = téléphone ; 电 + 脑 (cerveau) = 电脑 « cerveau électrique » = ordinateur. Cette logique est \textbf{régulière et prévisible} : connaître le sens de chaque brique permet de deviner le sens du mot composé.
\end{propriete}

\begin{exemplebox}[La famille des mots en 电]
\begin{itemize}
\item \zh{电话} \emph{diànhuà} : 电 (électricité) + 话 (parole) = « parole électrique » → \textbf{téléphone}.\\
\zh{我每天晚上给妈妈打电话。} \emph{Wǒ měitiān wǎnshang gěi māma dǎ diànhuà.} « J'appelle ma mère au téléphone tous les soirs. »
\item \zh{电视} \emph{diànshì} : 电 + 视 (regarder, vision) = « vision électrique » → \textbf{télévision}.\\
\zh{我们在家看电视。} \emph{Wǒmen zài jiā kàn diànshì.} « Nous regardons la télévision à la maison. »
\item \zh{电脑} \emph{diànnǎo} : 电 + 脑 (cerveau) = « cerveau électrique » → \textbf{ordinateur}.\\
\zh{我用电脑做作业。} \emph{Wǒ yòng diànnǎo zuò zuòyè.} « J'utilise l'ordinateur pour faire mes devoirs. »
\item \zh{电影} \emph{diànyǐng} : 电 + 影 (ombre, image) = « ombre électrique » → \textbf{film, cinéma}.\\
\zh{周末我们去看电影。} \emph{Zhōumò wǒmen qù kàn diànyǐng.} « Le week-end, nous allons voir un film. »
\item \zh{电子邮件} \emph{diànzǐ yóujiàn} : 电子 (électronique) + 邮件 (courrier) = \textbf{courriel, e-mail}.\\
\zh{老师给我发了电子邮件。} \emph{Lǎoshī gěi wǒ fāle diànzǐ yóujiàn.} « Le professeur m'a envoyé un courriel. »
\end{itemize}
\end{exemplebox}

\begin{popfigure}
\begin{center}\tcbox[colback=popblue,colframe=popblue,coltext=white,arc=9pt,boxsep=4pt,left=6pt,right=6pt]{\bfseries\small 电 diàn \emph{électricité}}\end{center}
\vspace{-2pt}
\begin{tcbraster}[raster columns=3,raster equal height=rows,raster column skip=6pt,raster row skip=6pt,enhanced,blankest]
\begin{tcolorbox}[enhanced,colback=popblue,colframe=popblue,coltext=white,boxrule=0.9pt,arc=3pt,left=4pt,right=4pt,top=3pt,bottom=3pt,boxsep=1pt,halign=left]\bfseries\scriptsize 电话 diànhuà \emph{téléphone} 话 = parole\end{tcolorbox}
\begin{tcolorbox}[enhanced,colback=poporange,colframe=poporange,coltext=white,boxrule=0.9pt,arc=3pt,left=4pt,right=4pt,top=3pt,bottom=3pt,boxsep=1pt,halign=left]\bfseries\scriptsize 电视 diànshì \emph{télévision} 视 = vision\end{tcolorbox}
\begin{tcolorbox}[enhanced,colback=popgreen,colframe=popgreen,coltext=white,boxrule=0.9pt,arc=3pt,left=4pt,right=4pt,top=3pt,bottom=3pt,boxsep=1pt,halign=left]\bfseries\scriptsize 电脑 diànnǎo \emph{ordinateur} 脑 = cerveau\end{tcolorbox}
\begin{tcolorbox}[enhanced,colback=poppurple,colframe=poppurple,coltext=white,boxrule=0.9pt,arc=3pt,left=4pt,right=4pt,top=3pt,bottom=3pt,boxsep=1pt,halign=left]\bfseries\scriptsize 电影 diànyǐng \emph{film} 影 = ombre\end{tcolorbox}
\begin{tcolorbox}[enhanced,colback=popgold,colframe=popgold,coltext=white,boxrule=0.9pt,arc=3pt,left=4pt,right=4pt,top=3pt,bottom=3pt,boxsep=1pt,halign=left]\bfseries\scriptsize 电子邮件 diànzǐ yóujiàn \emph{courriel}\end{tcolorbox}
\end{tcbraster}

\legende{Carte mentale lexicale : la famille des mots construits sur 电.}
\end{popfigure}

\begin{center}
\begin{tabularx}{\linewidth}{|X|X|X|X|}
\hline
\rowcolor{popblueL} Caractères & Pinyin & Nature & Traduction \\
\hline
电 & diàn & nom & électricité ; éclair (sens ancien) \\
\hline
闪电 & shǎndiàn & nom & éclair, foudre \\
\hline
电话 & diànhuà & nom & téléphone \\
\hline
打电话 & dǎ diànhuà & verbe + compl. & passer un appel téléphonique \\
\hline
电视 & diànshì & nom & télévision \\
\hline
电脑 & diànnǎo & nom & ordinateur \\
\hline
电影 & diànyǐng & nom & film, cinéma \\
\hline
电子邮件 & diànzǐ yóujiàn & nom & courriel, e-mail \\
\hline
原来 & yuánlái & adverbe & à l'origine, en fait \\
\hline
古人 & gǔrén & nom & les anciens, les gens de l'Antiquité \\
\hline
样子 & yàngzi & nom & apparence, forme \\
\hline
\end{tabularx}
\end{center}

\courssub{Écriture du caractère 电 : l'ordre des 5 traits}

Le caractère \zh{电} s'écrit en \textbf{5 traits}, dans un ordre précis qu'il faut respecter (règle générale : de haut en bas, de gauche à droite, l'extérieur avant l'intérieur, et le trait qui ferme en dernier).

\begin{enumerate}
\item \textbf{Trait 1 : 丨} — le trait vertical central, qui descend du haut vers le milieu.
\item \textbf{Trait 2 : (trait brisé)} — l'angle horizontal-vertical : on trace le haut du cadre puis on descend à droite, d'un seul mouvement.
\item \textbf{Trait 3 : 一} — l'horizontal médian à l'intérieur du cadre.
\item \textbf{Trait 4 : 一} — l'horizontal inférieur qui ferme le bas du cadre.
\item \textbf{Trait 5 : 乚} — le \textbf{crochet vertical final} : il part du centre, descend \emph{en dessous} du cadre et se termine par un petit crochet vers le haut. C'est ce trait qui « ferme » le caractère et qui rappelle le zigzag de l'éclair ancien.
\end{enumerate}

\begin{attention}[Erreur fréquente des francophones : le trait central qui ne dépasse pas]
Beaucoup d'élèves francophones écrivent \zh{电} avec le trait central qui \textbf{s'arrête au niveau du cadre}, sans dépasser en bas. Résultat : le caractère ressemble à \zh{田} (tián, champ) ou à \zh{由} (yóu) mal formé, et devient illisible pour un Chinois. Retenez : le 5\ieme{} trait \zh{乚} doit \textbf{descendre sous le cadre} et \textbf{crocheter} vers le haut, comme la pointe d'un éclair. Sans ce dépassement, il n'y a pas de caractère 电 !
\end{attention}

\courssub{Caractères proches à ne pas confondre}

Quatre caractères partagent le même « cadre » avec un trait central. Seule la façon dont le trait central sort du cadre les distingue. C'est un piège classique à l'examen.

\begin{center}
\begin{tabularx}{\linewidth}{|X|X|X|X|}
\hline
\rowcolor{popblueL} Caractère & Pinyin & Traduction & Différence de forme \\
\hline
电 & diàn & électricité & le trait central dépasse \textbf{en bas} avec un crochet 乚 \\
\hline
申 & shēn & déclarer, s'étendre & le trait vertical traverse et dépasse \textbf{en haut et en bas}, sans crochet \\
\hline
由 & yóu & depuis, à cause de & le trait vertical dépasse \textbf{en haut} seulement \\
\hline
甲 & jiǎ & premier, armure, ongle & le trait vertical dépasse \textbf{en bas} seulement, droit, sans crochet \\
\hline
\end{tabularx}
\end{center}

\textbf{Moyen mnémotechnique.} Imaginez le cadre comme le ciel nuageux : l'éclair \zh{电} descend et \emph{frappe le sol en crochettant} ; \zh{申} est la foudre qui \emph{s'étend} des deux côtés ; \zh{由} est une pousse qui \emph{sort vers le haut} ; \zh{甲} est un clou \emph{planté vers le bas}.

\begin{exemplebox}[Comparer dans des phrases]
\begin{itemize}
\item \zh{电} : \zh{我家有电视。} \emph{Wǒ jiā yǒu diànshì.} « Chez moi, il y a une télévision. »
\item \zh{申} : \zh{他申请了中国的大学。} \emph{Tā shēnqǐngle Zhōngguó de dàxué.} « Il a fait une demande d'inscription dans une université chinoise. »
\item \zh{由} : \zh{从雅温得到杜阿拉很远。} \emph{Cóng Yǎwēndé dào Dù'ālā hěn yuǎn.} « De Yaoundé à Douala, c'est loin. »
\end{itemize}
\end{exemplebox}

\begin{exoresolu}[Identifier et justifier]
\textbf{Énoncé.} Voici trois mots : \zh{电脑}, \zh{电影}, \zh{电话}. (1) Traduisez chaque mot en français. (2) Expliquez la logique de composition de chacun (que signifie la deuxième brique ?). (3) Pourquoi le caractère \zh{电} se trouve-t-il dans ces trois mots ?
\tcblower
\textbf{Solution détaillée.}
\begin{enumerate}
\item \zh{电脑} (diànnǎo) = ordinateur ; \zh{电影} (diànyǐng) = film ; \zh{电话} (diànhuà) = téléphone.
\item \zh{脑} (nǎo) signifie « cerveau » : l'ordinateur est un « cerveau électrique ». \zh{影} (yǐng) signifie « ombre, image » : le film est une « ombre électrique » (projection d'images). \zh{话} (huà) signifie « parole » : le téléphone transmet la « parole électrique ».
\item Le caractère \zh{电} signifie aujourd'hui « électricité ». D'après la règle de formation, 电 + nom d'organe ou de fonction forme le nom d'un appareil électrique moderne. Ces trois appareils fonctionnent à l'électricité, d'où la présence de \zh{电} en tête de mot.
\end{enumerate}
\textbf{Méthode à retenir.} Face à un mot composé inconnu, découpez-le en briques, traduisez chaque brique, puis devinez le sens global : c'est la stratégie de lecture la plus efficace en chinois.
\end{exoresolu}

\begin{exercice}[À vous de jouer : écriture et discrimination]
\begin{enumerate}
\item Écrivez le caractère \zh{电} cinq fois en respectant l'ordre des 5 traits ; numérotez chaque trait sur votre première copie.
\item Complétez avec le bon caractère (\zh{电}, \zh{申}, \zh{由}, \zh{甲}) : \zh{＿视} (télévision) ; \zh{＿请} (faire une demande) ; \zh{因＿} (en raison de).
\item Traduisez en français : \zh{我用电脑发电子邮件。} \emph{Wǒ yòng diànnǎo fā diànzǐ yóujiàn.}
\item Devinez le sens probable du mot \zh{电梯} (diàntī), sachant que \zh{梯} (tī) signifie « échelle, escalier ».
\end{enumerate}
\end{exercice}

\begin{corrige}[Exercice « À vous de jouer »]
\begin{enumerate}
\item Ordre des traits : 丨 (vertical central) → (trait brisé) (angle haut-droite) → 一 (horizontal médian) → 一 (horizontal du bas) → 乚 (crochet final qui descend sous le cadre). Vérifiez que le dernier trait dépasse bien en bas.
\item \zh{电视} (télévision) ; \zh{申请} (faire une demande, postuler) ; \zh{由于} (en raison de, à cause de).
\item « J'utilise l'ordinateur pour envoyer un courriel. » (用 + outil = « avec, au moyen de ».)
\item \zh{电梯} = « échelle électrique » → l'\textbf{ascenseur}. La logique 电 + fonction s'applique encore !
\end{enumerate}
\end{corrige}

\begin{aretenir}[L'essentiel de la section : 电]
\begin{itemize}
\item \zh{电} est né d'un \cle{pictogramme} : le dessin d'un éclair zébrant le ciel.
\item Forme traditionnelle \zh{電} = 雨 (pluie) + 申 (foudre qui s'étend) ; forme simplifiée \zh{电} (1956).
\item Glissement de sens : éclair → électricité → appareil électrique.
\item Règle d'or : \textbf{电 + organe/fonction = appareil moderne} (电话, 电视, 电脑, 电影, 电子邮件).
\item Écriture : 5 traits, le crochet final 乚 descend \emph{sous} le cadre.
\item Ne pas confondre : 电 (dépasse en bas avec crochet), 申 (dépasse des deux côtés), 由 (dépasse en haut), 甲 (dépasse en bas, droit).
\end{itemize}
\end{aretenir}

% ============================================================
% SECTION 3 : VOCABULAIRE DES RÉSEAUX SOCIAUX
% ============================================================
\courssub{Vocabulaire des réseaux sociaux : WeChat, Douyin et actions courantes}

Au Cameroun comme en Chine, les réseaux sociaux font partie du quotidien. En Chine, l'application reine est \zh{微信} (Wēixìn, WeChat) : c'est à la fois WhatsApp, Facebook, PayPal et bien plus. La version chinoise de TikTok s'appelle \zh{抖音} (Dǒuyīn). Maîtriser le vocabulaire ci-dessous vous permettra de comprendre des tutoriels, de discuter avec des correspondants chinois et de rédiger vos propres publications.

\begin{textebox}[Dialogue : « Qu'est-ce que tu fais sur ton téléphone ? »]
A: 你在看什么？\\
\emph{Nǐ zài kàn shénme?}\\
« Tu regardes quoi ? »\\[4pt]
B: 我在刷抖音，看短视频。\\
\emph{Wǒ zài shuā Dǒuyīn, kàn duǎnshìpín.}\\
« Je fais défiler Douyin, je regarde des vidéos courtes. »\\[4pt]
A: 我也喜欢刷视频。有时候我在微信朋友圈发照片。\\
\emph{Wǒ yě xǐhuan shuā shìpín. Yǒushíhou wǒ zài Wēixìn péngyouquān fā zhàopiàn.}\\
« Moi aussi j'aime regarder des vidéos. Parfois je poste des photos sur les Moments WeChat. »\\[4pt]
B: 我常发语音消息给朋友，不爱打字。\\
\emph{Wǒ cháng fā yǔyīn xiāoxi gěi péngyou, bù ài dǎzì.}\\
« J'envoie souvent des messages vocaux à mes amis, je n'aime pas taper du texte. »
\end{textebox}

\begin{center}
\begin{tabularx}{\linewidth}{|X|X|X|X|}
\hline
\rowcolor{popblueL} Caractères & Pinyin & Nature & Traduction \\
\hline
微信 & Wēixìn & nom propre & WeChat (application tout-en-un) \\
\hline
抖音 & Dǒuyīn & nom propre & Douyin (TikTok version Chine) \\
\hline
朋友圈 & péngyouquān & nom & « Cercle d'amis » (fil d'actualité WeChat) \\
\hline
发 / 发布 & fā / fābù & verbe & poster, publier, envoyer \\
\hline
消息 / 信息 & xiāoxi / xìnxī & nom & message, information \\
\hline
语音消息 & yǔyīn xiāoxi & nom & message vocal \\
\hline
视频通话 & shìpín tōnghuà & nom & appel vidéo \\
\hline
短视频 & duǎnshìpín & nom & vidéo courte (format TikTok/Reels) \\
\hline
直播 & zhíbō & nom/verbe & direct, live / faire un live \\
\hline
点赞 & diǎnzàn & verbe & liker, mettre un « j'aime » \\
\hline
评论 & pínglùn & verbe/nom & commenter / commentaire \\
\hline
转发 / 分享 & zhuǎnfā / fēnxiǎng & verbe & partager, transférer (repost) \\
\hline
关注 & guānzhù & verbe & suivre (un compte), s'abonner \\
\hline
粉丝 & fěnsī & nom & fans, abonnés (translit. de « fans ») \\
\hline
网红 & wǎnghóng & nom & influenceur, star du web (litt. « rouge sur le net ») \\
\hline
表情包 & biǎoqíngbāo & nom & pack d'emojis, autocollants, mèmes \\
\hline
红包 & hóngbāo & nom & enveloppe rouge (cadeau d'argent numérique) \\
\hline
上网 & shàngwǎng & verbe & aller sur Internet, se connecter \\
\hline
登录 & dēnglù & verbe & se connecter (login) \\
\hline
退出 & tuìchū & verbe & se déconnecter (logout) \\
\hline
账号 & zhànghào & nom & compte, identifiant \\
\hline
密码 & mìmǎ & nom & mot de passe \\
\hline
刷 & shuā & verbe & faire défiler (un fil), swiper \\
\hline
\end{tabularx}
\end{center}

\begin{propriete}[Structures utiles pour les actions sur les réseaux]
\begin{itemize}
\item \textbf{Publier / Envoyer} : \zh{发} (fā) + \textbf{objet} (照片, 消息, 视频, 红包, 朋友圈). Ex: \zh{发朋友圈} (poster sur Moments), \zh{发语音} (envoyer un vocal).
\item \textbf{Regarder / Consommer} : \zh{看} (kàn) / \zh{刷} (shuā) + \textbf{objet} (视频, 短视频, 直播, 朋友圈). \zh{刷} implique un défilement rapide, passif.
\item \textbf{Interagir} : \zh{点赞} (liker), \zh{评论} (commenter), \zh{转发} / \zh{分享} (partager), \zh{关注} (suivre).
\item \textbf{Communication} : \zh{打视频电话} / \zh{视频聊天} (appel vidéo), \zh{发语音} (envoyer un vocal), \zh{打字} (taper du texte).
\end{itemize}
\end{propriete}

\begin{attention}[Pièges fréquentes : 信息 vs 消息 ; 转发 vs 分享]
\begin{itemize}
\item \zh{信息} (xìnxī) = « information, données » (plus formel, technique). \zh{消息} (xiāoxi) = « message, nouvelle, info » (quotidien, personnel). Sur WeChat, on dit \zh{发消息} (envoyer un message) ou \zh{收到消息} (recevoir un message).
\item \zh{转发} (zhuǎnfā) = « transférer, renvoyer » (action technique : bouton « forward », garde l'auteur original). \zh{分享} (fēnxiǎng) = « partager » (sens large, peut être vers d'autres apps ou en story). À l'oral, on utilise souvent \zh{转发} pour le bouton spécifique de WeChat.
\item \zh{上网} (shàngwǎng) = « aller sur Internet » (verbe d'action). Ne pas confondre avec \zh{登录} (dēnglù, « se logger sur un compte »).
\end{itemize}
\end{attention}

\begin{exemplebox}[Phrases types pour s'exprimer sur les réseaux]
\begin{itemize}
\item \zh{我每天都上微信。} \emph{Wǒ měitiān dōu shàng Wēixìn.} « Je vais sur WeChat tous les jours. »
\item \zh{请关注我的账号。} \emph{Qǐng guānzhù wǒ de zhànghào.} « Merci de suivre mon compte. »
\item \zh{这个视频很火，我转发给你了。} \emph{Zhège shìpín hěn huǒ, wǒ zhuǎnfā gěi nǐ le.} « Cette vidéo est virale, je te l'ai transférée. » (\zh{火} huǒ = populaire, « en feu »).
\item \zh{我忘记密码了，登录不上。} \emph{Wǒ wàngjì mìmǎ le, dēnglù bù shàng.} « J'ai oublié mon mot de passe, je n'arrive pas à me connecter. » (Complément de résultat \zh{不上} = ne pas réussir à).
\item \zh{主播在直播带货。} \emph{Zhǔbō zài zhíbō dàihuò.} « Le streamer fait de la vente en direct. » (\zh{带货} dàihuò = vendre des produits en live).
\end{itemize}
\end{exemplebox}

% ============================================================
% SECTION 4 : MODÈLE DE TEXTE COMMENTÉ ET CONNECTEURS
% ============================================================
\courssub{Modèle de texte commenté : « Mon quotidien numérique »}

Voici un texte modèle (niveau HSK 3 / Première A) respectant les attendus de l'expression écrite : paragraphe unique, connecteurs logiques, vocabulaire du module, structures variées (SVO, \zh{把}, \zh{被}, compléments de durée/degré).

\begin{textebox}[Mon quotidien numérique]
我是喀麦隆的高中生，叫马里。\\
\emph{Wǒ shì Kāmàilóng de gāozhōngshēng, jiào Mǎlǐ.}\\
« Je suis un lycéen camerounais, je m'appelle Mari. »\\[4pt]
\textbf{首先}，我每天都用手机上微信。\\
\emph{Shǒuxiān, wǒ měitiān dōu yòng shǒujī shàng Wēixìn.}\\
« \textbf{D'abord}, j'utilise mon téléphone pour aller sur WeChat tous les jours. »\\[4pt]
\textbf{然后}，我习惯看朋友圈，给好友的照片点赞、评论。\\
\emph{Ránhòu, wǒ xíguàn kàn péngyouquān, gěi hǎoyǒu de zhàopiàn diǎnzàn, pínglùn.}\\
« \textbf{Ensuite}, j'ai l'habitude de regarder les Moments, de liker et commenter les photos de mes amis. »\\[4pt]
\textbf{有时候}，我也刷抖音看短视频，觉得很有意思。\\
\emph{Yǒushíhou, wǒ yě shuā Dǒuyīn kàn duǎnshìpín, juéde hěn yǒu yìsi.}\\
« \textbf{Parfois}, je fais aussi défiler Douyin pour regarder des vidéos courtes, je trouve ça très intéressant. »\\[4pt]
\textbf{因为}学习很忙，\textbf{所以}我不看直播，也不玩游戏。\\
\emph{Yīnwèi xuéxí hěn máng, suǒyǐ wǒ bù kàn zhíbō, yě bù wán yóuxì.}\\
« \textbf{Parce que} les études sont chargées, \textbf{donc} je ne regarde pas de lives et je ne joue pas aux jeux. »\\[4pt]
\textbf{如果}有空，我会给妈妈发视频通话。\\
\emph{Rúguǒ yǒukòng, wǒ huì gěi māma fā shìpín tōnghuà.}\\
« \textbf{Si} j'ai du temps, je ferai un appel vidéo à maman. »\\[4pt]
\textbf{总之}，手机让我的生活更方便，但我要控制时间。\\
\emph{Zǒngzhī, shǒujī ràng wǒ de shēnghuó gèng fāngbiàn, dàn wǒ yào kòngzhì shíjiān.}\\
« \textbf{En résumé}, le téléphone rend ma vie plus pratique, mais je dois contrôler mon temps. »
\end{textebox}

\begin{propriete}[Les connecteurs logiques indispensables (niveau Première)]
\begin{center}
\begin{tabularx}{\linewidth}{|X|X|X|X|}
\hline
\rowcolor{popblueL} Chinois & Pinyin & Français & Emploi / Position \\
\hline
首先 & shǒuxiān & d'abord, en premier & Début de phrase, énumération \\
然后 & ránhòu & ensuite, puis & Suite d'actions chronologiques \\
还有 & hái yǒu & de plus, en outre & Ajout d'un élément \\
另外 & lìngwài & par ailleurs, en plus & Changement d'aspect, ajout \\
因为…所以… & yīnwèi… suǒyǐ… & parce que… donc… & Cause \rightarrow Conséquence (corrélatif) \\
虽然…但是… & suīrán… dànshì… & bien que… mais… & Concession \rightarrow Opposition (corrélatif) \\
如果…就… & rúguǒ… jiù… & si… alors… & Condition \rightarrow Conséquence (corrélatif) \\
总之 & zǒngzhī & en résumé, bref & Conclusion \\
其实 & qíshí & en fait, en réalité & Rectification, précision \\
比如说 & bǐrúshuō & par exemple & Illustration \\
\hline
\end{tabularx}
\end{center}
\end{propriete}

\begin{methode}[Comment rédiger un paragraphe cohérent (Expression écrite)]
\begin{enumerate}
\item \textbf{Analyser le sujet :} Identifier le thème (ex: « Mon application préférée »), le temps (présent/habitudes), la personne (je/nous).
\item \textbf{Planifier en 3--4 idées :} 1. Présentation (Qui ? Quoi ?). 2. Fréquence / Habitudes (Connecteurs: 首先, 然后). 3. Raisons / Préférences (Connecteurs: 因为…所以…, 觉得). 4. Conclusion / Avenir (Connecteurs: 总之, 如果…).
\item \textbf{Choisir le vocabulaire :} Réutiliser les mots du module (微信, 发消息, 视频, 方便).
\item \textbf{Écrire en chinois directement :} Ne pas traduire mot à mot du français. Utiliser les structures chinoises (Sujet + Temps + Lieu + Verbe + Complément).
\item \textbf{Relire et corriger :} Vérifier l'ordre des mots, les mots-outils (\zh{了, 着, 过, 的, 得, 地}), les classificateurs, les tons notés en pinyin.
\end{enumerate}
\end{methode}

\begin{exoresolu}[Analyse du texte modèle]
\textbf{Énoncé.} Relevez dans le texte « Mon quotidien numérique » : (1) tous les connecteurs ; (2) deux structures \zh{把} ou \zh{被} ou complément de degré (s'il y en a) ; (3) le vocabulaire spécifique aux réseaux sociaux.
\tcblower
\textbf{Solution.}
\begin{enumerate}
\item Connecteurs : \zh{首先} (shǒuxiān), \zh{然后} (ránhòu), \zh{有时候} (yǒushíhou), \zh{因为…所以…} (yīnwèi… suǒyǐ…), \zh{如果…} (rúguǒ…), \zh{总之} (zǒngzhī).
\item Structures : Pas de \zh{把}/\zh{被} ici (texte descriptif). Complément de degré : \zh{很有意思} (hěn yǒu yìsi, \zh{很} + adj). Structure \zh{给...点赞} (construction causative/prépositionnelle). \zh{让...更方便} (causatif \zh{让}).
\item Vocabulaire réseaux : \zh{微信}, \zh{朋友圈}, \zh{点赞}, \zh{评论}, \zh{刷}, \zh{抖音}, \zh{短视频}, \zh{直播}, \zh{视频通话}.
\end{enumerate}
\end{exoresolu}

% ============================================================
% SECTION 5 : THÈME ET VERSION GUIDÉS
% ============================================================
\courssub{Thème et version guidés : méthode et entraînement}

La traduction (thème = français → chinois ; version = chinois → français) est une compétence évaluée à l'examen. Elle teste la maîtrise des structures, du vocabulaire et de la syntaxe.

\begin{methode}[Méthode du thème (Français → Chinois)]
\begin{enumerate}
\item \textbf{Segmenter} la phrase française en unités de sens (Sujet, Verbe, Compléments).
\item \textbf{Identifier la structure chinoise cible} (SVO, \zh{把}, \zh{被}, \zh{在}, \zh{有}, comparatif \zh{比}, etc.).
\item \textbf{Traduire les mots-clés} (vocabulaire précis, classificateurs).
\item \textbf{Assembler} dans l'ordre chinois : Sujet + (Temps) + (Lieu/Manière) + Verbe + (Compléments/Objet) + (Particules aspectuelles).
\item \textbf{Vérifier} : mots-outils (\zh{了} action terminée, \zh{过} expérience, \zh{着} état), \zh{的/得/地}, ordre des compléments (Temps avant Lieu).
\end{enumerate}
\end{methode}

\begin{methode}[Méthode de la version (Chinois → Français)]
\begin{enumerate}
\item \textbf{Segmenter} la phrase chinoise par les mots-outils et la ponctuation.
\item \textbf{Identifier la structure} (repérer le verbe principal, le sujet, les compléments).
\item \textbf{Traduire littéralement} mot à mot en gardant l'ordre chinois.
\item \textbf{Reformuler en français correct} (conjugaison, prépositions, articles, ordre SVO français).
\item \textbf{Vérifier} le sens global, les temps (passé/composé pour \zh{了}/\zh{过}), la voix (passif pour \zh{被}).
\end{enumerate}
\end{methode}

\begin{exoresolu}[Thème guidé : phrases isolées]
\textbf{Énoncé.} Traduisez en chinois (caractères + pinyin).
\begin{enumerate}
\item J'envoie souvent des messages vocaux à mes amis sur WeChat.
\item Hier soir, j'ai regardé un film sur l'ordinateur.
\item Si tu as le temps, ajoute-moi sur WeChat (scanne mon QR code).
\item Cette vidéo est très drôle, je l'ai partagée sur mon Moments.
\item Je n'arrive pas à me connecter car j'ai oublié mon mot de passe.
\end{enumerate}
\tcblower
\textbf{Solution détaillée.}
\begin{enumerate}
\item \zh{我经常在微信上给朋友发语音消息。} \emph{Wǒ jīngcháng zài Wēixìn shàng gěi péngyou fā yǔyīn xiāoxi.} (Structure : Sujet + Fréquence + Lieu + Verbe + Destinataire + Objet).
\item \zh{昨天晚上我在电脑上看了一个电影。} \emph{Zuótiān wǎnshang wǒ zài diànnǎo shàng kànle yí ge diànyǐng.} (Temps passé \zh{了}, classificateur \zh{个}, lieu \zh{在...上}).
\item \zh{如果你有空，加我的微信（扫我的二维码）。} \emph{Rúguǒ nǐ yǒukòng, jiā wǒ de Wēixìn (sǎo wǒ de èrwéimǎ).} (Condition \zh{如果}, verbe \zh{加} pour ajouter contact, \zh{二维码} èrwéimǎ = QR code).
\item \zh{这个视频很搞笑，我转发到朋友圈了。} \emph{Zhège shìpín hěn gǎoxiào, wǒ zhuǎnfā dào péngyouquān le.} (\zh{搞笑} gǎoxiào = drôle, \zh{转发到} = partager vers, \zh{了} action réalisée).
\item \zh{我登录不上，因为我忘记密码了。} \emph{Wǒ dēnglù bù shàng, yīnwèi wǒ wàngjì mìmǎ le.} (Complément de résultat \zh{不上} = ne pas réussir à, cause \zh{因为}).
\end{enumerate}
\end{exoresolu}

\begin{exoresolu}[Version guidée : court paragraphe]
\textbf{Énoncé.} Traduisez en français :
\zh{马里是我的中国朋友。他住在北京。上个周末，他发微信给我，说想看喀麦隆的足球比赛。我给他发了一个视频，是喀麦隆队对埃及队的比赛。他看了以后，觉得非常精彩，给我发了一个大红包。}
\tcblower
\textbf{Solution.}
« Mari est mon ami chinois. Il habite à Pékin. Le week-end dernier, il m'a envoyé un message sur WeChat pour dire qu'il voulait regarder un match de football du Cameroun. Je lui ai envoyé une vidéo, c'était le match de l'équipe du Cameroun contre l'équipe d'Égypte. Après l'avoir regardée, il a trouvé ça très palpitant, et m'a envoyé un gros enveloppe rouge (cadeau d'argent). »

\textbf{Notes de traduction :}
\begin{itemize}
\item \zh{上个周末} (shàng ge zhōumò) = « le week-end dernier » (passé récent, pas de \zh{了} nécessaire sur le temps).
\item \zh{说想看} (shuō xiǎng kàn) = « dire vouloir regarder » (verbe \zh{说} introduit le contenu).
\item \zh{对} (duì) = « contre » (pour les matchs).
\item \zh{看了以后} (kànle yǐhòu) = « après avoir regardé » (aspect accompli \zh{了} + \zh{以后}).
\item \zh{精彩} (jīngcǎi) = « splendide, palpitant, excitant » (pour un spectacle/match).
\item \zh{大红包} (dà hóngbāo) = « gros enveloppe rouge » (montant important).
\end{itemize}
\end{exoresolu}

\begin{exercice}[Thème et version : entraînement individuel]
\textbf{Thème (Français → Chinois). Traduisez en caractères + pinyin.}
\begin{enumerate}
\item Je suis en train de regarder un direct (live) sur Douyin.
\item Merci de liker et de partager ma vidéo.
\item Je n'ai pas de fans, je ne suis pas une star du web.
\item Nous discutons par appel vidéo tous les soirs.
\item Comme il pleut, je reste à la maison pour surfer sur Internet.
\end{enumerate}

\textbf{Version (Chinois → Français). Traduisez en français.}
\begin{enumerate}
\item \zh{请扫我的二维码，关注我的账号。}
\item \zh{我不喜欢打字，我只发语音。}
\item \zh{因为网不好，视频通话断了。}
\item \zh{这个网红有很多粉丝，她的直播很有意思。}
\item \zh{退出账号前，记得保存密码。}
\end{enumerate}
\end{exercice}

\begin{corrige}[Exercice « Thème et version »]
\textbf{Thème :}
\begin{enumerate}
\item \zh{我正在刷抖音看直播。} \emph{Wǒ zhèngzài shuā Dǒuyīn kàn zhíbō.} (\zh{正在} action en cours).
\item \zh{请给我的视频点赞、转发。} \emph{Qǐng gěi wǒ de shìpín diǎnzàn, zhuǎnfā.} (\zh{请} impératif poli).
\item \zh{我没有粉丝，不是网红。} \emph{Wǒ méiyǒu fěnsī, bù shì wǎnghóng.}
\item \zh{我们每天晚上都打视频电话。} \emph{Wǒmen měitiān wǎnshang dōu dǎ shìpín diànhuà.}
\item \zh{因为下雨，我在家上网。} \emph{Yīnwèi xiàyǔ, wǒ zài jiā shàngwǎng.}
\end{enumerate}

\textbf{Version :}
\begin{enumerate}
\item « Merci de scanner mon QR code et de suivre mon compte. »
\item « Je n'aime pas taper du texte, j'envoie seulement des vocaux. »
\item « Parce que le réseau est mauvais, l'appel vidéo a coupé. » (\zh{断了} duànle = coupé/rompu).
\item « Cet influenceur a beaucoup de fans, ses lives sont très intéressants. »
\item « Avant de te déconnecter, n'oublie pas de sauvegarder le mot de passe. » (\zh{记得} jìde = se souvenir de / ne pas oublier de).
\end{enumerate}
\end{corrige}

% ============================================================
% SECTION 6 : ENTRAÎNEMENT À L'EXAMEN ET GRILLE D'ÉVALUATION
% ============================================================
\courssub{Entraînement à l'examen et grille d'évaluation}

Cette section simule les épreuves types du baccalauréat (Première A, LV2 Chinois) : Compréhension de l'écrit, Expression écrite (rédaction courte), Traduction (Thème/Version).

\begin{exercice}[Épreuve type Bac — Durée : 1h30]
\textbf{Partie A : Compréhension de l'écrit (6 points)}
Lisez le texte suivant et répondez aux questions en français.

\zh{王明是喀麦隆留学生，在广州读大学。他很喜欢用微信。每天早上，他先看朋友圈，看中国朋友发的照片和视频。然后他给妈妈发语音，报平安。晚上，他常和同学视频通话，一起复习功课。上个星期天，他看了一场直播，是关于中国茶文化的。他觉得很有意思，就转发给了喀麦隆的同学。王明说：「微信让我和家人、朋友更近了。」}

\textbf{Questions :}
\begin{enumerate}
\item Où habite Wang Ming ? Quelle est sa nationalité ?
\item Citez trois actions qu'il fait sur WeChat le matin et le soir.
\item Qu'a-t-il regardé dimanche dernier ? Qu'a-t-il fait ensuite ?
\item Traduisez la dernière phrase : \zh{微信让我和家人、朋友更近了。}
\item Selon le texte, pourquoi WeChat est-il important pour lui ? (Inférence)
\end{enumerate}

\textbf{Partie B : Expression écrite (7 points)}
\textit{Sujet :} « Votre correspondant chinois vous demande : \zh{你平时怎么用手机上网？} (Comment utilises-tu ton téléphone pour aller sur Internet d'habitude ?). Répondez-lui par un message WeChat d'environ \textbf{80--100 caractères chinois}. Utilisez au moins \textbf{4 connecteurs} (d'abord, ensuite, parce que, donc, si, en résumé...) et \textbf{5 mots de vocabulaire} de la leçon (微信, 朋友圈, 短视频, 点赞, 视频通话, 红包, 登录...). »

\textbf{Partie C : Traduction (7 points)}
\begin{enumerate}
\item \textbf{Thème (3 pts) :} Traduisez en chinois.
\begin{enumerate}
\item Je n'ai pas l'habitude de poster des photos sur les Moments.
\item Hier, j'ai reçu un gros enveloppe rouge de mon ami chinois.
\item Si le Wi-Fi ne marche pas, je ne peux pas regarder de vidéos.
\end{enumerate}
\item \textbf{Version (4 pts) :} Traduisez en français.
\zh{现在的年轻人离不开手机。上网、刷视频、发消息，已经成为生活的一部分。虽然很方便，但是要注意保护眼睛，控制时间。建议大家多运动，少看屏幕。}
\end{enumerate}
\end{exercice}

\begin{corrige}[Épreuve type Bac — Corrigé détaillé]
\textbf{Partie A : Compréhension}
\begin{enumerate}
\item Il habite à Guangzhou (Canton), il est camerounais (étudiant camerounais \zh{喀麦隆留学生}).
\item Matin : regarder les Moments (\zh{看朋友圈}). Soir : envoyer un vocal à sa mère (\zh{发语音报平安}), appel vidéo avec camarades pour réviser (\zh{视频通话复习功课}).
\item Il a regardé un live sur la culture du thé chinois (\zh{中国茶文化}). Il l'a partagé/transféré à ses camarades camerounais (\zh{转发给了喀麦隆的同学}).
\item « WeChat me rend, ma famille et mes amis plus proches. » / « WeChat rapproche ma famille, mes amis et moi. » (\zh{让} causative, \zh{更近} comparatif « plus proche », \zh{了} changement d'état).
\item Parce qu'il est loin de son pays (étudiant à l'étranger), WeChat lui permet de garder le contact quotidien avec sa famille et ses amis (inférence : \zh{报平安} donner des nouvelles, \zh{更近} raccourcir la distance).
\end{enumerate}

\textbf{Partie B : Expression écrite (Exemple de production attendue)}
\zh{我通常先登录微信，然后看朋友圈，给朋友的照片点赞。有时候我刷短视频，觉得很有意思。因为学习忙，所以我不常看直播。如果周末有空，我会和家人打视频通话。总之，手机很方便，但我要控制时间。}
\emph{(98 caractères. Connecteurs : 通常/先/然后, 有时候, 因为...所以..., 如果...会..., 总之. Vocab : 登录, 微信, 朋友圈, 点赞, 短视频, 直播, 视频通话.)}

\textbf{Grille d'évaluation (Expression écrite) :}
\begin{center}
\begin{tabularx}{\linewidth}{|X|X|X|}
\hline
\rowcolor{popblueL} Critère & Indicateurs & Points \\
\hline
Respect consignes & Longueur (80-100 car.), 4 connecteurs, 5 mots vocab. & 2 \\
\hline
Grammaire / Syntaxe & Ordre des mots correct, structures variées, particules (\zh{了, 的, 得}) justes & 2 \\
\hline
Vocabulaire & Emploi précis mots leçon, pas de répétitions excessives & 1.5 \\
\hline
Cohérence / Logique & Idées liées, connecteurs utilisés à bon escient & 1.5 \\
\hline
\hline
\textbf{Total} & & \textbf{7} \\
\hline
\end{tabularx}
\end{center}

\textbf{Partie C : Traduction}
\begin{enumerate}
\item \textbf{Thème :}
\begin{enumerate}
\item \zh{我不习惯在朋友圈发照片。} \emph{Wǒ bù xíguàn zài péngyouquān fā zhàopiàn.}
\item \zh{昨天，我收到了中国朋友发的一个大红包。} \emph{Zuótiān, wǒ shōudàole Zhōngguó péngyou fā de yí ge dà hóngbāo.} (\zh{收到} recevoir, \zh{了} accompli).
\item \zh{如果网不好 / Wi-Fi 连不上，我就看不了视频。} \emph{Rúguǒ wǎng bù hǎo / Wi-Fi lián bù shàng, wǒ jiù kàn bù liǎo shìpín.} (\zh{看不了} ne pas pouvoir regarder / ne pas réussir à regarder).
\end{enumerate}
\item \textbf{Version :}
« Les jeunes d'aujourd'hui ne peuvent pas se passer de leur téléphone. Aller sur Internet, faire défiler des vidéos, envoyer des messages sont devenus une partie de la vie. Bien que ce soit très pratique, il faut faire attention à protéger ses yeux et contrôler son temps. Je conseille à tout le monde de faire plus de sport et de moins regarder les écrans. »
\begin{itemize}
\item \zh{离不开} (lí bù kāi) = ne pas pouvoir se passer de, être accro à.
\item \zh{成为...一部分} (chéngwéi... yībùfèn) = devenir une partie de...
\item \zh{虽然...但是...} (suīrán... dànshì...) = bien que... mais... (corrélatif obligatoire).
\item \zh{注意保护} (zhùyì bǎohù) = faire attention à protéger / veiller à.
\item \zh{建议} (jiànyì) = suggérer, conseiller.
\end{itemize}
\end{enumerate}
\end{corrige}

\begin{aretenir}[Bilan final de la Leçon 38 — Module 5]
\begin{itemize}
\item \textbf{Culture écriture :} \zh{电} : pictogramme éclair $\rightarrow$ traditionnel \zh{電} (雨+申) $\rightarrow$ simplifié \zh{电} (5 traits, crochet final descendant). Confusions : \zh{申, 由, 甲}.
\item \textbf{Lexique appareils :} Règle \zh{电} + organe/fonction (\zh{话, 视, 脑, 影, 子邮件, 梯}).
\item \textbf{Lexique réseaux sociaux :} \zh{微信, 抖音, 朋友圈, 发/刷/看/点赞/评论/转发/关注/直播/红包/登录/密码/二维码}.
\item \textbf{Grammaire / Connecteurs :} \zh{首先/然后/还有}, \zh{因为...所以...}, \zh{虽然...但是...}, \zh{如果...就...}, \zh{总之}.
\item \textbf{Compétences écrites :} Rédiger un message court (80-100 car.) structuré ; Traduire phrases isolées et court paragraphe (Thème/Version) en respectant l'ordre des mots et les particules.
\item \textbf{Stratégie examen :} Segmenter $\rightarrow$ Structurer $\rightarrow$ Vocabulaire $\rightarrow$ Assembler $\rightarrow$ Relire (mots-outils, classificateurs, tons).
\end{itemize}
\end{aretenir}

\coursec{Radicaux et écriture des caractères du thème}

Dans la section précédente, nous avons vu que le caractère \zh{电} (diàn, électricité) est né d'un dessin : un éclair dans un champ. Cette logique du dessin et de la combinaison n'est pas propre à \zh{电} : c'est le principe général de l'écriture chinoise. Pour écrire correctement les caractères des réseaux sociaux (\zh{话}, \zh{网}, \zh{信}, \zh{发}, \zh{友}…), il faut donc d'abord comprendre leurs \cle{radicaux} (les « clés » qui donnent le sens général) et respecter l'\cle{ordre des traits}. C'est l'objet de cette section : apprendre à voir, décomposer, écrire et mémoriser ces caractères.

\courssub{Observer avant d'écrire : les radicaux du thème}

\begin{exemplebox}[Observation : quatre caractères, un même radical]
Regardez ces quatre caractères et cherchez ce qu'ils ont en commun :\\
\zh{话} (huà, parole) — \zh{说} (shuō, parler) — \zh{语} (yǔ, langue) — \zh{谢} (xiè, remercier)\\
Tous les quatre commencent, à gauche, par le même élément : \zh{讠}. Et tous les quatre ont un rapport avec la \cle{parole}. Ce n'est pas un hasard : \zh{讠} est le radical de la parole.
\end{exemplebox}

\begin{definition}[Le radical 讠, clé de la parole]
\zh{讠} (yán) est la forme \emph{simplifiée} du radical \zh{言} (parole), utilisée quand il est placé à gauche d'un caractère. Il s'écrit en \textbf{2 traits} : un point (丶) en haut, puis un trait horizontal-replié-vertical. Tout caractère contenant \zh{讠} à gauche a un lien avec la parole, la langue ou la communication : \zh{说话} (shuōhuà, parler), \zh{汉语} (Hànyǔ, la langue chinoise), \zh{谢谢} (xièxie, merci).
\end{definition}

\begin{center}
\begin{tabularx}{\linewidth}{|X|X|X|X|}
\hline
\rowcolor{popblueL} Caractères & Pinyin & Nature & Traduction \\
\hline
话 & huà & nom & parole, propos \\
\hline
说 & shuō & verbe & parler, dire \\
\hline
语 & yǔ & nom & langue \\
\hline
谢 & xiè & verbe & remercier \\
\hline
信 & xìn & nom & message, lettre ; confiance \\
\hline
网 & wǎng & nom & filet ; réseau, Internet \\
\hline
发 & fā & verbe & envoyer, émettre, publier \\
\hline
友 & yǒu & nom & ami(e) \\
\hline
\end{tabularx}
\end{center}

\begin{savaistu}[网 : un filet qui attrape le monde]
Le caractère \zh{网} (wǎng) est un \cle{pictogramme} : il dessinait à l'origine un \textbf{filet de pêche}, avec ses mailles croisées à l'intérieur d'un cadre \zh{冂}. Aujourd'hui, il désigne le réseau, et Internet se dit \zh{互联网} (hùliánwǎng), littéralement « filet mutuellement connecté ». Pour un Chinois, Internet est donc un immense filet qui attrape le monde entier — une image que les pêcheurs de Douala ou de Limbe comprendraient très bien ! De même, \zh{网友} (wǎngyǒu), « ami du filet », désigne un ami rencontré en ligne, et \zh{上网} (shàngwǎng), « monter sur le filet », signifie « aller sur Internet ».
\end{savaistu}

\courssub{Décomposer les caractères : radicaux et histoires}

\begin{definition}[Les radicaux 亻 et 冂]
\zh{亻} (rén) est la forme debout, à gauche, du radical « homme » \zh{人}. \zh{信} (xìn, message, confiance) = \zh{亻} homme + \zh{言} parole : « la parole d'un homme est un message fiable ».\\
\zh{冂} est le cadre qui entoure \zh{网} : il représente les bords du filet.\\
\zh{朋} (péng, ami) est formé de deux \zh{月} côte à côte : deux compagnons qui marchent ensemble, comme dans \zh{朋友} (péngyou, ami).
\end{definition}

\begin{popfigure}
\begin{center}
\begin{tikzpicture}[node distance=0.6cm]
\node[draw=popblue, fill=popblueL, rounded corners, minimum width=1.4cm, minimum height=1.2cm, font=\Huge] (xin) {信};
\node[draw=popgreen, fill=popgreenL, rounded corners, minimum width=0.9cm, minimum height=1.2cm, font=\Huge, below left=0.9cm and 0.3cm of xin] (ren) {亻};
\node[draw=poporange, fill=poporangeL, rounded corners, minimum width=0.9cm, minimum height=1.2cm, font=\Huge, below right=0.9cm and 0.3cm of xin] (yan) {言};
\node[below=0.05cm of ren, font=\small, align=center] {homme\\(à gauche)};
\node[below=0.05cm of yan, font=\small, align=center] {parole\\(à droite)};
\draw[-{Stealth}, thick, popdark] (xin) -- (ren);
\draw[-{Stealth}, thick, popdark] (xin) -- (yan);
\node[draw=poppurple, fill=poppurpleL, rounded corners, minimum width=1.4cm, minimum height=1.2cm, font=\Huge, right=4.5cm of xin] (hua) {话};
\node[draw=popgreen, fill=popgreenL, rounded corners, minimum width=0.9cm, minimum height=1.2cm, font=\Huge, below left=0.9cm and 0.3cm of hua] (yi) {讠};
\node[draw=poporange, fill=poporangeL, rounded corners, minimum width=0.9cm, minimum height=1.2cm, font=\Huge, below right=0.9cm and 0.3cm of hua] (she) {舌};
\node[below=0.05cm of yi, font=\small, align=center] {parole\\(radical, 2 traits)};
\node[below=0.05cm of she, font=\small, align=center] {langue\\(6 traits)};
\draw[-{Stealth}, thick, popdark] (hua) -- (yi);
\draw[-{Stealth}, thick, popdark] (hua) -- (she);
\end{tikzpicture}
\end{center}
\legende{Décomposition des caractères : 信 = 亻 + 言 (« la parole d'un homme ») ; 话 = 讠 + 舌 (« la parole de la langue »).}
\end{popfigure}

\begin{methode}[Mémoriser un caractère en 4 étapes]
\begin{enumerate}
\item \textbf{Identifier le radical} : quel élément donne le sens ? (\zh{话} → \zh{讠}, parole.)
\item \textbf{Raconter son histoire} : \zh{信} = « la parole d'un homme est un message fiable » ; \zh{网} = un filet de pêche.
\item \textbf{L'écrire 5 fois en comptant les traits} à voix haute, dans le bon ordre.
\item \textbf{L'utiliser dans une phrase} : \zh{我给你发消息。} (Wǒ gěi nǐ fā xiāoxi. — Je t'envoie un message.)
\end{enumerate}
\end{methode}

\courssub{L'ordre des traits : règles et application}

\begin{propriete}[Règle générale de l'ordre des traits]
En chinois, on écrit : \textbf{de haut en bas}, \textbf{de gauche à droite}, l'\textbf{horizontal avant le vertical}, l'\textbf{extérieur avant l'intérieur}, et on \textbf{ferme le cadre en dernier}. Pour \zh{话} : on écrit d'abord le radical \zh{讠} à gauche (2 traits), puis la partie droite \zh{舌} de haut en bas (6 traits), soit \textbf{8 traits au total}. Pour \zh{网} : on trace le cadre \zh{冂} d'abord, puis les mailles à l'intérieur, soit \textbf{6 traits}.
\end{propriete}

\begin{popfigure}
\begin{center}
\begin{tikzpicture}[scale=0.9]
\foreach \i in {0,...,7}{
  \pgfmathsetmacro{\x}{mod(\i,4)*2.2}
  \pgfmathsetmacro{\y}{-floor(\i/4)*2.2}
  \draw[popline] (\x,\y) rectangle (\x+1.8,\y+1.8);
  \draw[popline, dashed] (\x+0.9,\y) -- (\x+0.9,\y+1.8);
  \draw[popline, dashed] (\x,\y+0.9) -- (\x+1.8,\y+0.9);
  \pgfmathtruncatemacro{\n}{\i+1}
  \node[popmuted, font=\footnotesize] at (\x+0.25,\y+1.55) {\n};
}
\node[font=\Large] at (0.9,0.9) {丶};
\node[font=\Large] at (3.1,0.9) {讠};
\node[font=\Large] at (5.3,0.9) {讠};
\node[font=\Large] at (7.5,0.9) {讠};
\node[font=\Large] at (0.9,-1.3) {话};
\node[font=\Large] at (3.1,-1.3) {话};
\node[font=\Large] at (5.3,-1.3) {话};
\node[font=\Large] at (7.5,-1.3) {话};
\node[align=center, font=\small] at (4.2,-3.1) {Étapes 1-2 : le radical 讠\ (point, puis trait replié) ;\\étapes 3-6 : la partie haute de 舌\ (丿, 一, 一, 丨) ;\\étapes 7-8 : la bouche 口\ en bas. Total : 8 traits.};
\end{tikzpicture}
\end{center}
\legende{Ordre des traits de 话 (8 traits) : d'abord 讠 à gauche, puis 舌 de haut en bas.}
\end{popfigure}

\begin{popfigure}
\begin{center}
\begin{tikzpicture}[scale=0.9]
\foreach \i in {0,...,5}{
  \pgfmathsetmacro{\x}{\i*2.2}
  \draw[popline] (\x,0) rectangle (\x+1.8,1.8);
  \draw[popline, dashed] (\x+0.9,0) -- (\x+0.9,1.8);
  \draw[popline, dashed] (\x,0.9) -- (\x+1.8,0.9);
  \pgfmathtruncatemacro{\n}{\i+1}
  \node[popmuted, font=\footnotesize] at (\x+0.25,1.55) {\n};
}
\node[font=\Large] at (0.9,0.9) {丨};
\node[font=\Large] at (3.1,0.9) {冂};
\node[font=\Large] at (5.3,0.9) {冂};
\node[font=\Large] at (7.5,0.9) {冂};
\node[font=\Large] at (9.7,0.9) {网};
\node[font=\Large] at (11.9,0.9) {网};
\node[align=center, font=\small] at (6.4,-0.7) {Trait 1 : vertical gauche ; trait 2 : le cadre 冂\ se referme ;\\traits 3-4 : première maille 乂 ; traits 5-6 : deuxième maille. Total : 6 traits.};
\end{tikzpicture}
\end{center}
\legende{Ordre des traits de 网 (6 traits) : le cadre d'abord, les mailles du filet ensuite.}
\end{popfigure}

\begin{aretenir}[Comptage des traits des caractères du thème]
\begin{itemize}
\item \zh{话} (huà) : \textbf{8 traits} = \zh{讠} (2) + \zh{舌} (6).
\item \zh{网} (wǎng) : \textbf{6 traits} = cadre (2) + mailles (4).
\item \zh{信} (xìn) : \textbf{9 traits} = \zh{亻} (2) + \zh{言} (7).
\item \zh{发} (fā) : \textbf{5 traits}.
\item \zh{友} (yǒu) : \textbf{4 traits}.
\end{itemize}
\end{aretenir}

\begin{exemplebox}[Les caractères en contexte]
\zh{我给你发消息。} (Wǒ gěi nǐ fā xiāoxi.) « Je t'envoie un message. »\\
\zh{他是我的网友。} (Tā shì wǒ de wǎngyǒu.) « C'est un ami rencontré en ligne. »\\
\zh{我每天晚上发朋友圈。} (Wǒ měitiān wǎnshang fā péngyouquān.) « Je publie sur mon fil chaque soir. »\\
\zh{谢谢你的信！} (Xièxie nǐ de xìn !) « Merci pour ton message ! »
\end{exemplebox}

\courssub{Attention aux caractères proches}

\begin{attention}[话 ou 活 ? Le radical change tout]
Erreur très fréquente des francophones : confondre \zh{话} (huà, parole) et \zh{活} (huó, vivre). Les deux contiennent \zh{舌} ; \textbf{seul le radical change} : \zh{讠} (parole) contre \zh{氵} (eau). Mais le sens est totalement différent : \zh{说话} (shuōhuà) = parler ; \zh{生活} (shēnghuó) = la vie. Autres pièges : \zh{网} (wǎng, filet) ≠ \zh{冈} (gāng, colline, sans les mailles) ; \zh{信} (xìn, message) ≠ \zh{唁} (yàn, condoléances, avec \zh{口}) ; \zh{友} (yǒu, ami) ≠ \zh{反} (fǎn, contraire).
\end{attention}

\begin{center}
\begin{tabularx}{\linewidth}{|X|X|X|X|}
\hline
\rowcolor{popblueL} Paire & Pinyin & Différence & Traduction \\
\hline
话 / 活 & huà / huó & 讠 parole / 氵 eau & parole / vivre \\
\hline
网 / 冈 & wǎng / gāng & avec / sans mailles & filet, réseau / colline \\
\hline
信 / 唁 & xìn / yàn & 亻 homme / 口 bouche & message / condoléances \\
\hline
友 / 反 & yǒu / fǎn & trait final différent & ami / contraire \\
\hline
\end{tabularx}
\end{center}

\begin{exoresolu}[Décomposer et écrire]
Énoncé. 1) Décomposez les caractères \zh{信} et \zh{话} en radicaux et donnez leur « histoire ». 2) Indiquez le nombre de traits de \zh{网} et de \zh{发}. 3) Traduisez : \zh{他是我的网友。}
\tcblower
Solution détaillée. 1) \zh{信} = \zh{亻} (homme) + \zh{言} (parole) : « la parole d'un homme est un message fiable » → message, confiance. \zh{话} = \zh{讠} (parole) + \zh{舌} (langue) : « ce que dit la langue » → parole. 2) \zh{网} a 6 traits (cadre 2 + mailles 4) ; \zh{发} a 5 traits. 3) \zh{他是我的网友。} (Tā shì wǒ de wǎngyǒu.) = « C'est un ami rencontré en ligne » : \zh{他} il, \zh{是} est, \zh{我的} mon, \zh{网友} ami du réseau.
\end{exoresolu}

\begin{exercice}[À vous d'écrire]
1) Écrivez \zh{话} cinq fois en comptant les traits à voix haute. 2) Complétez avec le bon caractère (\zh{话} ou \zh{活}) : \zh{他说中国\_\_。} / \zh{我们在雅温得生\_\_。} 3) Traduisez en français : \zh{我给你发消息。}
\end{exercice}

\begin{corrige}[Exercice « À vous d'écrire »]
1) Vérification : 8 traits, \zh{讠} d'abord puis \zh{舌} de haut en bas. 2) \zh{他说中国话。} (Tā shuō Zhōngguó huà. — Il parle chinois.) / \zh{我们在雅温得生活。} (Wǒmen zài Yǎwēndé shēnghuó. — Nous vivons à Yaoundé.) 3) « Je t'envoie un message » — littéralement « je à toi envoie message », avec la structure Sujet + 给 + personne + verbe.
\end{corrige}

\coursec{Leçon 38 — Expression écrite : origine du caractère 电 et radicaux ; caractères des réseaux sociaux}

\courssub{Origine et évolution du caractère \zh{电} (diàn)}

\begin{textebox}[Évolution graphique : de l'éclair à l'électricité]
\textbf{甲骨文 (Jiǎgǔwén) :} \zh{} (représentation pictographique d'un éclair zigzaguant dans le ciel). \\
\textbf{金文 (Jīnwén) :} \zh{} (l'éclair sort d'un nuage \zh{云} / \zh{雨}). \\
\textbf{小篆 (Xiǎozhuàn) :} \zh{⚡} (structure fixée : \zh{雨} « pluie/nuage » en haut, \zh{申} « étendre/éclair » en bas — \zh{申} donne aussi la prononciation). \\
\textbf{隶书 (Lìshū) :} \zh{⚡} (linéarisation, trait unique pour l'éclair). \\
\textbf{繁体字 (Fántǐzì) :} \zh{電} (radical \zh{雨} yǔ, 13 traits). \\
\textbf{简体字 (Jiǎntǐzì) :} \zh{电} (radical \zh{田} tián, 5 traits — simplification cursive où la partie supérieure \zh{雨} se réduit à \zh{田}).
\end{textebox}

\begin{propriete}[Analyse du caractère simplifié \zh{电}]
\begin{itemize}
    \item \textbf{Radical (部首) :} \zh{田} (tián, « champ »). \emph{Attention : ce n'est pas le radical sémantique d'origine (qui était \zh{雨}), mais le radical de classement dans les dictionnaires modernes.}
    \item \textbf{Composant phonétique :} \zh{申} (shēn, « étendre, expliquer » → ici « éclair qui s'étend »).
    \item \textbf{Sens originel :} « éclair, foudre ». Sens moderne étendu : « électricité », puis par métonymie « télé- » (téléphone, télévision, télégramme).
    \item \textbf{Ordre des traits (5 traits) :} 
    \begin{enumerate}
        \item \zh{丨} (verticale centrale)
        \item \zh{(trait brisé)} (horizontale + crochet bas à droite → forme le haut de \zh{田})
        \item \zh{一} (horizontale médiane)
        \item \zh{一} (horizontale basse → ferme \zh{田})
        \item \zh{乙} (trait courbe final « \zh{乙} » partant du bas de la verticale centrale vers la gauche)
    \end{enumerate}
    \emph{Règle :} Pour \zh{田}, on écrit le contour (traits 1-2-3-4) avant l'intérieur (vide ici), puis le trait ajouté (5).
\end{itemize}
\end{propriete}

\begin{exemplebox}[Famille lexicale de \zh{电} (HSK 1-4)]
\begin{center}
\begin{tabularx}{\linewidth}{|X|X|X|X|}
\hline
\rowcolor{popblueL} \textbf{Caractères} & \textbf{Pinyin} & \textbf{Nature} & \textbf{Traduction} \\
\hline
电话 & diànhuà & n. & téléphone \\
\hline
电脑 & diànnǎo & n. & ordinateur (lit. « cerveau électrique ») \\
\hline
电视 & diànshì & n. & télévision (lit. « vision électrique ») \\
\hline
电影 / 电影 & diànyǐng & n. & cinéma, film (lit. « ombre électrique ») \\
\hline
电子邮件 / 电邮 & diànzǐ yóujiàn / diànyóu & n. & courriel, email \\
\hline
充电 & chōngdiàn & v. & recharger (batterie) \\
\hline
短信 & duǎnxìn & n. & SMS (message court) \\
\hline
\end{tabularx}
\end{center}
\end{exemplebox}

\courssub{Radicaux et écriture des caractères du thème}

Les caractères du vocabulaire « télécommunications / numérique » se construisent autour de radicaux sémantiques transparents. Les identifier permet de mémoriser l'écriture et le sens par familles.

\begin{propriete}[Quatre radicaux clés du module]
\begin{center}
\begin{tabularx}{\linewidth}{|X|X|X|X|}
\hline
\rowcolor{poporangeL} \textbf{Radical (Clé)} & \textbf{Sens du radical} & \textbf{Caractères de la leçon} & \textbf{Explication} \\
\hline
\zh{讠} (yán, 言) & parole, langage, message & \zh{信} (xìn), \zh{讯} (xùn), \zh{课} (kè), \zh{读} (dú) & Côté gauche : « messages échangés » (\zh{信} = personne \zh{人} + parole \zh{言}). \\
\hline
\zh{钅} (jīn, 金) & métal, or, hardware & \zh{针} (zhēn), \zh{钱} (qián), \zh{链} (liàn) & Côté gauche : composants matériels, paiement (\zh{钱} = or + deux traits/tenues). \\
\hline
\zh{扌} (shǒu, 手) & main, action manuelle & \zh{打} (dǎ), \zh{拍} (pāi), \zh{发} (fā), \zh{扫} (sǎo) & Côté gauche : action de la main (\zh{发} = main qui envoie/lance). \\
\hline
\zh{目} (mù) & œil, voir & \zh{视} (shì), \zh{看} (kàn), \zh{眠} (mián) & En haut ou à gauche : vision, vidéo (\zh{视} = voir + montrer \zh{示}). \\
\hline
\end{tabularx}
\end{center}
\end{propriete}

\begin{attention}[Pièges d'écriture : caractères proches (形近字)]
\begin{itemize}
    \item \zh{微} (wēi, micro) \textbf{vs} \zh{徽} (huī, emblème) : \zh{微} a \zh{彳} (pas) en bas à gauche ; \zh{徽} a \zh{心} (cœur).
    \item \zh{信} (xìn, message/confiance) \textbf{vs} \zh{侠} (xiá, chevalier) : \zh{信} = \zh{人} + \zh{言} ; \zh{侠} = \zh{人} + \zh{夹} (phonétique).
    \item \zh{博} (bó, vaste/blog) \textbf{vs} \zh{搏} (bó, lutter) : \zh{博} = \zh{十} + \zh{尃} (phonétique) ; \zh{搏} = \zh{扌} + \zh{尃}.
    \item \zh{视} (shì, vidéo/regarder) \textbf{vs} \zh{试} (shì, essayer) : \zh{视} = \zh{示} + \zh{见} (voir) ; \zh{试} = \zh{讠} + \zh{式} (style/test).
    \item \zh{群} (qún, groupe) \textbf{vs} \zh{裙} (qún, jupe) : \zh{群} = \zh{羊} + \zh{君} (phonétique) ; \zh{裙} = \zh{衤} (vêtement) + \zh{君}.
\end{itemize}
\end{attention}

\begin{experience}[Atelier d'écriture guidée (15 min)]
\begin{enumerate}
    \item \textbf{Dictée de radicaux :} L'enseignant dicte : « message », « ordinateur », « vidéo », « groupe », « payer ». Les élèves écrivent le radical correspondant (\zh{讠}, \zh{钅}, \zh{目}, \zh{羊}, \zh{钅}) sur ardoise.
    \item \textbf{Reconstitution :} À partir des composants donnés (\zh{人}, \zh{言}, \zh{微}, \zh{博}, \zh{扌}, \zh{发}, \zh{目}, \zh{示}, \zh{见}), recomposer \zh{信}, \zh{微}, \zh{博}, \zh{发}, \zh{视}.
    \item \textbf{Chasse aux intrus :} Dans une liste de 10 caractères (ex: \zh{电, 微, 信, 博, 抖, 音, 网, 站, 群, 裙}), identifier celui qui n'appartient pas au champ lexical « réseaux sociaux / télécom » (\zh{裙}) et justifier par le radical.
\end{enumerate}
\end{experience}

\courssub{Vocabulaire des réseaux sociaux}

Après avoir étudié l'origine du caractère \zh{电} et les radicaux qui structurent le vocabulaire des télécommunications (\zh{讠}, \zh{钅}, \zh{扌}, \zh{目}), nous disposons maintenant des clés pour comprendre l'écriture des mots du monde numérique. Cette section propose une acquisition systématique du lexique des réseaux sociaux, organisée par champs sémantiques, avec une attention particulière aux collocations verbales et aux pièges d'interférence du français.

\courssub{Plateformes et applications principales}

Observons d'abord les noms des applications les plus utilisées en Chine. Remarquez que beaucoup sont des mots composés transparents : \zh{微} (micro) + \zh{信} (message/lettre), \zh{微} + \zh{博} (micro-blog), \zh{抖} (secouer/vibrer) + \zh{音} (son).

\begin{center}
\begin{tabularx}{\linewidth}{|X|X|X|X|}
\hline
\rowcolor{popblueL} \textbf{Caractères} & \textbf{Pinyin} & \textbf{Nature} & \textbf{Traduction} \\
\hline
微信 & Wēixìn & n. & WeChat (application « tout-en-un » : messagerie, paiement, réseau social) \\
\hline
微博 & Wēibó & n. & Weibo (micro-blog style Twitter/X) \\
\hline
抖音 & Dǒuyīn & n. & Douyin (version chinoise de TikTok, vidéos courtes) \\
\hline
QQ & QQ & n. & QQ (messagerie historique de Tencent, encore très utilisée par les étudiants) \\
\hline
网站 & wǎngzhàn & n. & site web \\
\hline
应用 / 应用程序 & yìngyòng / yìngyòng chéngxù & n. & application (appli) \\
\hline
\end{tabularx}
\end{center}

\begin{aretenir}[Équivalences utiles pour la traduction]
\begin{itemize}
    \item \textbf{WhatsApp} $\approx$ \zh{微信} (Wēixìn) — \emph{Mais WeChat est bien plus qu'une messagerie : c'est un écosystème (paiement, mini-programmes, administration).}
    \item \textbf{TikTok (international)} $\approx$ \zh{抖音} (Dǒuyīn) — \emph{Contenu et algorithme distincts de la version internationale.}
    \item \textbf{« Liker / mettre un like »} = \zh{点赞} (diǎnzàn, litt. « cliquer louer/approuver »).
    \item \textbf{« Suivre / s'abonner »} = \zh{关注} (guānzhù).
\end{itemize}
\end{aretenir}

\courssub{Actions courantes sur les réseaux sociaux}

En chinois, les actions sur Internet s'expriment souvent par des structures \textbf{Verbe + Objet} (VO) ou \textbf{Préposition + Lieu + Verbe}. La particule \zh{上} (shàng, « sur ») indique le support numérique.

\begin{center}
\begin{tabularx}{\linewidth}{|X|X|X|X|}
\hline
\rowcolor{poporangeL} \textbf{Caractères} & \textbf{Pinyin} & \textbf{Nature} & \textbf{Traduction} \\
\hline
上网 & shàngwǎng & v. (VO) & aller sur Internet, se connecter (litt. « monter sur le filet ») \\
\hline
发消息 & fā xiāoxi & v. (VO) & envoyer un message \\
\hline
发朋友圈 & fā péngyouquān & v. (VO) & publier un statut / poster sur son fil (Moments) \\
\hline
点赞 & diǎnzàn & v. (VO) & liker, aimer (litt. « cliquer-approuver ») \\
\hline
评论 & pínglùn & v. / n. & commenter / un commentaire \\
\hline
分享 & fēnxiǎng & v. & partager \\
\hline
关注 & guānzhù & v. & suivre, s'abonner à (un compte) \\
\hline
聊天 & liáotiān & v. (VO) & discuter, chatter (litt. « bavarder-ciel ») \\
\hline
看视频 & kàn shìpín & v. (VO) & regarder des vidéos \\
\hline
刷视频 & shuā shìpín & v. (VO) & scroller des vidéos (action rapide, passive) \\
\hline
\end{tabularx}
\end{center}

\begin{propriete}[Collocations à mémoriser : Verbe + Support / Verbe + Objet]
Les prépositions \zh{在} (zài) et \zh{用} (yòng) structurent l'action technique.
\begin{center}
\begin{tabularx}{\linewidth}{|X|X|X|}
\hline
\rowcolor{popgreenL} \textbf{Structure} & \textbf{Exemple (Caractères + Pinyin)} & \textbf{Traduction} \\
\hline
用 + Outil + Verbe & \zh{用手机上网} (yòng shǒujī shàngwǎng) & aller sur Internet \emph{avec son téléphone} \\
\hline
在 + Plateforme + 上 + Verbe & \zh{在微信上聊天} (zài Wēixìn shàng liáotiān) & discuter \emph{sur} WeChat \\
\hline
给 + Personne + Verbe & \zh{给朋友点赞} (gěi péngyou diǎnzàn) & liker \emph{un ami} (son contenu) \\
\hline
发 + Contenu & \zh{发照片} (fā zhàopiàn) / \zh{发视频} (fā shìpín) & poster une photo / une vidéo \\
\hline
\end{tabularx}
\end{center}
\end{propriete}

\begin{exemplebox}[Exemples d'actions courantes]
\begin{enumerate}
    \item \zh{我每天用手机上网。} (Wǒ měitiān yòng shǒujī shàngwǎng.) — Je vais sur Internet avec mon téléphone chaque jour.
    \item \zh{她在微博上发了一个视频。} (Tā zài Wēibó shàng fā le yí gè shìpín.) — Elle a posté une vidéo sur Weibo.
    \item \zh{请给我的照片点赞！} (Qǐng gěi wǒ de zhàopiàn diǎnzàn !) — Like ma photo, s'il te plaît !
    \item \zh{我们在QQ群里聊天。} (Wǒmen zài QQ qún lǐ liáotiān.) — Nous discutons dans le groupe QQ.
\end{enumerate}
\end{exemplebox}

\courssub{Personnes, contenus et groupes}

Le vocabulaire suivant permet de décrire les acteurs et les objets de l'interaction en ligne.

\begin{center}
\begin{tabularx}{\linewidth}{|X|X|X|X|}
\hline
\rowcolor{poppurpleL} \textbf{Caractères} & \textbf{Pinyin} & \textbf{Nature} & \textbf{Traduction} \\
\hline
网友 & wǎngyǒu & n. & ami(e) interne, internaute (litt. « ami du filet ») \\
\hline
朋友圈 & péngyouquān & n. & cercle d'amis = fil d'actualité / « Mur » (WeChat) \\
\hline
视频 & shìpín & n. & vidéo \\
\hline
照片 & zhàopiàn & n. & photo \\
\hline
群 & qún & n. & groupe (de discussion) ; classif. pour les groupes \\
\hline
大群 & dàqún & n. & grand groupe (souvent > 100 personnes) \\
\hline
红包 & hóngbāo & n. & enveloppe rouge (cadeau d'argent numérique) \\
\hline
直播 & zhíbō & n. / v. & direct (live) / faire un direct \\
\hline
主播 & zhǔbō & n. & streamer, animateur de direct \\
\hline
\end{tabularx}
\end{center}

\begin{definition}[朋友圈  : le « Mur » chinois]
\zh{朋友圈} (péngyouquān) se compose de \zh{朋友} (ami) + \zh{圈} (cercle). C'est l'équivalent exact du \emph{Fil d'actualité} (News Feed) ou du \emph{Mur} sur Facebook. Seuls les contacts acceptés (bidirectionnel sur WeChat) voient les publications. On dit \zh{发朋友圈} (publier dedans) et \zh{看朋友圈} (le consulter). Attention : ce n'est pas un « groupe » (\zh{群}), c'est un espace personnel agrégé.
\end{definition}

\begin{textebox}[Texte support : Mon quotidien numérique]
\zh{我每天用微信跟朋友聊天，有时候也发朋友圈。} \\
\zh{Wǒ měitiān yòng Wēixìn gēn péngyou liáotiān, yǒu shíhou yě fā péngyouquān.} \\
\zh{每天早上，我先看微信消息，再刷抖音看视频。周末我在大群里分享照片，大家都给我点赞、评论。网络很方便，但不能花太多时间。} \\
\\
\textbf{Traduction :} Chaque jour, j'utilise WeChat pour discuter avec mes amis, et parfois je publie aussi sur mon fil. Chaque matin, je regarde d'abord les messages WeChat, puis je scrolle Douyin pour voir des vidéos. Le week-end, je partage des photos dans le grand groupe, tout le monde me like et commente. Internet est pratique, mais on ne doit pas y passer trop de temps.
\\
\textbf{Glossaire clé :} \zh{跟...聊天} (gēn... liáotiān : discuter avec...), \zh{先...再...} (xiān... zài... : d'abord... puis...), \zh{大家} (dàjiā : tout le monde), \zh{但} (dàn : mais), \zh{不能} (bù néng : ne pas pouvoir / ne pas devoir).
\end{textebox}

\courssub{Fréquence, durée, avantages et risques}

Pour exprimer une opinion nuancée (compétence visée : expression écrite), il faut maîtriser les indicateurs de fréquence, la structure de la durée \zh{花时间} et les adjectifs d'évaluation.

\begin{center}
\begin{tabularx}{\linewidth}{|X|X|X|X|}
\hline
\rowcolor{poppinkL} \textbf{Caractères} & \textbf{Pinyin} & \textbf{Nature} & \textbf{Traduction} \\
\hline
每天 & měitiān & adv. & chaque jour, tous les jours \\
\hline
有时候 / 有时 & yǒu shíhou / yǒushí & adv. & parfois, de temps en temps \\
\hline
经常 & jīngcháng & adv. & souvent, régulièrement \\
\hline
花时间 / 花...时间 & huā shíjiān / huā... shíjiān & v. (VO) & passer du temps / passer ... temps (structure : 花 + Durée + 在 + Action / 花 + Durée + Verb) \\
\hline
方便 & fāngbiàn & adj. & pratique, commode \\
\hline
有意思 & yǒu yìsi & adj. & intéressant, amusant (litt. « avoir du sens ») \\
\hline
浪费时间 & làngfèi shíjiān & v. (VO) & perdre du temps / gaspiller du temps \\
\hline
影响学习 & yǐngxiǎng xuéxí & v. (VO) & nuire aux études, affecter les études \\
\hline
上瘾 & shàngyǐn & v. / adj. & être accro, addictif \\
\hline
近视 & jìnshì & n. / adj. & myopie / myope \\
\hline
\end{tabularx}
\end{center}

\begin{propriete}[Structure de la durée : \zh{花} (huā) + Durée + \zh{在} (zài) + V / + V + \zh{上} (shàng)]
Pour dire « passer du temps À faire qqch », le chinois utilise le verbe \zh{花} (dépenser).
\begin{center}
\begin{tabularx}{\linewidth}{|X|X|X|}
\hline
\rowcolor{popblueL} \textbf{Chinois (Caractères + Pinyin)} & \textbf{Structure} & \textbf{Traduction} \\
\hline
\zh{他花两个小时上网。} (Tā huā liǎng gè xiǎoshí shàngwǎng.) & Sujet + 花 + Durée + Verbe & Il passe deux heures sur Internet. \\
\hline
\zh{他在手机上花了三个小时。} (Tā zài shǒujī shàng huā le sān gè xiǎoshí.) & Sujet + 在 + Support + 花 + Durée & Il a passé trois heures sur son téléphone. \\
\hline
\zh{不要在游戏上花太多时间。} (Búyào zài yóuxì shàng huā tài duō shíjiān.) & Négatif + 在 + Objet + 花 + Durée & Ne passe pas trop de temps sur les jeux. \\
\hline
\end{tabularx}
\end{center}
\emph{Note :} Le classifieur pour l'heure est \zh{个} (gè) pour \zh{小时} (xiǎoshí). \zh{两} (liǎng) s'utilise pour « deux » devant un classifieur, pas \zh{二} (èr).
\end{propriete}

\begin{exemplebox}[Avantages et risques : phrases modèles]
\begin{enumerate}
    \item \zh{微信很方便，可以发消息、付钱、叫车。} (Wēixìn hěn fāngbiàn, kěyǐ fā xiāoxi, fù qián, jiào chē.) — WeChat est très pratique : on peut envoyer des messages, payer, appeler un taxi.
    \item \zh{看短视频很有意思，但是容易上瘾。} (Kàn duǎn shìpín hěn yǒu yìsi, dànshì róngyì shàngyǐn.) — Regarder des vidéos courtes est amusant, mais c'est facile d'y devenir accro.
    \item \zh{他每天花四个小时刷抖音，浪费时间，影响学习。} (Tā měitiān huā sì gè xiǎoshí shuā Dǒuyīn, làngfèi shíjiān, yǐngxiǎng xuéxí.) — Il passe 4h/jour à scroller Douyin, il perd son temps et nuit à ses études.
    \item \zh{上网太久对眼睛不好，容易近视。} (Shàngwǎng tài jiǔ duì yǎnjing bù hǎo, róngyì jìnshì.) — Rester trop longtemps en ligne est mauvais pour les yeux, risque de myopie.
\end{enumerate}
\end{exemplebox}

\begin{attention}[Erreur fréquente : « Ouvrir Internet »]
\begin{itemize}
    \item \textbf{Français :} « J'ouvre Internet / J'allume Internet. »
    \item \textbf{Chinois incorrect :} \zh{*开网} (kāi wǎng) ou \zh{*打开网} (dǎkāi wǎng).
    \item \textbf{Chinois correct :} \zh{上网} (shàngwǎng) — Verbe intransitif (VO figé) signifiant « accéder au réseau / se connecter ».
    \item \textbf{Explication :} \zh{上} (monter) + \zh{网} (filet/réseau). L'image est celle de « monter à bord du filet ». On ne « ouvre » pas le réseau, on y « monte ».
    \item \textbf{Autres collocations :} \zh{上课} (aller en cours), \zh{上班} (aller au travail), \zh{上车} (monter en voiture).
\end{itemize}
\end{attention}

\begin{savaistu}[Cameroun vs Chine : l'écosystème applicatif]
Au \textbf{Cameroun}, \textbf{WhatsApp} (propriété de Meta) domine outrageusement la communication quotidienne (texte, audio, appels, groupes de classe, associations, commerce informel). Facebook et Messenger restent très présents. Le paiement mobile (Mobile Money : MTN MoMo, Orange Money) fonctionne \emph{à côté} de WhatsApp, pas \emph{dedans}.

En \textbf{Chine}, \textbf{WeChat (\zh{微信})} est une « \textbf{super-application} » (\zh{超级应用}) : elle intègre la messagerie, le réseau social (\zh{朋友圈}), le paiement universel (\zh{微信支付} WeChat Pay — utilisé par tous, des centres commerciaux aux vendeurs de rue), la commande de taxi (\zh{滴滴} Didi via mini-programme), la prise de RDV médical, le paiement des factures, les démarches administratives, etc. Un Chinois peut vivre des semaines sans sortir ni carte bancaire ni espèces, seulement son téléphone avec WeChat. \zh{抖音} (Douyin) domine la vidéo courte et le « social commerce » (achat en direct \zh{直播带货}), bien plus intégré que TikTok Shop en Europe/Afrique.
\end{savaistu}

\courssub{Carte mentale lexicale : 社交媒体}

La figure ci-dessous synthétise les quatre piliers du vocabulaire de cette leçon. Utilisez-la pour réviser les familles de mots et préparer vos productions écrites.

\begin{popfigure}
\begin{center}\tcbox[colback=popblue,colframe=popblue,coltext=white,arc=9pt,boxsep=4pt,left=6pt,right=6pt]{\bfseries\small 社交媒体 (Shèjiāo méitǐ)}\end{center}
\vspace{-2pt}
\begin{tcbraster}[raster columns=2,raster equal height=rows,raster column skip=6pt,raster row skip=6pt,enhanced,blankest]
\begin{tcolorbox}[enhanced,colback=white,colframe=poporange,boxrule=0.9pt,arc=3pt,title={Plateformes 平台},fonttitle=\bfseries\scriptsize,coltitle=white,colbacktitle=poporange,left=4pt,right=4pt,top=2pt,bottom=2pt,boxsep=1pt,toptitle=1.5pt,bottomtitle=1.5pt,halign=left]
\begin{itemize}[nosep,leftmargin=*,itemsep=1pt,font=\scriptsize]
\item 微信 Wēixìn
\item 微博 Wēibó
\item 抖音 Dǒuyīn
\item QQ
\end{itemize}
\end{tcolorbox}
\begin{tcolorbox}[enhanced,colback=white,colframe=popgreen,boxrule=0.9pt,arc=3pt,title={Actions 动作},fonttitle=\bfseries\scriptsize,coltitle=white,colbacktitle=popgreen,left=4pt,right=4pt,top=2pt,bottom=2pt,boxsep=1pt,toptitle=1.5pt,bottomtitle=1.5pt,halign=left]
\begin{itemize}[nosep,leftmargin=*,itemsep=1pt,font=\scriptsize]
\item 上网 shàngwǎng
\item 发消息 fā xiāoxi
\item 点赞 diǎnzàn
\item 评论 pínglùn
\item 分享 fēnxiǎng
\item 关注 guānzhù
\item 刷视频 shuā shìpín
\end{itemize}
\end{tcolorbox}
\begin{tcolorbox}[enhanced,colback=white,colframe=poppurple,boxrule=0.9pt,arc=3pt,title={Personnes/Contenu 人与内容},fonttitle=\bfseries\scriptsize,coltitle=white,colbacktitle=poppurple,left=4pt,right=4pt,top=2pt,bottom=2pt,boxsep=1pt,toptitle=1.5pt,bottomtitle=1.5pt,halign=left]
\begin{itemize}[nosep,leftmargin=*,itemsep=1pt,font=\scriptsize]
\item 网友 wǎngyǒu
\item 朋友圈 péngyouquān
\item 视频 shìpín
\item 照片 zhàopiàn
\item 群 qún
\item 主播 zhǔbō
\end{itemize}
\end{tcolorbox}
\begin{tcolorbox}[enhanced,colback=white,colframe=poppink,boxrule=0.9pt,arc=3pt,title={Opinions/Impact 评价与影响},fonttitle=\bfseries\scriptsize,coltitle=white,colbacktitle=poppink,left=4pt,right=4pt,top=2pt,bottom=2pt,boxsep=1pt,toptitle=1.5pt,bottomtitle=1.5pt,halign=left]
\begin{itemize}[nosep,leftmargin=*,itemsep=1pt,font=\scriptsize]
\item 方便 fāngbiàn
\item 有意思 yǒu yìsi
\item 浪费时间 làngfèi shíjiān
\item 影响学习 yǐngxiǎng xuéxí
\item 上瘾 shàngyǐn
\item 近视 jìnshì
\end{itemize}
\end{tcolorbox}
\end{tcbraster}

\legende{Carte mentale du vocabulaire des réseaux sociaux (HSK 3-4). Branches : Plateformes, Actions, Personnes/Contenu, Opinions/Impact.}
\end{popfigure}

\courssub{Modèle de texte commenté et connecteurs}

\begin{propriete}[Connecteurs logiques pour l'expression écrite (Niveau Première)]
Ces connecteurs permettent d'articuler un paragraphe cohérent (narration, argumentation, concession).
\begin{center}
\begin{tabularx}{\linewidth}{|X|X|X|}
\hline
\rowcolor{popblueL} \textbf{Connecteur (Caractères + Pinyin)} & \textbf{Fonction} & \textbf{Exemple d'emploi} \\
\hline
\zh{先...再...} (xiān... zài...) & Séquence temporelle & \zh{我先看消息，再刷视频。} \\
\hline
\zh{虽然...但是/可是...} (suīrán... dànshì/kěshì...) & Concession / Opposition & \zh{虽然微信方便，但是别上瘾。} \\
\hline
\zh{因为...所以...} (yīnwèi... suǒyǐ...) & Cause / Conséquence & \zh{因为影响学习，所以我控制时间。} \\
\hline
\zh{而且} (érqiě) & Addition (renforcement) & \zh{又快而且方便。} \\
\hline
\zh{另外} (lìngwài) & Ajout d'un argument & \zh{另外，保护眼睛也很重要。} \\
\hline
\zh{因此} (yīncǐ) & Conséquence (formel) & \zh{因此，我们要注意休息。} \\
\hline
\end{tabularx}
\end{center}
\end{propriete}

\begin{textebox}[Texte modèle commenté : « Mon rapport aux écrans »]
\zh{虽然手机很方便，可以上网、聊天、看视频，但是上网太久对眼睛不好，容易近视。} \\
\zh{Suīrán shǒujī hěn fāngbiàn, kěyǐ shàngwǎng, liáotiān, kàn shìpín, dànshì shàngwǎng tài jiǔ duì yǎnjing bù hǎo, róngyì jìnshì.} \\
\zh{因为我要准备考试，所以我决定每天只花一个小时刷抖音。} \\
\zh{Yīnwèi wǒ yào zhǔnbèi kǎoshì, suǒyǐ wǒ juédìng měitiān zhǐ huā yí gè xiǎoshí shuā Dǒuyīn.} \\
\zh{先做完作业，再上网，这样既不耽误学习，也能放松。} \\
\zh{Xiān zuò wán zuòyè, zài shàngwǎng, zhèyàng jì bù dānwù xuéxí, yě néng fàngsōng.} \\
\\
\textbf{Traduction :} Bien que le téléphone soit pratique (Internet, chat, vidéos), rester trop longtemps en ligne est mauvais pour les yeux (myopie). Comme je dois préparer les examens, je décide de ne scroller Douyin qu'une heure par jour. D'abord finir les devoirs, puis Internet : ainsi ni les études ne sont retardées, ni la détente n'est oubliée. \\
\textbf{Analyse structurelle :} 1. Concession (\zh{虽然...但是...}) pour nuancer. 2. Cause/Effet (\zh{因为...所以...}) pour justifier une décision. 3. Séquence + Correlative (\zh{先...再...} / \zh{既...也...}) pour proposer une solution équilibrée.
\end{textebox}

\courssub{Thème et version guidés}

\begin{exoresolu}[Thème guidé : Français $\rightarrow$ Chinois (Caractères + Pinyin)]
\textbf{Énoncé :}
\begin{enumerate}
    \item Je passe souvent deux heures sur WeChat pour discuter avec des amis en ligne.
    \item Cette vidéo est très intéressante, je veux la partager sur mon fil (Moments).
    \item Ne passe pas trop de temps à scroller des vidéos, c'est mauvais pour les yeux.
    \item Il a liké ma photo et m'a suivi (abonné).
\end{enumerate}

\tcblower

\textbf{Solution détaillée :}

\begin{enumerate}
    \item \textbf{Analyse :} « Je » = \zh{我} ; « souvent » = \zh{经常} ; « passer deux heures » = \zh{花两个小时} ; « sur WeChat » = \zh{在微信上} ; « discuter avec » = \zh{跟...聊天} ; « amis en ligne » = \zh{网友}.\par
    \zh{我经常在微信上花两个小时跟网友聊天。} (Wǒ jīngcháng zài Wēixìn shàng huā liǎng gè xiǎoshí gēn wǎngyǒu liáotiān.)\par
    \emph{Ordre : Sujet + Adv(Fréq) + Préposition(Lieu) + Verbe Durée + Verbe Principal + Objet.}
    
    \item \textbf{Analyse :} « Cette vidéo » = \zh{这个视频} ; « très intéressant » = \zh{很有意思} ; « je veux » = \zh{我想} ; « la partager » = \zh{分享它} (souvent omis) / \zh{发到朋友圈} (poster sur Moments).\par
    \zh{这个视频很有意思，我想发到朋友圈。} (Zhè gè shìpín hěn yǒu yìsi, wǒ xiǎng fā dào péngyouquān.)\par
    \emph{Structure : \zh{发到} (fā dào) = poster vers / sur. Plus naturel que \zh{分享到朋友圈} à l'oral.}
    
    \item \textbf{Analyse :} Impératif négatif = \zh{别} / \zh{不要} ; « scroller des vidéos » = \zh{刷视频} ; « trop de temps » = \zh{太多时间} ; « c'est mauvais pour les yeux » = \zh{对眼睛不好}.\par
    \zh{别花太多时间刷视频，对眼睛不好。} (Bié huā tài duō shíjiān shuā shìpín, duì yǎnjing bù hǎo.)\par
    \emph{Note : \zh{别} plus courant à l'oral pour l'impératif négatif. \zh{花时间做某事} structure VO.}
    
    \item \textbf{Analyse :} « Il » = \zh{他} ; « a liké » = \zh{给...点赞了} (aspect accompli \zh{了}) ; « ma photo » = \zh{我的照片} ; « et » = \zh{还} / \zh{又} ; « suivi » = \zh{关注了}.\par
    \zh{他给我的照片点赞了，还关注了我。} (Tā gěi wǒ de zhàopiàn diǎnzàn le, hái guānzhù le wǒ.)\par
    \emph{Structure : \zh{给 + Objet + 点赞} ; \zh{关注} prend l'objet direct (la personne). Double \zh{了} pour deux actions accomplies successives.}
\end{enumerate}
\end{exoresolu}

\courssub{Entraînement à l'examen et grille d'évaluation}

\begin{methode}[Méthode : Réussir l'expression écrite (Thème / Production) en 4 étapes]
\begin{enumerate}
    \item \textbf{Analyser le sujet :} Identifier le temps (présent, passé \zh{了}, futur \zh{要/会}), le registre (familier \zh{别} vs formel \zh{请勿}), les points de grammaire imposés (comparatif \zh{比}, durée \zh{花}, concession \zh{虽然}).
    \item \textbf{Construire le squelette :} Rédiger les idées principales en français, puis les traduire mot-à-mot en structures chinoises (SVO, Compléments avant le verbe). Vérifier l'ordre : \textbf{Sujet - Temps - Lieu - Manière - Verbe - Objet - Particules}.
    \item \textbf{Rédiger en caractères :} Écrire la phrase complète. Vérifier : radicaux, traits, classificateurs (\zh{个, 条, 个}), mots-outils (\zh{的, 得, 地, 了, 过}).
    \item \textbf{Relire (Auto-correction) :} \textit{Check-list :} Tons notés ? Ponctuation chinoise ? \zh{了} position (V+了 vs Phrase+了) ? \zh{在...上} vs \zh{用...} ? \zh{给} pour le bénéficiaire ? Connecteurs logiques présents ?
\end{enumerate}
\end{methode}

\begin{aretenir}[Grille d'évaluation simplifiée — Expression écrite (sur 20 pts)]
\begin{center}
\begin{tabularx}{\linewidth}{|X|X|X|}
\hline
\rowcolor{popgoldL} \textbf{Critère} & \textbf{Indicateurs (Niveau Première A)} & \textbf{Points} \\
\hline
\textbf{Respect de la consigne} & Longueur (50-80 car.), sujet traité, registre adapté & /4 \\
\hline
\textbf{Vocabulaire} & Utilisation correcte de 8+ mots du thème (plateformes, actions, opinions), variété & /4 \\
\hline
\textbf{Grammaire / Structures} & Ordre des mots (COV), \zh{花} durée, \zh{在...上}, \zh{给} VO, connecteurs (\zh{虽然/因为/先...再...}), aspects (\zh{了/过}) & /6 \\
\hline
\textbf{Caractères / Écriture} & Caractères corrects (pas de \zh{开网}, pas de confusion \zh{微/徽}, \zh{视/试}), ponctuation pleine chasse & /3 \\
\hline
\textbf{Cohérence / Style} & Enchaînement logique, phrases liées, ton naturel (pas de traduction mot-à-mot française) & /3 \\
\hline
\end{tabularx}
\end{center}
\end{aretenir}

\begin{exercice}[Mise en application — Simulation d'examen (Durée : 30 min)]
\begin{enumerate}
    \item \textbf{Vocabulaire (4 pts) :} Complétez en caractères et pinyin.
    \begin{itemize}
        \item \zh{\_\_\_\_\_\_} (shàngwǎng) : aller sur Internet.
        \item \zh{\_\_\_\_\_\_} (diǎnzàn) : liker.
        \item \zh{\_\_\_\_\_\_} (péngyouquān) : fil d'actualité WeChat.
        \item \zh{\_\_\_\_\_\_} (làngfèi shíjiān) : perdre du temps.
    \end{itemize}
    \item \textbf{Traduction — Version (6 pts) :} Traduisez en français.
    \begin{enumerate}
        \item 我妈妈每天在微信群里发照片。
        \item 这个主播很有意思，我关注他了。
        \item 别在手机上花太多时间，影响学习。
    \end{enumerate}
    \item \textbf{Traduction — Thème (6 pts) :} Traduisez en chinois (caractères + pinyin).
    \begin{enumerate}
        \item Parfois, je regarde des directs (live) sur Douyin.
        \item Nous avons partagé la vidéo dans le groupe de la classe.
        \item Mon ami est accro aux jeux en ligne, il joue chaque jour.
    \end{enumerate}
    \item \textbf{Production écrite (4 pts) :} Écrivez un court paragraphe (5-6 phrases, 60-80 caractères) pour décrire vos habitudes sur les réseaux sociaux (plateforme, fréquence, actions, opinion). Utilisez \textbf{au moins 8 mots} du vocabulaire de cette section et les connecteurs \zh{先...再...}, \zh{虽然...但是...}, \zh{因为...所以...}.
\end{enumerate}
\end{exercice}

\begin{corrige}[Exercice — Vocabulaire des réseaux sociaux]
\begin{enumerate}
    \item \zh{上网} (shàngwǎng) ; \zh{点赞} (diǎnzàn) ; \zh{朋友圈} (péngyouquān) ; \zh{浪费时间} (làngfèi shíjiān).
    \item \textbf{Version :}
        \begin{enumerate}
            \item Ma mère poste des photos dans le groupe WeChat chaque jour.
            \item Ce streamer est très intéressant, je l'ai suivi (abonné).
            \item Ne passe pas trop de temps sur ton téléphone, cela nuit aux études.
        \end{enumerate}
    \item \textbf{Thème :}
        \begin{enumerate}
            \item \zh{有时候我在抖音看直播。} (Yǒu shíhou wǒ zài Dǒuyīn kàn zhíbō.)
            \item \zh{我们在班级群里分享了视频。} (Wǒmen zài bānjí qún lǐ fēnxiǎng le shìpín.)
            \item \zh{我朋友上网玩游戏上瘾了，每天都玩。} (Wǒ péngyou shàngwǎng wán yóuxì shàngyǐn le, měitiān dōu wán.)
        \end{enumerate}
    \item \textbf{Production écrite (Modèle de réponse — 72 caractères) :}
    \zh{我经常用微信跟网友聊天，每天大概花一个小时。虽然微信很方便，可以发朋友圈、发红包，但是刷视频太久对眼睛不好。所以我现在控制时间，不想浪费时间，因为我要好好学习。} \\
    (Wǒ jīngcháng yòng Wēixìn gēn wǎngyǒu liáotiān, měitiān dàgài huā yí gè xiǎoshí. Suīrán Wēixìn hěn fāngbiàn, kěyǐ fā péngyouquān, fā hóngbāo, dànshì shuā shìpín tài jiǔ duì yǎnjing bù hǎo. Suǒyǐ wǒ xiànzài kòngzhì shíjiān, bù xiǎng làngfèi shíjiān, yīnwèi wǒ yào hǎohǎo xuéxí.) \\
    \emph{Traduction :} Je discute souvent avec des amis en ligne sur WeChat, j'y passe environ une heure par jour. Bien que WeChat soit pratique (publier sur Moments, envoyer des enveloppes rouges), scroller des vidéos trop longtemps est mauvais pour les yeux. Donc maintenant je contrôle mon temps, je ne veux pas perdre de temps, car je dois bien étudier. \\
    \emph{Validation grille :} Vocab (微信, 网友, 聊天, 花时间, 方便, 朋友圈, 红包, 刷视频, 眼睛, 浪费时间, 学习) = 11 mots. Connecteurs : \zh{虽然...但是...}, \zh{所以}, \zh{因为...}. Grammaire : \zh{用...}, \zh{跟...聊天}, \zh{花...时间}, \zh{对...不好}, \zh{控制}, \zh{不想}, \zh{要}. Caractères corrects.
\end{enumerate}
\end{corrige}

\coursec{Modèle de texte commenté et connecteurs}

Après avoir acquis le vocabulaire clé des réseaux sociaux (noms d'applications, verbes d'action, adjectifs d'appréciation), nous passons à l'étape cruciale de la \textbf{production écrite} : assembler ces briques lexicales dans une structure cohérente, fluide et grammaticalement correcte. En chinois, comme en français, un bon paragraphe ne se limite pas à une suite de phrases ; il repose sur une \textbf{architecture logique} (Introduction – Développement – Conclusion) et sur des \textbf{connecteurs} qui guident le lecteur. Cette section vous donne le modèle-type, décortique les outils de liaison et vous fournit la méthode pour rédiger votre propre texte sur ce thème.

\courssub{Le texte modèle : mon usage des réseaux sociaux}

Observons d'abord un texte complet, écrit par une élève camerounaise fictive, Marie. Il respecte parfaitement la structure en trois parties attendue à l'examen (HSK 3 / Bac LV2) : présentation de l'outil et de la fréquence, description des activités concrètes, opinion nuancée.

\begin{textebox}[Modèle : Mon usage des réseaux sociaux (Marie, élève à Yaoundé)]
我叫玛丽，是雅温得一所中学的学生。\\
Wǒ jiào Mǎlì, shì Yǎwēndé yī suǒ zhōngxué de xuésheng.\\
Je m'appelle Marie, je suis élève dans un collège de Yaoundé.

我每天都用手机上网，因为我喜欢跟朋友聊天。\\
Wǒ měitiān dōu yòng shǒujī shàngwǎng, yīnwèi wǒ xǐhuan gēn péngyou liáotiān.\\
Chaque jour, j'utilise mon téléphone pour aller sur Internet, car j'aime discuter avec mes amis.

我也在微信上发照片，朋友们常常给我点赞。\\
Wǒ yě zài Wēixìn shàng fā zhàopiàn, péngyoumen chángcháng gěi wǒ diǎnzàn.\\
Je publie aussi des photos sur WeChat, et mes amis me likent souvent.

网络很方便，但是花太多时间会影响学习。\\
Wǎngluò hěn fāngbiàn, dànshì huā tài duō shíjiān huì yǐngxiǎng xuéxí.\\
Internet est pratique, mais y passer trop de temps nuit aux études.

所以我觉得应该少上网，多看书。\\
Suǒyǐ wǒ juéde yīnggāi shǎo shàngwǎng, duō kàn shū.\\
Donc je pense qu'il faut moins aller en ligne et lire davantage.
\end{textebox}

\begin{center}
\begin{tabularx}{\linewidth}{|X|X|X|X|}
\hline
\rowcolor{popblueL} \textbf{Caractères} & \textbf{Pinyin} & \textbf{Nature} & \textbf{Traduction} \\
\hline
雅温得 & Yǎwēndé & n.p. & Yaoundé (capitale du Cameroun) \\
\hline
中学 & zhōngxué & n. & collège / lycée \\
\hline
每天 & měitiān & adv. de temps & chaque jour \\
\hline
都 & dōu & adv. & tous (marque la totalité) \\
\hline
因为 & yīnwèi & conj. & parce que \\
\hline
跟 & gēn & prép. & avec \\
\hline
聊天 & liáotiān & v. & discuter, bavarder \\
\hline
发 & fā & v. & envoyer, publier \\
\hline
照片 & zhàopiàn & n. & photo \\
\hline
常常 & chángcháng & adv. & souvent \\
\hline
点赞 & diǎnzàn & v. & liker, mettre un « j'aime » \\
\hline
方便 & fāngbiàn & adj. & pratique, commode \\
\hline
但是 & dànshì & conj. & mais, cependant \\
\hline
花 & huā & v. & passer, dépenser (temps/argent) \\
\hline
影响 & yǐngxiǎng & v. & influencer, affecter, nuire à \\
\hline
所以 & suǒyǐ & conj. & donc, par conséquent \\
\hline
觉得 & juéde & v. & penser, trouver (avis subjectif) \\
\hline
应该 & yīnggāi & v. mod. & devoir, falloir \\
\hline
少 & shǎo & adv./adj. & peu, moins (devant un verbe) \\
\hline
多 & duō & adv./adj. & beaucoup, plus (devant un verbe) \\
\hline
\end{tabularx}
\end{center}

\courssub{Analyse de la structure en trois parties}

Le texte de Marie n'est pas une liste aléatoire ; il suit le plan standard de l'\textbf{expression écrite argumentative simple} exigée au Bac. Voici le découpage précis :

\begin{popfigure}
\begin{tikzpicture}[
    node distance=1.2cm and 1.6cm,
    box/.style={rectangle, rounded corners, draw=popdark, thick, fill=popcream, text width=3.4cm, align=center, minimum height=1.8cm, font=\small},
    arrow/.style={-Stealth, thick, popblue},
    labelbox/.style={rectangle, draw=poporange, fill=poporangeL, text width=3.4cm, align=center, font=\small\bfseries, minimum height=0.9cm}
]
% Nodes
\node[box, align=center] (intro){\textbf{1. INTRODUCTION}\\Présentation + outil + fréquence\\\zh{我叫玛丽……我每天都用手机上网}};
\node[box, right=of intro, align=center] (dev){\textbf{2. DÉVELOPPEMENT}\\Activités + cause + réaction\\\zh{因为我喜欢跟朋友聊天……发照片……点赞}};
\node[box, right=of dev, align=center] (conc){\textbf{3. OPINION / CONCLUSION}\\Nuance + conseil\\\zh{网络很方便，但是……所以我觉得……}};

% Labels above
\node[labelbox, above=0.3cm of intro] (lintro) {Qui ? Quoi ? Quand ?};
\node[labelbox, above=0.3cm of dev] (ldev) {Pourquoi ? Comment ? Avec qui ?};
\node[labelbox, above=0.3cm of conc] (lconc) {Bilan ? Avenir ?};

% Arrows
\draw[arrow] (intro.east) -- (dev.west) node[midway, above, font=\tiny, popdark] {因为};
\draw[arrow] (dev.east) -- (conc.west) node[midway, above, font=\tiny, popdark] {但是 / 所以};
\end{tikzpicture}
\legende{Organigramme de la structure type du paragraphe argumentatif (Introduction $\rightarrow$ Développement $\rightarrow$ Conclusion).}
\end{popfigure}

\begin{definition}[Les trois parties obligatoires]
\textbf{1. Introduction (phrases 1--2) :} \textbf{Qui} parle ? \textbf{Quel} outil ? \textbf{Quelle} fréquence ? \\
\emph{Exemple :} \zh{我叫玛丽，是雅温得一所中学的学生。我每天都用手机上网。} (Présentation de l'énonciateur, localisation camerounaise, outil \zh{手机}, fréquence \zh{每天}).

\textbf{2. Développement (phrases 2--3) :} \textbf{Pourquoi} (cause \zh{因为}) ? \textbf{Quelles actions} précises (verbes + objets) ? \textbf{Avec qui} ? \textbf{Résultat social} (\zh{点赞}) ? \\
\emph{Exemple :} \zh{因为我喜欢跟朋友聊天。我也在微信上发照片，朋友们常常给我点赞。} (Cause, action 1, action 2, réaction des autres).

\textbf{3. Opinion / Conclusion (phrases 4--5) :} \textbf{Constat} (avantage \zh{方便}), \textbf{Nuance} (risque \zh{但是}), \textbf{Conclusion logique} (conséquence \zh{所以} + conseil \zh{应该}). \\
\emph{Exemple :} \zh{网络很方便，但是花太多时间会影响学习。所以我觉得应该少上网，多看书。}
\end{definition}

\courssub{Les connecteurs logiques essentiels : 因为...所以..., 但是, 也/还, 对...来说}

Les connecteurs sont le \textbf{ciment} du texte. Sans eux, vos phrases restent isolées. Le programme de 1\iere{} A exige la maîtrise de quatre types de liaisons logiques pour ce thème.

\begin{propriete}[Règle d'or : 因为... 所以... (Cause $\rightarrow$ Conséquence)]
En français, on évite « Parce que..., donc... » (redondance). \textbf{En chinois, la structure \cle{因为...，所以...} est la norme standard} pour lier explicitement une cause à sa conséquence dans une phrase complexe. Les deux mots se répondent l'un à l'autre (on peut aussi n'en garder qu'un seul, mais jamais traduire « parce que... donc... » mot à mot en français).

\emph{Structure :} \textbf{因为} + Cause (Suj + V) \textbf{，} \textbf{所以} + Conséquence (Suj + V).

\begin{center}
\begin{tabularx}{\linewidth}{|X|X|X|}
\hline
\rowcolor{popblueL} \textbf{Structure} & \textbf{Exemple (Caractères + Pinyin)} & \textbf{Traduction} \\
\hline
因为 + Cause，所以 + Conséquence & \zh{因为微信很方便，所以我每天都用。} (Yīnwèi Wēixìn hěn fāngbiàn, suǒyǐ wǒ měitiān dōu yòng.) & Parce que WeChat est pratique, je l'utilise tous les jours. \\
\hline
因为 + Cause，所以 + Conséquence & \zh{因为下雨，所以我不去踢球。} (Yīnwèi xiàyǔ, suǒyǐ wǒ bù qù tī qiú.) & Parce qu'il pleut, je ne vais pas jouer au foot. \\
\hline
Sujet commun en tête & \zh{我因为喜欢中文，所以努力学习。} (Wǒ yīnwèi xǐhuan Zhōngwén, suǒyǐ nǔlì xuéxí.) & Moi, parce que j'aime le chinois, j'étudie dur. \\
\hline
\end{tabularx}
\end{center}

\textbf{Point de grammaire :} la virgule chinoise \zh{，} sépare obligatoirement la proposition de cause (après \zh{因为}) de la proposition de conséquence (avant \zh{所以}). Ne l'oubliez jamais à l'écrit.
\end{propriete}

\begin{propriete}[La concession et l'opposition : 但是 (dànshì)]
\zh{但是} signifie « mais, cependant ». Il introduit une idée qui contraste avec la précédente. Il se place en \textbf{début de la deuxième proposition} (après une virgule ou en début de phrase).

\begin{center}
\begin{tabularx}{\linewidth}{|X|X|X|}
\hline
\rowcolor{poporangeL} \textbf{Structure} & \textbf{Exemple} & \textbf{Traduction} \\
\hline
Phrase 1，但是 + Phrase 2 & \zh{网络很方便，但是浪费时间。} (Wǎngluò hěn fāngbiàn, dànshì làngfèi shíjiān.) & Internet est pratique, mais cela fait perdre du temps. \\
\hline
Suj + V...，但是 + ... & \zh{我想去，但是没时间。} (Wǒ xiǎng qù, dànshì méi shíjiān.) & Je veux y aller, mais je n'ai pas le temps. \\
\hline
\end{tabularx}
\end{center}
\emph{Variante :} \zh{不过} (bùguò) — « toutefois », un peu plus léger, fréquent à l'oral.
\end{propriete}

\begin{propriete}[La concession renforcée : 虽然... 但是... (bien que... mais...)]
Pour exprimer « bien que... », le chinois utilise la paire \cle{虽然...，但是...}. En français on ne dit pas « bien que... mais... », mais en chinois les deux mots se combinent naturellement.

\emph{Structure :} \textbf{虽然} + Concession \textbf{，} \textbf{但是} + Idée principale.

\emph{Exemple :} \zh{虽然抖音很有意思，但是太浪费时间。} (Suīrán Dǒuyīn hěn yǒu yìsi, dànshì tài làngfèi shíjiān.) « Bien que Douyin soit amusant, cela fait trop perdre de temps. »
\end{propriete}

\begin{propriete}[L'addition d'idées : 也 (yě) vs 还 (hái)]
Les deux signifient « aussi / encore », mais leur emploi diffère :

\begin{center}
\begin{tabularx}{\linewidth}{|X|X|X|X|}
\hline
\rowcolor{popgreenL} \textbf{Mot} & \textbf{Règle principale} & \textbf{Exemple} & \textbf{Traduction} \\
\hline
\zh{也} (yě) & \textbf{Même sujet} (ou sujet différent), action ou état identique/supplémentaire. & \zh{我用微信，我也用抖音。} (Wǒ yòng Wēixìn, wǒ yě yòng Dǒuyīn.) & J'utilise WeChat, j'utilise \textbf{aussi} Douyin. \\
\hline
\zh{还} (hái) & \textbf{Ajout d'un élément nouveau} (objet, action) : « en plus ». & \zh{我发照片，还发视频。} (Wǒ fā zhàopiàn, hái fā shìpín.) & J'envoie des photos, et \textbf{en plus} j'envoie des vidéos. \\
\hline
\end{tabularx}
\end{center}

\emph{Astuce :} dans le texte modèle, \zh{我也在微信上发照片} $\rightarrow$ le sujet \zh{我} est identique à la phrase précédente $\rightarrow$ \zh{也}. Les deux adverbes se placent \textbf{après le sujet et avant le verbe}.
\end{propriete}

\begin{propriete}[La personnalisation de l'avis : 对...来说 (duì... lái shuō)]
Structure : \textbf{对} + Personne/Groupe + \textbf{来说} = « Pour..., du point de vue de... ». Placez cet ensemble en \textbf{début de phrase} (comme un thème annoncé).

\begin{center}
\begin{tabularx}{\linewidth}{|X|X|X|}
\hline
\rowcolor{poppurpleL} \textbf{Structure} & \textbf{Exemple} & \textbf{Traduction} \\
\hline
对 + [Personne] + 来说，... & \zh{对我来说，上网很有用。} (Duì wǒ lái shuō, shàngwǎng hěn yǒuyòng.) & \textbf{Pour moi}, aller sur Internet est très utile. \\
\hline
对 + [Groupe] + 来说，... & \zh{对学生来说，手机很重要。} (Duì xuésheng lái shuō, shǒujī hěn zhòngyào.) & \textbf{Pour les élèves}, le portable est très important. \\
\hline
\end{tabularx}
\end{center}

\emph{Structure voisine utile :} \zh{对 + personne/chose + 好/不好} = « bon / mauvais pour... » : \zh{对眼睛不好} (duì yǎnjing bù hǎo) « mauvais pour les yeux ».
\end{propriete}

\courssub{Méthode pas à pas pour rédiger votre paragraphe}

\begin{methode}[Rédiger un texte « Réseaux sociaux » en 4 étapes]
\emph{Suivez cet ordre pour éviter le syndrome de la page blanche et garantir la cohérence.}

\begin{enumerate}
\item \textbf{Planifier en 3 blocs (brouillon) :} notez 1 mot-clé par bloc.
      \begin{itemize}
      \item \textbf{Intro :} Qui ? Outil ? Fréquence ? (\emph{Ex. : moi, téléphone, chaque jour})
      \item \textbf{Développement :} 2 actions précises + 1 cause (\zh{因为}) + 1 réaction (\zh{点赞}).
      \item \textbf{Conclusion :} 1 avantage + 1 risque (\zh{但是}) + 1 conseil (\zh{应该}).
      \end{itemize}
\item \textbf{Rédiger 1 phrase simple par idée :} écrivez des phrases \textbf{courtes} (Suj + Adv + V + Obj). Pas de phrases à rallonge.
      \emph{Ex. :} \zh{我每天上网。} \zh{我喜欢聊天。} \zh{我发照片。}
\item \textbf{Souder avec les connecteurs :} insérez \zh{因为/所以}, \zh{但是}, \zh{也/还}, \zh{对...来说} entre vos phrases brutes.
      \emph{Transformation :} \zh{我每天上网。因为我喜欢聊天。} $\rightarrow$ \zh{我每天都上网，因为我喜欢聊天。}
\item \textbf{Vérifier l'ordre des mots (S-Temps-Lieu-V-O) et la ponctuation :}
      \begin{itemize}
      \item Complément de temps (\zh{每天}, \zh{昨天}) $\rightarrow$ \textbf{avant} le verbe.
      \item Complément de lieu (\zh{在微信上}) $\rightarrow$ \textbf{avant} le verbe.
      \item Ponctuation : \zh{。} (fin de phrase), \zh{，} (séparation des propositions), \zh{、} (énumération de noms).
      \end{itemize}
\end{enumerate}
\end{methode}

\courssub{Point de vigilance majeur : la place du complément de temps}

C'est l'erreur \textbf{numéro 1 des francophones} (transfert de la structure française « Je vais chaque jour » $\rightarrow$ *\zh{我去每天}).

\begin{attention}[Piège : position de « 每天 » et des compléments de temps]
En chinois, la règle absolue est : \textbf{Sujet + Temps + Lieu + Verbe + Objet}.
Le complément de temps (\zh{每天}, \zh{昨天}, \zh{现在}, \zh{常常}) se place \textbf{immédiatement après le sujet (ou en tête de phrase pour l'emphase), mais JAMAIS après le verbe}.

\begin{center}
\begin{tabularx}{\linewidth}{|X|X|X|}
\hline
\rowcolor{poppinkL} \textbf{Correct} & \textbf{Incorrect} & \textbf{Explication} \\
\hline
\zh{我每天上网。} (Wǒ měitiān shàngwǎng.) & *\zh{我上网每天。} & \zh{每天} (fréquence) se place avant \zh{上网}. \\
\hline
\zh{他在微信上聊天。} (Tā zài Wēixìn shàng liáotiān.) & *\zh{他聊天在微信上。} & \zh{在微信上} (lieu) se place avant le verbe. \\
\hline
\zh{昨天我去了图书馆。} (Zuótiān wǒ qù le túshūguǎn.) & *\zh{我去了图书馆昨天。} & \zh{昨天} en tête de phrase ou après le sujet. \\
\hline
\end{tabularx}
\end{center}

\emph{Rappel :} \zh{都} (dōu) suit souvent le complément de temps pour insister sur la totalité : \zh{我每天都上网。} « Je vais en ligne tous les jours sans exception. »
\end{attention}

\courssub{Exercice résolu : transformation et connecteurs}

\begin{exoresolu}[Du français au chinois : structurer ses idées]
\textbf{Consigne :} traduisez le paragraphe suivant en chinois en respectant la structure en 3 parties et en utilisant \zh{因为...所以...}, \zh{但是}, \zh{还}.

\textit{« Je m'appelle Paul. Je suis lycéen à Douala. J'utilise souvent mon téléphone pour aller sur Internet parce que je veux regarder des vidéos. En plus, je like les photos de mes amis. Internet est amusant, mais c'est mauvais pour les yeux. Donc je pense qu'il faut moins regarder l'écran. »}

\tcblower
\textbf{Solution détaillée et commentée :}

\textbf{Étape 1 : découpage en 3 parties.}
\begin{itemize}
\item \textit{Intro :} Paul, Douala, lycée, téléphone, souvent.
\item \textit{Développement :} cause (regarder des vidéos), action 1 (regarder), action 2 (liker les photos).
\item \textit{Conclusion :} amusant / mauvais pour les yeux (mais) $\rightarrow$ moins regarder (donc).
\end{itemize}

\textbf{Étape 2 : vocabulaire clé.}
\zh{保罗} (Bǎoluó, Paul), \zh{杜阿拉} (Dù'ālā, Douala), \zh{高中} (gāozhōng, lycée), \zh{看视频} (kàn shìpín, regarder des vidéos), \zh{点赞} (diǎnzàn, liker), \zh{有意思} (yǒu yìsi, amusant), \zh{眼睛} (yǎnjing, yeux), \zh{屏幕} (píngmù, écran), \zh{少看} (shǎo kàn, moins regarder).

\textbf{Étape 3 : assemblage avec connecteurs.}

1. \zh{我叫保罗，是杜阿拉一所高中的学生。} (Intro : identité + lieu + statut.)
2. \zh{我经常用手机上网，\textbf{因为}我想看视频。} (Cause : \zh{因为} + sujet \zh{我} + verbe \zh{想看}.)
3. \zh{我\textbf{还}给朋友的照片点赞。} (Addition d'une action nouvelle $\rightarrow$ \zh{还} ; structure \zh{给 + personne + 的 + chose + 点赞} « liker les photos de quelqu'un ».)
4. \zh{网络很有意思，\textbf{但是}对眼睛不好。} (Concession : \zh{有意思} contre \zh{不好} ; structure \zh{对...不好} « mauvais pour... ».)
5. \zh{\textbf{所以}我觉得应该少看屏幕。} (Conséquence : \zh{所以} + \zh{觉得} + \zh{应该} + \zh{少看}.)

\textbf{Texte final assemblé :}
\zh{我叫保罗，是杜阿拉一所高中的学生。我经常用手机上网，因为我想看视频。我还给朋友的照片点赞。网络很有意思，但是对眼睛不好。所以我觉得应该少看屏幕。}

\textbf{Pinyin :}
Wǒ jiào Bǎoluó, shì Dù'ālā yī suǒ gāozhōng de xuésheng. Wǒ jīngcháng yòng shǒujī shàngwǎng, yīnwèi wǒ xiǎng kàn shìpín. Wǒ hái gěi péngyou de zhàopiàn diǎnzàn. Wǎngluò hěn yǒu yìsi, dànshì duì yǎnjing bù hǎo. Suǒyǐ wǒ juéde yīnggāi shǎo kàn píngmù.

\textbf{Traduction littérale de contrôle :} « Je m'appelle Paul, je suis élève dans un lycée de Douala. J'utilise souvent mon téléphone pour aller sur Internet, parce que je veux regarder des vidéos. En plus, je like les photos de mes amis. Internet est très amusant, mais c'est mauvais pour les yeux. Donc je pense qu'il faut moins regarder l'écran. »
\end{exoresolu}

\courssub{Entraînement guidé : à vous de jouer !}

\begin{exercice}[Production écrite : mon réseau social préféré]
\textbf{Consigne :} rédigez un paragraphe de 5 à 6 phrases en chinois (caractères + pinyin) sur \textbf{votre} usage d'UNE application (WeChat, Douyin/TikTok, WhatsApp, QQ...). Vous êtes élève à Yaoundé, Douala, Bafoussam, Ngaoundéré ou Buea.

\textbf{Contraintes obligatoires (grille d'évaluation) :}
\begin{enumerate}
\item \textbf{Intro (1 phrase) :} nom, ville, école, application, fréquence (\zh{每天 / 经常 / 有时}).
\item \textbf{Développement (2 phrases) :} 1 phrase avec \zh{因为...} (raison), 1 phrase avec \zh{也} ou \zh{还} (action supplémentaire : \zh{发视频}, \zh{看直播}, \zh{加群}, \zh{买东西}).
\item \textbf{Conclusion (2 phrases) :} 1 phrase avec \zh{虽然...但是...} ou \zh{...，但是...} (avantage vs inconvénient : \zh{方便 / 有意思} vs \zh{浪费时间 / 影响睡眠 / 近视}). 1 phrase avec \zh{所以} + \zh{应该 / 想} + conseil (\zh{少玩手机}, \zh{多运动}, \zh{早睡}).
\item \textbf{Forme :} caractères soignés, pinyin avec tons, ponctuation chinoise (\zh{。，}).
\end{enumerate}

\textbf{Aide au vocabulaire :}
\zh{直播} (zhíbò, n./v. : direct, live), \zh{加群} (jiā qún, v. : rejoindre un groupe), \zh{近视} (jìnshì, adj./n. : myope, myopie), \zh{睡眠} (shuìmián, n. : sommeil), \zh{运动} (yùndòng, v./n. : faire du sport, sport), \zh{短视频} (duǎn shìpín, n. : vidéo courte), \zh{浪费时间} (làngfèi shíjiān, v. : perdre du temps).
\end{exercice}

\begin{corrige}[Exercice : mon réseau social préféré]
\textbf{Exemple de production attendue (niveau HSK 3 / Bac A) :}

\zh{我叫让，是巴富萨姆一所高中的学生。我每天都用抖音，因为我看短视频很开心。我也喜欢看直播，主播很有趣。抖音虽然很有意思，但是太浪费时间，还影响睡眠。所以我决定以后少玩手机，多去踢球。}

\textbf{Pinyin :}
Wǒ jiào Ràng, shì Bāfùsàmǔ yī suǒ gāozhōng de xuésheng. Wǒ měitiān dōu yòng Dǒuyīn, yīnwèi wǒ kàn duǎn shìpín hěn kāixīn. Wǒ yě xǐhuan kàn zhíbò, zhǔbō hěn yǒuqù. Dǒuyīn suīrán hěn yǒu yìsi, dànshì tài làngfèi shíjiān, hái yǐngxiǎng shuìmián. Suǒyǐ wǒ juédìng yǐhòu shǎo wán shǒujī, duō qù tī qiú.

\textbf{Traduction :}
Je m'appelle Jean, je suis élève dans un lycée de Bafoussam. Chaque jour, j'utilise Douyin, car regarder des vidéos courtes me rend joyeux. J'aime aussi regarder les lives, les animateurs sont très intéressants. Bien que Douyin soit très amusant, cela fait trop perdre de temps et, en plus, cela nuit au sommeil. Donc je décide qu'à l'avenir, je jouerai moins avec mon téléphone et j'irai plus souvent jouer au foot.

\textbf{Analyse des points forts (pour votre auto-évaluation) :}
\begin{itemize}
\item Structure en 3 parties respectée (intro / développement / conclusion).
\item \zh{因为} (phrase 2) bien placé avant le verbe \zh{看}.
\item \zh{也} (phrase 3) pour l'ajout avec le même sujet.
\item \zh{虽然...但是...} (phrase 4) : structure de concession avancée (bonus HSK 4).
\item \zh{太...} (trop...) + \zh{还} (en plus) pour l'accumulation des inconvénients.
\item \zh{所以} + \zh{决定} (décider) + \zh{以后} (à l'avenir) + couple \zh{少/多} (moins/plus).
\item Ordre des mots correct : \zh{我每天都用} (Suj + temps + \zh{都} + V), \zh{影响睡眠} (V-O), \zh{去踢球} (V-O).
\end{itemize}
\end{corrige}

\coursec{Thème et version guidés}

Après avoir analysé la structure d'un texte modèle et identifié les connecteurs logiques pour organiser nos idées (\emph{Modèle de texte commenté et connecteurs}), nous passons à l'étape cruciale de la production écrite : la traduction. Passer du français au chinois (thème) ou du chinois au français (version) ne consiste pas à remplacer mot à mot, mais à \textbf{changer de système de pensée}. Le français place le complément de temps \emph{après} le verbe (SVO + Temps) ; le chinois place le temps \emph{avant} le verbe (S + Temps + V + O). Maîtriser ce basculement, c'est réussir son expression écrite.

\courssub{Méthode du thème : du français vers le chinois}

\begin{methode}[Les 5 étapes du thème réussi]
\begin{enumerate}
    \item \textbf{Analyser la phrase française :} souligner le \cle{sujet}, encadrer l'\cle{expression de temps} (jour, heure, fréquence, durée), identifier le \cle{verbe} et ses \cle{compléments} (COD, COI, lieu, manière).
    \item \textbf{Traduire les unités lexicales :} chercher le mot chinois exact pour chaque mot-clé (attention aux classificateurs, aux verbes composés, aux particules).
    \item \textbf{Réordonner selon la syntaxe chinoise :} appliquer la règle d'or \textbf{Sujet + Temps + Verbe + Compléments}. Placer le lieu avant le verbe (avec \zh{在}), la manière avant le verbe, le COI de personne (avec \zh{给}/\zh{跟}/\zh{对}) avant le verbe.
    \item \textbf{Assembler et vérifier la grammaire :} ajouter les particules (\cle{了} pour l'accompli, \cle{过} pour l'expérientiel, \cle{的}/\cle{得}/\cle{地}), vérifier l'ordre des compléments, choisir le bon classificateur.
    \item \textbf{Relire en chinois :} vérifier chaque caractère (traits, radical), la ponctuation pleine chasse (，。) et la cohérence du sens.
\end{enumerate}
\end{methode}

\begin{popfigure}
\begin{tikzpicture}[
    node distance=1.2cm and 0.9cm,
    box/.style={rectangle, draw=popblue, thick, fill=popblueL, rounded corners, minimum height=1.1cm, text width=2.4cm, align=center, font=\small},
    arrow/.style={-Stealth, thick, popdark},
    labelbox/.style={rectangle, draw=poporange, fill=poporangeL, rounded corners, minimum height=0.8cm, text width=3.2cm, align=center, font=\small\bfseries}
]
% French structure
\node[box, align=center] (fr_s){Sujet\\\emph{(Je)}};
\node[box, right=of fr_s, align=center] (fr_v){Verbe\\\emph{(utilise)}};
\node[box, right=of fr_v, align=center] (fr_cod){COD\\\emph{(mon téléphone)}};
\node[box, right=of fr_cod, align=center] (fr_temp){Temps / But\\\emph{(tous les jours, pour discuter)}};

% Arrow
\node[labelbox, below=of fr_v, yshift=-0.6cm, align=center] (trans){TRANSFORMATION\\Réordonnancement};
\draw[arrow] (fr_v.south) -- (trans.north);

% Chinese structure
\node[box, below=of trans, yshift=-0.6cm, align=center] (zh_s){Sujet\\\zh{我} (wǒ)};
\node[box, right=of zh_s, align=center] (zh_t){Temps / Fréquence\\\zh{每天} (měitiān)};
\node[box, right=of zh_t, align=center] (zh_v){Verbe (+ Particule)\\\zh{用} (yòng) / \zh{花} (huā)};
\node[box, right=of zh_v, align=center] (zh_c){Compléments\\(COI, COD, Durée)};

\draw[arrow] (fr_s) -- (fr_v);
\draw[arrow] (fr_v) -- (fr_cod);
\draw[arrow] (fr_cod) -- (fr_temp);
\draw[arrow] (zh_s) -- (zh_t);
\draw[arrow] (zh_t) -- (zh_v);
\draw[arrow] (zh_v) -- (zh_c);

% Labels
\node[above=0.2cm of fr_v, font=\bfseries\small, text=popdark] {FRANÇAIS : SVO + Temps};
\node[above=0.2cm of zh_t, font=\bfseries\small, text=popdark] {CHINOIS : S + Temps + V + O};
\end{tikzpicture}
\legende{Schéma de réordonnancement : le bloc « Temps » migre de la fin de la phrase (français) vers la position préverbale (chinois).}
\end{popfigure}

\begin{propriete}[Règle d'or de l'ordre des mots : S + Temps + Verbe + Compléments]
En chinois, l'expression de temps (moment, fréquence) se place \textbf{systématiquement après le sujet et avant le verbe}.
\begin{center}
\begin{tabularx}{\linewidth}{|X|X|X|}\hline
\rowcolor{popblueL} \textbf{Français (SVO + Temps)} & \textbf{Chinois (S + Temps + V + O)} & \textbf{Pinyin} \\ \hline
Je vais \textbf{tous les jours} à l'école. & 我\textbf{天天}去学校。 & Wǒ \textbf{tiāntiān} qù xuéxiào. \\ \hline
Il a étudié \textbf{pendant deux heures}. & 他学了\textbf{两个小时}。 & Tā xué le \textbf{liǎng gè xiǎoshí}. \\ \hline
Nous mangeons \textbf{au restaurant} (lieu). & 我们\textbf{在饭馆}吃饭。 & Wǒmen \textbf{zài fànguǎn} chīfàn. \\ \hline
\end{tabularx}
\end{center}
\emph{Remarque importante :} le \textbf{moment} et la \textbf{fréquence} (天天， 每天， 昨天) précèdent le verbe ; la \textbf{durée} (combien de temps a duré l'action) suit le verbe : Verbe + 了 + durée.
\end{propriete}

\courssub{Application : trois thèmes guidés}

\begin{exemplebox}[Thème 1 : Action simple + instrument]
\textbf{Phrase française :} « J'utilise mon téléphone pour discuter avec mes amis. »

\textbf{Étape 1 -- Analyse :}
\begin{itemize}
    \item Sujet : \emph{Je} $\rightarrow$ \zh{我} (wǒ).
    \item Verbe : \emph{utiliser} $\rightarrow$ \zh{用} (yòng).
    \item Instrument : \emph{mon téléphone} $\rightarrow$ \zh{手机} (shǒujī). Structure : \zh{用} + instrument + verbe d'action.
    \item Action principale : \emph{discuter avec mes amis} $\rightarrow$ \zh{跟朋友聊天} (gēn péngyou liáotiān). \zh{跟} (gēn) = avec (introduit la personne avant le verbe).
\end{itemize}

\textbf{Étape 2 -- Réordonnancement (S + V + Instrument + Action) :}
\zh{我} + \zh{用} + \zh{手机} + \zh{跟朋友聊天}。

\textbf{Résultat :}
\begin{center}
\zh{我用手机跟朋友聊天。}\\
\emph{Wǒ yòng shǒujī gēn péngyou liáotiān.}\\
« J'utilise mon téléphone pour discuter avec mes amis. »
\end{center}
\end{exemplebox}

\begin{attention}[Piège : \zh{用} et \zh{跟} se placent avant le verbe d'action]
\zh{用} (yòng) introduit l'instrument \textbf{sans préposition supplémentaire} : \zh{用手机} (avec le téléphone). Ne dites pas \zh{*用着手机} ni \zh{*用在手机}.

\zh{跟} (gēn) « avec » introduit la personne \textbf{avant} le verbe d'action : \zh{跟朋友聊天} (discuter \emph{avec} des amis), et non \zh{*聊天跟朋友}, calqué sur le français.
\end{attention}

\begin{exemplebox}[Thème 2 : Cause / Conséquence + Fréquence]
\textbf{Phrase française :} « Parce que WeChat est pratique, je m'en sers tous les jours. »

\textbf{Étape 1 -- Analyse :}
\begin{itemize}
    \item Proposition 1 (Cause) : \emph{Parce que WeChat est pratique} $\rightarrow$ \zh{因为} (yīnwèi) + sujet \zh{微信} (Wēixìn) + \zh{很方便} (hěn fāngbiàn).
    \item Proposition 2 (Conséquence) : \emph{je m'en sers tous les jours} $\rightarrow$ sujet \zh{我} (wǒ) -- temps \zh{天天} (tiāntiān) -- verbe \zh{用} (yòng).
\end{itemize}

\textbf{Étape 2 -- Structure \zh{因为}... \zh{所以}... (yīnwèi... suǒyǐ...) :}
En chinois standard, \zh{所以} (suǒyǐ) « donc / c'est pourquoi » est \textbf{fortement recommandé} dans la seconde proposition pour marquer la conséquence (alors que le français évite « parce que... donc... »).

\textbf{Résultat :}
\begin{center}
\zh{因为微信很方便，所以我天天用。}\\
\emph{Yīnwèi Wēixìn hěn fāngbiàn, suǒyǐ wǒ tiāntiān yòng.}\\
« Parce que WeChat est très pratique, je l'utilise tous les jours. »
\end{center}
\end{exemplebox}

\begin{aretenir}[Connecteurs de cause et de conséquence]
\begin{itemize}
    \item \zh{因为}... \zh{所以}... : tournure standard, à l'écrit comme à l'oral.
    \item \zh{因为}... \zh{所以}... peut être allégée à l'oral : on garde souvent un seul des deux connecteurs, mais à l'\textbf{examen écrit}, utilisez la paire complète.
    \item \zh{既然}... \zh{就}... (jìrán... jiù...) : « puisque... alors... » (inférence logique, niveau HSK 4).
\end{itemize}
À l'écrit, ne laissez jamais la seconde proposition sans \zh{所以} ou \zh{就} après une proposition en \zh{因为}.
\end{aretenir}

\begin{exemplebox}[Thème 3 : Durée + Jugement de valeur]
\textbf{Phrase française :} « Il passe deux heures par jour sur Internet, c'est perdre du temps. »

\textbf{Étape 1 -- Analyse :}
\begin{itemize}
    \item Sujet : \emph{Il} $\rightarrow$ \zh{他} (tā).
    \item Temps : \emph{par jour / chaque jour} $\rightarrow$ \zh{每天} (měitiān), placé avant le verbe.
    \item Verbe + durée : \emph{passer deux heures} $\rightarrow$ verbe \zh{花} (huā, « dépenser » du temps ou de l'argent) + durée.
    \item Activité : \emph{sur Internet} $\rightarrow$ \zh{上网} (shàngwǎng).
    \item Commentaire : \emph{c'est perdre du temps} $\rightarrow$ \zh{这是浪费时间} (zhè shì làngfèi shíjiān).
\end{itemize}

\textbf{Étape 2 -- Choix de la structure pour « passer du temps » :} voir la règle ci-dessous.

\textbf{Résultat :}
\begin{center}
\zh{他每天花两个小时上网，这是浪费时间。}\\
\emph{Tā měitiān huā liǎng gè xiǎoshí shàngwǎng, zhè shì làngfèi shíjiān.}\\
« Il passe deux heures par jour sur Internet, c'est une perte de temps. »
\end{center}
\end{exemplebox}

\begin{propriete}[Le verbe \zh{花} (huā) : dépenser du temps ou de l'argent]
\textbf{Schéma A :} Sujet + \zh{花} + Durée + Verbe d'activité. \quad \textbf{Schéma B :} Sujet + Verbe + \zh{了} + Durée.
\begin{center}
\begin{tabularx}{\linewidth}{|X|X|}\hline
\rowcolor{popgreenL} \textbf{Structure A : \zh{花} + durée + activité} & \textbf{Structure B : Verbe + \zh{了} + durée} \\ \hline
\zh{他每天花两个小时上网。} & \zh{他每天上网两个小时。} \\
\emph{Tā měitiān huā liǎng gè xiǎoshí shàngwǎng.} & \emph{Tā měitiān shàngwǎng liǎng gè xiǎoshí.} \\ \hline
Accent sur la \textbf{dépense} (le temps coûte). & Description \textbf{factuelle} de la durée de l'action. \\ \hline
\end{tabularx}
\end{center}
Ici, le sens « c'est perdre du temps » appelle la structure A (\zh{花}), qui insiste sur le temps investi (et perdu). \zh{花} s'emploie aussi pour l'argent : \zh{我花了一百块钱买手机壳。} (Wǒ huā le yìbǎi kuài qián mǎi shǒujīké. « J'ai dépensé cent yuans pour acheter une coque de téléphone. »)
\end{propriete}

\courssub{Méthode de la version : du chinois vers le français}

\begin{methode}[Les 4 étapes de la version naturelle]
\begin{enumerate}
    \item \textbf{Segmenter :} repérer les limites de phrases (。 ！ ？) et les groupes syntaxiques (sujet, temps, verbe, COD, particules).
    \item \textbf{Glosser mot à mot (brouillon) :} écrire le sens littéral de chaque mot sous le chinois (\emph{ne pas encore rédiger la phrase française finale}).
    \item \textbf{Reconstruire la syntaxe française :} appliquer l'ordre SVO, conjuguer le verbe (temps, mode), placer les compléments (COD après le verbe, temps souvent en fin de phrase), choisir les prépositions (\emph{à, de, pour, sur}).
    \item \textbf{Polir :} vérifier les accords, le vocabulaire idiomatique (\emph{nuire à} plutôt que \emph{influencer} si le contexte est négatif), la ponctuation française.
\end{enumerate}
\end{methode}

\begin{exemplebox}[Version 1 : Structure \zh{给} (gěi) + Personne + Verbe]
\textbf{Phrase chinoise :} \zh{我的网友给我发了很多照片。}\\
\emph{Wǒ de wǎngyǒu gěi wǒ fā le hěn duō zhàopiàn.}

\textbf{Étape 1 -- Glossaire :}
\begin{center}
\begin{tabularx}{\linewidth}{|X|X|X|X|}\hline
\rowcolor{popblueL} Caractères & Pinyin & Nature & Sens littéral \\ \hline
\zh{我的} & wǒ de & pronom + particule & mon / de moi \\ \hline
\zh{网友} & wǎngyǒu & nom & ami en ligne \\ \hline
\zh{给} & gěi & préposition & à / pour \\ \hline
\zh{我} & wǒ & pronom & moi / me \\ \hline
\zh{发} & fā & verbe & envoyer \\ \hline
\zh{了} & le & particule & accompli \\ \hline
\zh{很多} & hěn duō & quantifieur & beaucoup (de) \\ \hline
\zh{照片} & zhàopiàn & nom & photos \\ \hline
\end{tabularx}
\end{center}

\textbf{Étape 2 -- Réordonnancement français :}
Sujet : \emph{mon ami en ligne}. Verbe : \emph{a envoyé} (passé composé, car \zh{了} marque une action réalisée). COI : \emph{m'}. COD : \emph{beaucoup de photos}.

\textbf{Résultat :}
\begin{center}
« Mon ami en ligne m'a envoyé beaucoup de photos. »
\end{center}
\end{exemplebox}

\begin{attention}[Piège majeur : la place de \zh{给} (gěi) + personne]
En français : \emph{Je t'envoie un message} (Verbe + COI + COD). En chinois : \zh{给} + personne se place \textbf{avant le verbe}.
\begin{itemize}
    \item \textbf{Correct :} \zh{我给你发消息。} (Wǒ gěi nǐ fā xiāoxi. « Je t'envoie un message. »)
    \item \textbf{À éviter :} \zh{*我发消息给你。}, calque du français \emph{j'envoie un message à toi} : tournure lourde et non idiomatique dans ce contexte.
\end{itemize}
\emph{Astuce :} visualisez \zh{给} comme une flèche $\rightarrow$ \zh{给我} qui dirige l'action vers le destinataire \emph{avant} que le verbe n'agisse.
\end{attention}

\begin{exemplebox}[Version 2 : Modal \zh{会} (huì) + verbe \zh{影响} (yǐngxiǎng)]
\textbf{Phrase chinoise :} \zh{抖音很有意思，但是会影响学习。}\\
\emph{Dǒuyīn hěn yǒuyìsi, dànshì huì yǐngxiǎng xuéxí.}

\textbf{Étape 1 -- Glossaire :}
\begin{center}
\begin{tabularx}{\linewidth}{|X|X|X|X|}\hline
\rowcolor{poppurpleL} Caractères & Pinyin & Nature & Sens \\ \hline
\zh{抖音} & Dǒuyīn & nom propre & TikTok (nom chinois) \\ \hline
\zh{很} & hěn & adverbe de degré & très \\ \hline
\zh{有意思} & yǒuyìsi & adjectif & intéressant, amusant \\ \hline
\zh{但是} & dànshì & conjonction & mais, cependant \\ \hline
\zh{会} & huì & verbe modal & risque de, peut (probabilité) \\ \hline
\zh{影响} & yǐngxiǎng & verbe & influencer, affecter \\ \hline
\zh{学习} & xuéxí & verbe / nom & étudier / les études \\ \hline
\end{tabularx}
\end{center}

\textbf{Étape 2 -- Reconstruction :}
\zh{抖音很有意思} $\rightarrow$ \emph{TikTok est très amusant.}
\zh{但是会影响学习} $\rightarrow$ \zh{会} exprime ici une \textbf{conséquence probable} : « risque de ». L'opposition \zh{但是} impose une traduction négative de \zh{影响} : \emph{nuire à} plutôt qu'\emph{influencer}.

\textbf{Résultat :}
\begin{center}
« TikTok est amusant, mais cela peut nuire aux études. »
\end{center}
\end{exemplebox}

\begin{propriete}[Nuances de \zh{影响} (yǐngxiǎng) et de \zh{会} (huì)]
\begin{itemize}
    \item \zh{影响} (yǐngxiǎng) : verbe transitif direct. \zh{影响学习} = affecter les études.
    \item \textbf{Traduction contextuelle :}
        \begin{itemize}
            \item Contexte neutre ou positif : \emph{influencer} -- \zh{这件事影响了我的决定。} (Zhè jiàn shì yǐngxiǎng le wǒ de juédìng. « Cette affaire a influencé ma décision. »)
            \item Contexte négatif (avec \zh{但是}, \zh{不好}, \zh{有害}) : \emph{nuire à}, \emph{affecter négativement}.
        \end{itemize}
    \item \zh{会} (huì) + verbe : « il y a de fortes chances que / cela risque de ». Plus affirmatif que \zh{可能} (kěnéng, « peut-être »).
\end{itemize}
\end{propriete}

\courssub{Synthèse comparative : tableau des pièges fréquents}

\begin{center}
\begin{tabularx}{\linewidth}{|X|X|X|X|}\hline
\rowcolor{poporangeL} \textbf{Point de grammaire} & \textbf{Erreur fréquente (francophone)} & \textbf{Forme correcte} & \textbf{Explication} \\ \hline
\zh{给} (gěi) + COI & \zh{*我发消息给你} & \zh{我给你发消息} & \zh{给} + personne est préverbal. \\ \hline
\zh{花} (huā) + durée & \zh{*我上网花了两个小时} & \zh{我花两个小时上网} / \zh{我上网两个小时} & \zh{花} + durée précède l'activité ; ou durée après le verbe sans \zh{花}. \\ \hline
\zh{影响} (yǐngxiǎng) & traduire toujours par « influencer » & adapter : « nuire à » & le contexte (\zh{但是}, \zh{不好}) impose le sens négatif. \\ \hline
Ordre du temps & \zh{*我上网每天} & \zh{我每天上网} & règle absolue : temps avant le verbe. \\ \hline
Cause / conséquence & oublier \zh{所以} & \zh{因为}... \zh{所以}... & paire corrélative attendue à l'écrit. \\ \hline
\end{tabularx}
\end{center}

\courssub{Entraînement intégré : thème et version sur un texte court}

\begin{textebox}[Dialogue WeChat : les devoirs]
\zh{-- 小李：王伟，你作业做完了吗？}\\
\emph{Xiǎo Lǐ : Wáng Wěi, nǐ zuòyè zuò wán le ma ?}\\
« Xiao Li : Wang Wei, as-tu fini tes devoirs ? »\\
\zh{-- 王伟：还没。我刚才花了一个小时刷抖音，现在得赶紧写。}\\
\emph{Wáng Wěi : Hái méi. Wǒ gāngcái huā le yí gè xiǎoshí shuā Dǒuyīn, xiànzài děi gǎnjǐn xiě.}\\
« Wang Wei : Pas encore. Je viens de passer une heure à faire défiler TikTok, maintenant je dois me dépêcher d'écrire. »\\
\zh{-- 小李：刷视频太久影响效率。放下手机，专心写吧！}\\
\emph{Xiǎo Lǐ : Shuā shìpín tài jiǔ yǐngxiǎng xiàolǜ. Fàngxià shǒujī, zhuānxīn xiě ba !}\\
« Xiao Li : Faire défiler des vidéos trop longtemps nuit à l'efficacité. Pose ton téléphone et écris concentré ! »\\
\zh{-- 王伟：好，我听你的。给我发个加油表情包吧！}\\
\emph{Wáng Wěi : Hǎo, wǒ tīng nǐ de. Gěi wǒ fā gè jiāyóu biǎoqíngbāo ba !}\\
« Wang Wei : D'accord, je suis ton conseil. Envoie-moi un sticker "Courage" ! »
\medskip

\textbf{Glossaire clé :} \zh{作业} (zuòyè) n. devoirs ; \zh{做完} (zuò wán) v. finir de faire ; \zh{刚才} (gāngcái) adv. tout à l'heure ; \zh{刷} (shuā) v. faire défiler (l'écran) ; \zh{得} (děi) v. modal, devoir (obligation) ; \zh{赶紧} (gǎnjǐn) adv. vite, sans tarder ; \zh{效率} (xiàolǜ) n. efficacité ; \zh{放下} (fàngxià) v. poser, lâcher ; \zh{专心} (zhuānxīn) adv. avec concentration ; \zh{听你的} (tīng nǐ de) expr. suivre ton conseil ; \zh{表情包} (biǎoqíngbāo) n. sticker, émoji animé ; \zh{加油} (jiāyóu) expr. courage, vas-y !
\end{textebox}

\begin{exoresolu}[Traduction bidirectionnelle du dialogue]
\textbf{Consigne :} traduire les répliques de \zh{王伟} en français, puis retrouver le chinois de la réplique de \zh{小李} à partir du français.

\textbf{1. Version :} \zh{王伟：还没。我刚才花了一个小时刷抖音，现在得赶紧写。}
\begin{itemize}
    \item \zh{还没} (hái méi) : « pas encore » (négation de l'accompli : \zh{没} sans \zh{了} final).
    \item \zh{我刚才花了一个小时刷抖音} : structure \zh{花} + durée + activité, avec \zh{了} d'accompli : « je viens de passer une heure à faire défiler TikTok ».
    \item \zh{现在得赶紧写} : \zh{得} (děi) = devoir (contrainte, urgence) : « maintenant je dois me dépêcher d'écrire ».
\end{itemize}
\textbf{Proposition :} « Pas encore. Je viens de passer une heure à faire défiler TikTok, maintenant je dois me dépêcher d'écrire. »

\textbf{2. Version :} \zh{王伟：好，我听你的。给我发个加油表情包吧！}
\begin{itemize}
    \item \zh{我听你的} : « je t'écoute / je suis ton conseil ».
    \item \zh{给我发个...} : structure \zh{给} + moi + verbe + classificateur \zh{个} + objet : « envoie-moi un... ».
    \item \zh{吧} (ba) : particule finale de suggestion ou de demande adoucie.
\end{itemize}
\textbf{Proposition :} « D'accord, je suis ton conseil. Envoie-moi un sticker "Courage" ! »

\textbf{3. Thème :} « Faire défiler des vidéos trop longtemps nuit à l'efficacité. Pose le téléphone, écris avec concentration ! »
\begin{itemize}
    \item \emph{Faire défiler des vidéos} : \zh{刷视频} (shuā shìpín), groupe verbal employé comme sujet.
    \item \emph{Trop longtemps} : \zh{太久} (tài jiǔ), placé après le groupe verbal : \zh{刷视频太久}.
    \item \emph{Nuit à l'efficacité} : \zh{影响效率} (yǐngxiǎng xiàolǜ), verbe transitif direct.
    \item \emph{Pose le téléphone} : \zh{放下手机} (fàngxià shǒujī), impératif.
    \item \emph{Écris avec concentration} : \zh{专心写} (zhuānxīn xiě), la manière précède le verbe ; \zh{吧} adoucit l'ordre.
\end{itemize}
\textbf{Proposition chinoise :}
\begin{center}
\zh{刷视频太久影响效率。放下手机，专心写吧！}\\
\emph{Shuā shìpín tài jiǔ yǐngxiǎng xiàolǜ. Fàngxià shǒujī, zhuānxīn xiě ba !}
\end{center}
\end{exoresolu}

\courssub{Exercices d'application (devoirs et entraînement à l'examen)}

\begin{exercice}[Thème : français $\rightarrow$ chinois]
Traduire les phrases suivantes en chinois (caractères + pinyin). Respecter l'ordre S + Temps + V + O et les structures \zh{给}, \zh{花}, \zh{因为}... \zh{所以}...
\begin{enumerate}
    \item Je passe trois heures par semaine à apprendre le chinois.
    \item Parce que mon téléphone est cassé, je ne peux pas aller sur WeChat.
    \item Elle m'a envoyé un message vocal hier soir.
    \item Regarder des vidéos courtes (\zh{短视频} duǎn shìpín) nuit au sommeil (\zh{睡眠} shuìmián).
    \item Envoie-moi une photo de Yaoundé, s'il te plaît (\zh{请} qǐng).
\end{enumerate}
\end{exercice}

\begin{corrige}[Exercice de thème]
\begin{enumerate}
    \item \zh{我每周花三个小时学中文。} (Wǒ měi zhōu huā sān gè xiǎoshí xué Zhōngwén.) -- Structure \zh{花} + durée + activité ; fréquence \zh{每周} avant le verbe.
    \item \zh{因为我的手机坏了，所以我上不了微信。} (Yīnwèi wǒ de shǒujī huài le, suǒyǐ wǒ shàng bù liǎo Wēixìn.) -- \zh{坏了} (huài le) : « est cassé » (adjectif + \zh{了} de changement d'état) ; \zh{上不了} (shàng bù liǎo) : « ne peux pas accéder », complément de capacité négatif.
    \item \zh{她昨天晚上给我发了一个语音。} (Tā zuótiān wǎnshang gěi wǒ fā le yí gè yǔyīn.) -- \zh{给我} avant \zh{发} ; \zh{语音} (yǔyīn) : message vocal ; classificateur \zh{个}.
    \item \zh{看短视频影响睡眠。} (Kàn duǎn shìpín yǐngxiǎng shuìmián.) -- Le groupe verbal \zh{看短视频} sert de sujet ; \zh{影响} est transitif direct.
    \item \zh{请给我发一张雅温得的照片。} (Qǐng gěi wǒ fā yì zhāng Yāwēndé de zhàopiàn.) -- \zh{请} en tête de phrase ; \zh{给我发} ; classificateur \zh{张} (zhāng) pour les objets plats comme les photos ; \zh{雅温得} (Yāwēndé) : Yaoundé.
\end{enumerate}
\end{corrige}

\begin{exercice}[Version : chinois $\rightarrow$ français]
Traduire en français naturel et correct.
\begin{enumerate}
    \item 我的同学每天用电脑做作业，很方便。
    \item 别花太多时间在游戏上，会影响成绩。
    \item 老师给我们发了一个视频，请大家看。
    \item 微信和抖音都很有意思，但是上网太久不好。
\end{enumerate}
\end{exercice}

\begin{corrige}[Exercice de version]
\begin{enumerate}
    \item « Mon camarade de classe utilise l'ordinateur tous les jours pour faire ses devoirs, c'est très pratique. » (\zh{同学} : camarade de classe ; \zh{做作业} : faire les devoirs ; \zh{每天} avant le verbe \zh{用}.)
    \item « Ne passe pas trop de temps sur les jeux, cela risque de nuire aux résultats. » (\zh{别} : impératif négatif « ne... pas » ; \zh{在...上} : « sur » ; \zh{成绩} chéngjì : notes, résultats scolaires ; \zh{会影响} : risque d'affecter.)
    \item « Le professeur nous a envoyé une vidéo, il nous demande à tous de la regarder. » (\zh{老师} : professeur ; \zh{给我们发了} : nous a envoyé ; \zh{请大家看} : invite tout le monde à regarder.)
    \item « WeChat et TikTok sont tous les deux très amusants, mais passer trop de temps en ligne n'est pas bon. » (\zh{和}... \zh{都} : « et... tous les deux » ; \zh{上网太久} : passer trop de temps en ligne ; \zh{不好} : pas bon, mauvais.)
\end{enumerate}
\end{corrige}

\begin{savaistu}[Le clavier chinois : pinyin, écriture manuscrite et Wubi]
Au Cameroun comme en Chine, on saisit le chinois sur ordinateur et sur smartphone principalement avec une \textbf{méthode de saisie} (\zh{输入法} shūrùfǎ) fondée sur le \textbf{pinyin} : on tape \texttt{wo}, puis on choisit \zh{我} dans la liste des candidats. Il existe aussi la méthode \textbf{Wubi} (\zh{五笔} wǔbǐ), fondée sur la structure graphique des caractères (touches correspondant à des types de traits) : très rapide pour les experts, mais longue à apprendre. Sur mobile, la \textbf{reconnaissance de l'écriture manuscrite} (\zh{手写} shǒuxiě) est très populaire : on trace le caractère au doigt et le système le reconnaît -- ce qui suppose de connaître l'ordre des traits ! Quel que soit l'outil, la \textbf{mémorisation des caractères} (écriture, radical, structure) reste indispensable pour choisir le bon candidat à l'écran : les homophones sont nombreux (par exemple \texttt{shi} propose \zh{是}, \zh{十}, \zh{事}, \zh{时}...).
\end{savaistu}

\coursec{Entraînement à l'examen et grille d'évaluation}

Après avoir travaillé la traduction thème/version dans la section précédente, nous passons maintenant à la mise en situation d'examen. L'épreuve d'expression écrite du Baccalauréat (série A) demande de produire un texte court (5 à 8 phrases), cohérent, lexicalement riche et \textbf{exempt d'erreurs de caractères}. Cette section vous donne la grille officielle, la méthode de gestion du temps, les pièges à éviter et un entraînement complet corrigé.

\courssub{Sujet type Bac : Consignes et attentes}

\begin{textebox}[Sujet d'expression écrite — Session type]
\textbf{Consigne :} Rédigez un texte de 5 à 8 phrases en chinois sur le thème : \zh{« 我的社交媒体使用情况 »} (Mon utilisation des réseaux sociaux).

\textbf{Temps imparti :} 25 minutes.

\textbf{Contraintes obligatoires :}
\begin{enumerate}
    \item Utiliser au moins \textbf{5 mots de vocabulaire} de la leçon (ex. \zh{微信}、\zh{直播}、\zh{粉丝}、\zh{点赞}、\zh{分享}、\zh{登录}、\zh{账号}).
    \item Intégrer au moins \textbf{2 connecteurs logiques} parmi : \zh{因为}、\zh{所以}、\zh{但是}、\zh{也}.
    \item Respecter la \textbf{ponctuation chinoise} (。 ， 、 ； ： ？ ！).
    \item Écrire des \textbf{caractères corrects} (traits, ordre, radicaux).
\end{enumerate}
\end{textebox}

\begin{methode}[Gestion des 25 minutes à l'examen]
\begin{enumerate}
    \item \textbf{5 min — Planification :} Lisez le sujet. Cochez les mots du vocabulaire que vous allez utiliser (visez 6-7 pour la marge). Choisissez vos 2 connecteurs. Écrivez le squelette en français : Phrase d'ouverture $\rightarrow$ Fréquence $\rightarrow$ Activité 1 $\rightarrow$ Activité 2 $\rightarrow$ Avis/Nuance $\rightarrow$ Conclusion.
    \item \textbf{15 min — Rédaction :} Écrivez directement en chinois (pas de traduction mot à mot depuis le français). Appliquez l'ordre des mots : \cle{Sujet + Temps + Lieu + Verbe + Objet}. Vérifiez chaque caractère \emph{au moment où vous l'écrivez} (le crochet de \zh{电}, le trait vertical de \zh{王} dans \zh{主}, etc.).
    \item \textbf{5 min — Relecture systématique (ordre impératif) :}
    \begin{enumerate}
        \item \textbf{Caractères :} Aucun trait manquant, aucun caractère inventé, radicaux corrects.
        \item \textbf{Place du complément de temps :} Devant le verbe (pas à la fin comme en français).
        \item \textbf{Connecteurs :} Présents, bien placés, ponctués (virgule après la proposition \zh{因为}, \zh{所以} en tête de proposition principale).
        \item \textbf{Ponctuation :} Point final 。, virgules ，, pas de point français . ni de virgule française ,.
    \end{enumerate}
\end{enumerate}
\end{methode}

\courssub{Grille d'évaluation détaillée (sur 20 points)}

Le correcteur utilise une grille à 5 critères de 4 points chacun. Voici le détail des indicateurs pour chaque niveau.

\begin{popfigure}
\legende{Grille d'évaluation officielle — Expression écrite (20 points)}
\begin{tikzpicture}[
    cell/.style={rectangle, draw=popline, minimum height=1.1cm, align=center, text width=3.2cm, font=\small},
    header/.style={cell, fill=popblue, text=white, font=\small\bfseries},
    crit/.style={cell, fill=popblueL, font=\small\bfseries, text width=2.8cm},
    desc/.style={cell, fill=popcream, text width=4.5cm},
    score/.style={cell, fill=popgreenL, font=\bfseries, text width=1.5cm},
    indic/.style={cell, fill=white, text width=4.5cm, font=\small}
]
% Header
\node[header] (h1) at (0,0) {Critère};
\node[header] (h2) at (3.2,0) {Descriptif};
\node[header] (h3) at (7.7,0) {Points};
\node[header] (h4) at (9.2,0) {Indicateurs de réussite (4/4)};

% Row 1
\node[crit, align=center] (c1) at (0,-1.1){1. Exactitude\\des caractères};
\node[desc, align=center] (d1) at (3.2,-1.1){Traits, ordre, radicaux,\\aucun caractère inventé};
\node[score] (s1) at (7.7,-1.1) {4 pts};
\node[indic, align=center] (i1) at (9.2,-1.1){0 faute de trait/radical ;\\\zh{电} crochet final ok, \zh{王} pas \zh{主} ;\\écriture lisible et standard.};

% Row 2
\node[crit, align=center] (c2) at (0,-2.2){2. Syntaxe \&\\ordre des mots};
\node[desc, align=center] (d2) at (3.2,-2.2){Sujet-Temps-Lieu-Verbe-Objet;\\compléments bien placés};
\node[score] (s2) at (7.7,-2.2) {4 pts};
\node[indic, align=center] (i2) at (9.2,-2.2){Temps avant verbe (\zh{我昨天上网}) ;\\\zh{在} lieu avant verbe ;\\pas de calque français (SVO strict).};

% Row 3
\node[crit, align=center] (c3) at (0,-3.3){3. Connecteurs\\logiques};
\node[desc, align=center] (d3) at (3.2,-3.3){Au moins 2 parmi\\\zh{因为/所以/但是/也}};
\node[score] (s3) at (7.7,-3.3) {4 pts};
\node[indic, align=center] (i3) at (9.2,-3.3){2 connecteurs différents utilisés ;\\structure \zh{因为...所以...} correcte ;\\\zh{但是} en début de 2e proposition.};

% Row 4
\node[crit, align=center] (c4) at (0,-4.4){4. Vocabulaire\\du thème};
\node[desc, align=center] (d4) at (3.2,-4.4){Au moins 5 mots\\de la leçon 38};
\node[score] (s4) at (7.7,-4.4) {4 pts};
\node[indic, align=center] (i4) at (9.2,-4.4){5 mots distincts minimum ;\\ex. \zh{微信、直播、粉丝、点赞、分享、登录、账号} ;\\emploi naturel (verbe/nom).};

% Row 5
\node[crit, align=center] (c5) at (0,-5.5){5. Ponctuation \&\\présentation};
\node[desc, align=center] (d5) at (3.2,-5.5){Ponctuation chinoise ;\\5-8 phrases ;\\propreté};
\node[score] (s5) at (7.7,-5.5) {4 pts};
\node[indic, align=center] (i5) at (9.2,-5.5){Points 。 virgules ， ;\\5 à 8 phrases complètes ;\\pas de ratures, écriture soignée.};

% Total line
\node[header, minimum height=0.8cm] (tot) at (0,-6.6) {TOTAL};
\node[header, minimum height=0.8cm] (tot2) at (3.2,-6.6) {};
\node[header, minimum height=0.8cm] (tot3) at (7.7,-6.6) {\textbf{20 pts}};
\node[header, minimum height=0.8cm, align=center] (tot4) at (9.2,-6.6){Seuil réussite : 10/20 ;\\Excellence : 16+/20};
\end{tikzpicture}
\end{popfigure}

\begin{aretenir}[Règle d'or pour la note]
Un texte de \textbf{5 phrases parfaites} (caractères justes, syntaxe HSK3, connecteurs, vocabulaire ciblé, ponctuation chinoise) obtient \textbf{20/20}. Un texte de 8 phrases avec 3 caractères mal écrits (ex. \zh{电} sans crochet, \zh{微} sans le haut \zh{⺌}, \zh{号} trait manquant) perd \textbf{au minimum 4 points au Critère 1} et souvent des points au Critère 5 (présentation). \cle{La qualité prime sur la quantité.}
\end{aretenir}

\courssub{Complément utile à l'examen : La phrase d'ouverture « passe-partout »}

Pour gagner du temps et montrer un niveau syntaxique élevé (structure prépositive \zh{随着...}), mémorisez cette phrase type. Elle s'adapte à tous les sujets « technologie / société / médias ».

\begin{exemplebox}[Phrase d'ouverture valorisante]
\zh{随着科技的发展，网络越来越重要。} \\
Suízhe kējì de fāzhǎn, wǎngluò yuèláiyuè zhòngyào. \\
« Avec le développement de la technologie, Internet devient de plus en plus important. »

\textbf{Analyse structurelle :}
\begin{itemize}
    \item \zh{随着} (suízhe) : préposition « au fur et à mesure que / avec » + Nom/Groupe verbal. Structure figée \zh{随着 + Cause/Contexte + 主语 + 越来越 + Adj}.
    \item \zh{科技的发展} (kējì de fāzhǎn) : « le développement de la technologie » (thème du Module 5).
    \item \zh{越来越重} (yuèláiyuè zhòngyào) : structure de progression « de plus en plus important ».
\end{itemize}

\textbf{Variantes à adapter :}
\begin{itemize}
    \item \zh{随着手机的普及，微信已成为必需品。} (Avec la généralisation des smartphones, WeChat est devenu indispensable.)
    \item \zh{随着网络的发展，直播很受欢迎。} (Avec le développement du réseau, le direct est très populaire.)
\end{itemize}
\end{exemplebox}

\begin{savaistu}[Qualité vs Quantité à l'examen]
À l'examen du Baccalauréat camerounais comme au HSK, les correcteurs sont instruits de \textbf{pénaliser lourdement les caractères « approximatifs »}. Un candidat qui écrit 6 phrases avec \zh{电} (crochet final manquant $\rightarrow$ devient \zh{田} + \zh{乙} incorrect), \zh{微} (écrit \zh{徽} ou haut raté), \zh{号} (traits dans le désordre) verra son Critère 1 noté 0 ou 1/4, ce qui fait chuter la note finale sous la moyenne, même si les idées sont bonnes. \textbf{Un texte court, propre, sans faute de caractère, vaut toujours mieux qu'un long texte illisible.} Entraînez-vous à écrire les 20 caractères clés de ce module les yeux fermés, dans l'ordre des traits.
\end{savaistu}

\courssub{Piège majeur : Le caractère « de mémoire approximative »}

\begin{attention}[Erreur fatale : Caractères « qui ressemblent » mais sont faux]
Les francophones ont tendance à écrire « de tête » sans vérifier les radicaux et le nombre de traits. En examen, \textbf{un trait manquant = caractère faux = 0 point pour ce mot}.
\begin{center}
\begin{tabularx}{\linewidth}{|X|X|X|X|}
\hline
\rowcolor{popblueL} \textbf{Caractère visé} & \textbf{Erreur fréquente (Francophone)} & \textbf{Analyse de l'erreur} & \textbf{Correction (Rappel traits)} \\
\hline
\zh{电} (diàn) & \zh{田} + trait horizontal simple en bas & Manque le \textbf{crochet final} (提 \emph{tí}) qui part du trait horizontal central vers le bas-droite. C'est le trait 5 (le dernier). & 5 traits : 丶 一 丿 一 \textcolor{poporange}{\textbf{乛}} (crochet). Radicale : \zh{田} (champ). \\
\hline
\zh{微} (wēi) & Haut écrit \zh{彳} (chì) ou \zh{徂} & Le haut est \zh{⺌} (petit) + \zh{⺌} (petit) + \zh{彳} (pas). Pas \zh{徽} (huī, insigne) qui a \zh{心} au milieu. & 13 traits. Haut : deux « petit » \zh{⺌⺌} + \zh{彳}. Bas : \zh{心}. \\
\hline
\zh{号} (hào) & Trait vertical central qui traverse tout & Le trait vertical (丨) \textbf{s'arrête} au trait horizontal du milieu. Il ne traverse pas le bas \zh{口}. & 5 traits : 丨 (trait brisé) 一 丨 (s'arrête) 一. Radicale : \zh{口}. \\
\hline
\zh{直} (zhí) & \zh{真} (zhēn) écrit à la place / œil \zh{目} manquant & \zh{直} = \zh{十} + \zh{目} (œil) + \zh{一}. \zh{真} = \zh{十} + \zh{目} + \zh{匕} (cuillère). Ne confondez pas le bas. & 8 traits. Ordre : 十 $\rightarrow$ 目 $\rightarrow$ 一 final. \\
\hline
\zh{播} (bō) & Radicale \zh{手} (扌) oubliée ou mal formée & Verbe d'action $\rightarrow$ radicale main \zh{扌} (3 traits) obligatoire à gauche. & 15 traits. \zh{扌} + \zh{数} (shù). Écrire \zh{扌} d'abord. \\
\hline
\end{tabularx}
\end{center}
\textbf{Conseil :} Pendant les 5 min de relecture, pointez du stylo \textbf{chaque radicale} de vos 5 mots imposés. Si vous n'êtes pas sûr à 100\% d'un caractère, \textbf{remplacez-le par un synonyme connu} (ex. \zh{发视频} $\rightarrow$ \zh{发帖子} / \zh{上传视频}).
\end{attention}

\courssub{Règle de relecture : L'ordre de vérification en 4 étapes}

\begin{propriete}[Check-list de relecture (5 minutes chrono)]
Appliquez cet ordre \textbf{strict} (ne mélangez pas les étapes) :
\begin{enumerate}
    \item \textbf{\textcolor{poporange}{Caractères (Critère 1) :}} Relisez \emph{uniquement} les caractères. Vérifiez les 5 mots imposés + les radicaux (\zh{扌} \zh{氵} \zh{讠} \zh{钅} \zh{⺌} \zh{田} \zh{口}). Cochez chaque trait mentalement.
    \item \textbf{\textcolor{poporange}{Place du Temps / Lieu (Critère 2) :}} Entourez les mots de temps (\zh{每天}、\zh{昨天}、\zh{现在}、\zh{一小时}). Sont-ils \textbf{avant le verbe} ? (Ex. \zh{我每天上网} $\checkmark$ ; \zh{我上网每天} $\times$).
    \item \textbf{\textcolor{popgreen}{Connecteurs (Critère 3) :}} Cherchez \zh{因为}、\zh{所以}、\zh{但是}、\zh{也}. Sont-ils au bon endroit ? Virgule après la proposition \zh{因为} ; \zh{所以} en tête de proposition principale.
    \item \textbf{\textcolor{popblue}{Ponctuation (Critère 5) :}} Remplacez tous les points français . par 。, virgules , par ，. Vérifiez le nombre de phrases (5 à 8). Propreté de la copie.
\end{enumerate}
\end{propriete}

\courssub{Auto-évaluation et évaluation croisée (Pair-work)}

\begin{experience}[Atelier « Correcteur du Bac » — 20 min]
\textbf{Objectif :} S'approprier la grille en corrigeant une copie anonyme.

\textbf{Déroulement :}
\begin{enumerate}
    \item \textbf{Échange :} Donnez votre brouillon final (ou la production de l'exercice résolu ci-dessous) à votre voisin. Pas de nom dessus.
    \item \textbf{Grille en main :} Chaque correcteur a la grille (5 critères $\times$ 4 pts). Il note \textbf{uniquement} avec un stylo vert : $\checkmark$ (acquis), $\sim$ (approximatif), $\times$ (manquant/faux).
    \item \textbf{Focus Caractères :} Le correcteur entoure \textbf{au crayon rouge} tout caractère suspect (trait manquant, radicale bizarre). Il ne devine pas : si doute $\rightarrow$ $\times$ au Critère 1.
    \item \textbf{Retour :} Rend la copie. L'auteur lit les annotations, corrige \textbf{au stylo bleu} sur sa copie définitive.
    \item \textbf{Bilan classe :} Le prof relève 3 erreurs types vues sur les copies (ex. \zh{号} trait traversant, \zh{电} sans crochet, temps en fin de phrase) et les corrige au tableau.
\end{enumerate}
\end{experience}

\courssub{Exercice résolu : Simulation d'examen (25 min)}

\begin{exoresolu}[Sujet : « 我的社交媒体使用情况 »]
\textbf{Énoncé :} Rédigez 6 phrases. Utilisez : \zh{微信}、\zh{直播}、\zh{粉丝}、\zh{点赞}、\zh{分享}、\zh{因为}、\zh{所以}. Incluez la phrase d'ouverture.

\textbf{Plan de rédaction (5 min) :}
1. Ouverture (\zh{随着...}) 2. Fréquence (\zh{每天}) 3. Action principale \zh{微信} 4. Action secondaire \zh{看直播} 5. Interactions \zh{点赞/分享} + \zh{因为/所以} 6. Conclusion modérée \zh{但是}.

\textbf{Production écrite (15 min) :}
\zh{随着科技的发展，网络越来越重要。} \\
\zh{我每天都登录微信。} \\
\zh{因为我喜欢看直播，所以我经常在手机上看足球直播。} \\
\zh{我也关注很多明星，他们有很多粉丝。} \\
\zh{如果视频有意思，我就点赞和分享给朋友。} \\
\zh{但是上网太久对眼睛不好，所以我控制时间。}

\textbf{Pinyin :}
Suízhe kējì de fāzhǎn, wǎngluò yuèláiyuè zhòngyào. \\
Wǒ měitiān dōu dēnglù Wēixìn. \\
Yīnwèi wǒ xǐhuan kàn zhíbò, suǒyǐ wǒ jīngcháng zài shǒujī shàng kàn zúqiú zhíbò. \\
Wǒ yě guānzhù hěnduō míngxīng, tāmen yǒu hěnduō fěnsī. \\
Rúguǒ shìpín yǒu yìsi, wǒ jiù diǎnzàn hé fēnxiǎng gěi péngyou. \\
Dànshì shàngwǎng tài jiǔ duì yǎnjing bù hǎo, suǒyǐ wǒ kòngzhì shíjiān.

\textbf{Traduction :}
Avec le développement de la technologie, Internet devient de plus en plus important. Je me connecte à WeChat tous les jours. Parce que j'aime regarder les directs, je regarde souvent des directs de foot sur mon téléphone. Je suis aussi beaucoup de stars, elles ont beaucoup de fans. Si la vidéo est intéressante, je like et je partage à des amis. Mais surfer trop longtemps n'est pas bon pour les yeux, donc je contrôle mon temps.

\textbf{Analyse par la grille (Auto-évaluation) :}
\begin{itemize}
    \item \textbf{Critère 1 (4/4) :} Caractères vérifiés : \zh{电} (crochet ok), \zh{微} (haut ok), \zh{号} (vertical ne traverse pas), \zh{直} (bas \zh{一} ok), \zh{播} (radicale \zh{扌} ok). Zéro faute.
    \item \textbf{Critère 2 (4/4) :} \zh{每天} devant \zh{登录} ; \zh{在手机上} devant \zh{看} ; \zh{因为...所以...} structure correcte.
    \item \textbf{Critère 3 (4/4) :} Connecteurs : \zh{因为}、\zh{所以} (paire), \zh{也}、\zh{但是}、\zh{所以} (2e). Total 5 connecteurs $>$ 2 requis.
    \item \textbf{Critère 4 (4/4) :} Mots leçon : \zh{微信}、\zh{直播}$\times$2、\zh{粉丝}、\zh{点赞}、\zh{分享}、\zh{登录} = 7 mots distincts $>$ 5 requis.
    \item \textbf{Critère 5 (4/4) :} 6 phrases. Ponctuation 。 ， 。 。 。 。 Présentation claire.
\end{itemize}
\textbf{Note finale : 20/20.}
\tcblower
\textbf{Solution détaillée (Commentaire du prof) :}
Ce texte est un modèle de réponse « excellence ». Il utilise la phrase d'ouverture fournie (gain de temps, structure complexe \zh{随着...越来越...}). Le vocabulaire est précis et varié (verbes \zh{登录}、\zh{看}、\zh{关注}、\zh{点赞}、\zh{分享}、\zh{控制}). La syntaxe respecte scrupuleusement l'ordre chinois (Temps/lieu avant verbe). La concession \zh{但是...所以...} en fin de texte montre une maturité langagière (niveau HSK 3/4) très valorisée. \textbf{Astuce :} La phrase \zh{如果视频有意思，我就...} (Si vidéo intéressante, alors je...) utilise la structure conditionnelle \zh{如果...就...} apprise en Module 3, ce qui prouve la réactivation des acquis.
\end{exoresolu}

\courssub{Exercice d'entraînement individuel}

\begin{exercice}[Simulation chronométrée — 25 minutes]
\textbf{Consigne :} Rédigez votre propre texte (5-8 phrases) sur \zh{« 我的社交媒体使用情况 »}.
\textbf{Mots imposés (choisissez-en 5 minimum) :} \zh{账号}、\zh{密码}、\zh{视频}、\zh{评论}、\zh{关注}、\zh{发布}、\zh{网红}、\zh{流量}.
\textbf{Connecteurs imposés (2 minimum) :} \zh{因为}、\zh{所以}、\zh{但是}、\zh{也}.
\textbf{Contrainte :} Utilisez la phrase d'ouverture \zh{随着科技的发展，网络越来越重要。} ou une variante.

\textbf{Procédure :}
\begin{enumerate}
    \item Mettez un minuteur sur 25 min.
    \item 5 min : Plan + choix vocabulaire (cochez sur la liste).
    \item 15 min : Rédaction sur copie double (brouillon $\rightarrow$ propre).
    \item 5 min : Relecture \textbf{strictement dans l'ordre} : Caractères $\rightarrow$ Temps/Lieu $\rightarrow$ Connecteurs $\rightarrow$ Ponctuation.
    \item Échange avec un pair pour évaluation croisée (grille 5 $\times$ 4).
\end{enumerate}

\textbf{Objectif visé :} Obtenir \textbf{16/20 minimum} (niveau « Bien ») de manière autonome.
\end{exercice}

\begin{corrige}[Exercice simulation chronométrée]
\textit{Correction individuelle :} Pas de correction unique. Utilisez la grille d'évaluation (Figure ci-dessus) et la check-list de relecture (Propriété 6.4).
\textbf{Points de vigilance pour l'auto-correction :}
\begin{itemize}
    \item Avez-vous écrit \zh{账号} (compte) et non \zh{帐号} (variante traditionnelle acceptée mais \zh{账} est le simplifié standard HSK) ? Radicale \zh{贝} (coquillage/argent) à gauche.
    \item \zh{密码} (mot de passe) : \zh{密} (secret) radicale \zh{宀} (toit) + \zh{必} ; \zh{码} (code) radicale \zh{石} (pierre) + \zh{马}.
    \item \zh{发布} (publier) : \zh{发} (envoyer) + \zh{布} (tissu/étaler). Ne pas confondre \zh{布} et \zh{希} (espérer).
    \item Structure \zh{因为...所以...} : La virgule sépare les deux propositions. Standard : \zh{因为下雨，所以我不去。} (Virgule après la proposition \zh{因为}).
\end{itemize}
\textbf{Si note $<$ 14/20 :} Refaite l'exercice 48h plus tard (effet d'espacement) en ciblant le critère le plus faible.
\end{corrige}

\newpage

\coursec{Activité d'intégration}

\coursec{Leçon 38 — Expression écrite : origine du caractère 电 et radicaux ; caractères des réseaux sociaux}

\courssub{Objectifs de l'activité}
\begin{itemize}
    \item Mobiliser le vocabulaire des réseaux sociaux (\zh{微信}, \zh{抖音}, \zh{上网}, \zh{聊天}, \zh{发照片}, \zh{点赞}, \zh{视频}, \zh{关注}) et les connecteurs logiques (\zh{因为…所以}, \zh{虽然…但是}, \zh{而且}, \zh{也}) pour produire un texte personnel cohérent (niveau HSK 2-3).
    \item Écrire correctement les caractères clés du thème : \zh{电} (électricité/numérique), \zh{话} (parole/téléphone), \zh{网} (réseau/Internet), \zh{信} (message/lettre), en respectant l'ordre des traits, la structure radicale (radicaux \zh{田}, \zh{讠}, \zh{纟}, \zh{亻}) et la proportion des composants.
    \item Appliquer la structure syntaxique canonique \textbf{Sujet + Temps + Lieu + Manière + Verbe + Objet + Durée/Complément} et exprimer une opinion nuancée avec \zh{但是} (\emph{dànshì}) / \zh{可是} (\emph{kěshì}).
    \item Développer l'autonomie en expression écrite et l'esprit critique par l'auto-évaluation et l'évaluation par les pairs à l'aide d'une grille critériée alignée sur les attendus du baccalauréat (LV2).
\end{itemize}

\courssub{Situation de communication : Projet « Pont Numérique Yaoundé–Pékin »}
\begin{textebox}[Contexte authentique : Échange scolaire Cameroun–Chine]
Dans le cadre d'un jumelage entre le \textbf{Lycée de la Salle} (Yaoundé) et le \textbf{Lycée n°4 de Pékin}, votre classe participe à l'opération « \zh{云端交友} » (\emph{Yúnduān jiāoyǔ} : Amis dans le nuage). Chaque élève doit créer une « \zh{微信朋友圈简介} » (\emph{Wēixìn péngyǒuquán jiǎnjiè} : Bio pour le cercle d'amis WeChat) qui sera imprimée et affichée sur le « \zh{云端墙} » (\emph{Yúnduān qiáng} : Mur du nuage) du CDI. Les correspondants chinois viendront lire vos profils lors de la visioconférence de fin d'année.

\textbf{Votre rôle :} Vous êtes \zh{保罗} (\emph{Bǎoluó}, Paul) ou \zh{玛丽} (\emph{Mǎlì}, Marie), élève de Première A à Yaoundé. Vous voulez vous présenter simplement, montrer votre culture numérique camerounaise (football, ndolé, musique) et inviter au dialogue interculturel.
\end{textebox}

\courssub{Origine et évolution du caractère \zh{电} (\emph{diàn})}
\begin{definition}[De l'os oracle au caractère simplifié : l'éclair dans le champ]
Le caractère \zh{电} (\emph{diàn}, électricité, éclair) est la forme simplifiée du caractère traditionnel \zh{電}.
\begin{itemize}
    \item \textbf{Écriture sur os oracle (甲骨文) / Bronze (金文) :} Représentation pictographique d'un éclair \zh{⚡} zigzaguant sous un nuage/pluie \zh{雨} au-dessus d'un champ \zh{田}. Sens originel : « la foudre frappe le champ ».
    \item \textbf{Petit sceau (小篆) :} Structure stabilisée : \zh{雨} (pluie/ciel) au-dessus, \zh{田} (champ) au milieu, \zh{乙} (\emph{yǐ}, trait zigzag) en bas pour l'éclair. Clé radicale traditionnelle : \zh{雨} (\emph{yǔ}, pluie).
    \item \textbf{Simplification (années 1950) :} Suppression de la clé \zh{雨} (jugée redondante) et conservation de \zh{田} + \zh{乙}. La clé radicale devient \zh{田} (\emph{tián}, champ). Le trait \zh{乙} figure l'éclair jaillissant du champ.
\end{itemize}
\textbf{Mnémotechnique :} \emph{« Dans le champ (田), un éclair (乙) jaillit : c'est l'électricité. »}
\end{definition}

\begin{attention}[Piège graphique : \zh{电} vs \zh{田} vs \zh{由}]
Ne pas confondre :
\begin{itemize}
    \item \zh{田} (\emph{tián}) : champ, carré parfait, 5 traits.
    \item \zh{由} (\emph{yóu}) : \zh{田} + trait vertical central qui traverse (sort par le haut et le bas), 5 traits.
    \item \zh{电} (\emph{diàn}) : \zh{田} + \zh{乙} (dernier trait : horizontal court + crochet vertical descendant \zh{乙}, ne traverse \textbf{pas} le haut de \zh{田}). Total : 5 traits.
\end{itemize}
Ordre des traits \zh{电} : \zh{横} (1), \zh{竖} (2), \zh{横} (3), \zh{横} (4) [cadre \zh{田}], \zh{横折弯钩/乙} (5).
\end{attention}

\courssub{Radicaux, structure et écriture des 4 caractères obligatoires}
\begin{definition}[Analyse radicale et phono-sémantique (HSK 2-3)]
\begin{center}
\begin{tabularx}{\linewidth}{|X|X|X|X|X|}
\hline
\rowcolor{popblueL} \textbf{Car.} & \textbf{Pinyin} & \textbf{Radical (Clé)} & \textbf{Type / Composants} & \textbf{Origine / Mnémotechnique} \\ \hline
\zh{电} & diàn & \zh{田} (tián) & Simplifié : Sémantique (\zh{田}) + Phonétique/Symbolique (\zh{乙}) & Éclair (\zh{乙}) dans le champ (\zh{田}). Énergie brute. \\ \hline
\zh{话} & huà & \zh{讠} (yán, parole) & Phono-sémantique : \zh{讠} (sens) + \zh{舌} (shé, langue / phonétique) & La parole (\zh{讠}) sort de la langue (\zh{舌}). \\ \hline
\zh{网} & wǎng & \zh{纟} (mì, fil/soie) & Phono-sémantique : \zh{纟} (sens : filet) + \zh{亡} (wáng, phonétique) & Filet en soie (\zh{纟}) qui capture (\zh{亡} son phonétique). \\ \hline
\zh{信} & xìn & \zh{亻} (rén, homme) & Associatif (会意) : \zh{亻} (homme) + \zh{言} (yán, parole) & La parole (\zh{言}) de l'homme (\zh{亻}) debout = confiance, message, lettre. \\ \hline
\end{tabularx}
\end{center}
\end{definition}

\begin{methode}[Ordre des traits clés (Rappel calligraphique)]
\begin{enumerate}
    \item \zh{话} (13 traits) : \zh{讠} (2 traits : point, courbe horizontale) → \zh{舌} (6 traits : \zh{千} + \zh{口} + trait vertical final). \textbf{Règle :} Clé \zh{讠} à gauche, étroite ; composant droit plus large.
    \item \zh{网} (6 traits) : \zh{纟} (3 traits : diagonale, diagonale, courbe) → \zh{亡} (3 traits : horizontale, horizontale, crochet oblique \zh{撇}). \textbf{Règle :} \zh{纟} ne touche pas \zh{亡} idéalement.
    \item \zh{信} (9 traits) : \zh{亻} (2 traits : trait court, trait long) → \zh{言} (7 traits : \zh{辛} sans le trait final + \zh{口}). \textbf{Règle :} \zh{亻} étroit, \zh{言} large. Ne pas écrire \zh{立} + \zh{日}.
\end{enumerate}
\end{methode}

\courssub{Vocabulaire des réseaux sociaux et de la vie numérique (Module 5)}
\begin{aretenir}[Lexique de base — Fréquence élevée, HSK 2-3]
\begin{center}
\begin{tabularx}{\linewidth}{|X|X|X|X|}
\hline
\rowcolor{popblueL} \textbf{Caractères} & \textbf{Pinyin} & \textbf{Nature} & \textbf{Traduction / Contexte} \\ \hline
\zh{微信} & Wēixìn & Nom propre & WeChat (super-appli chinoise : messagerie, paiement, réseau social) \\ \hline
\zh{抖音} & Dǒuyīn & Nom propre & Douyin (version chinoise de TikTok, vidéo courte) \\ \hline
\zh{QQ} & Q Q & Nom propre & QQ (messagerie historique, encore utilisée par étudiants) \\ \hline
\zh{网络} & wǎngluò & Nom & Internet, le réseau, la toile \\ \hline
\zh{上网} & shàngwǎng & Verbe (VO) & Aller sur Internet, se connecter (\zh{上} monter/sur + \zh{网} filet/réseau) \\ \hline
\zh{聊天} & liáotiān & Verbe (VO) & Discuter, bavarder (\zh{聊} causer + \zh{天} ciel/jour) \\ \hline
\zh{发照片} & fā zhàopiàn & V + N & Poster/envoyer des photos (\zh{发} envoyer/émettre) \\ \hline
\zh{发视频} & fā shìpín & V + N & Poster/envoyer des vidéos \\ \hline
\zh{点赞} & diǎnzàn & Verbe (VO) & Liker, mettre un « j'aime » (\zh{点} cliquer + \zh{赞} louer/approuver) \\ \hline
\zh{关注} & guānzhù & Verbe & Suivre (un compte), s'intéresser à, prêter attention \\ \hline
\zh{直播} & zhíbō & Verbe/Nom & Diffusion en direct, live-streaming \\ \hline
\zh{短视频} & duǎn shìpín & Nom & Vidéo courte (format Reels/TikTok/Douyin) \\ \hline
\zh{方便} & fāngbiàn & Adjectif & Pratique, commode, facile \\ \hline
\zh{花时间} & huā shíjiān & V + N & Passer du temps (collocation fixe) \\ \hline
\zh{雅温得} & Yǎwēndé & Nom propre & Yaoundé (capitale politique) \\ \hline
\zh{杜阿拉} & Dù'ālā & Nom propre & Douala (capitale économique) \\ \hline
\zh{恩多莱} & Ēnduōlái & Nom & Ndolé (plat national camerounais, feuilles amères + arachides) \\ \hline
\zh{喀麦隆队} & Kāmàilóng duì & Nom propre & Équipe du Cameroun (Lions Indomptables) \\ \hline
\end{tabularx}
\end{center}
\end{aretenir}

\begin{savaistu}[Socioculturel : Écosystème numérique Chine vs Cameroun]
\begin{itemize}
    \item \textbf{Chine :} Écosystème fermé (\zh{长城防火墙} Grand Pare-feu). Pas de Google, Facebook, WhatsApp, Instagram, YouTube (bloqués). Dominent : \zh{微信} (WeChat, 1,3 Md utilisateurs, paiement \zh{微信支付} omniprésent), \zh{抖音} (Douyin, algorithme ultra-performant), \zh{微博} (Weibo, type Twitter), \zh{B站} (Bilibili, vidéo jeune/ACG), \zh{小红书} (Xiaohongshu, lifestyle/avis).
    \item \textbf{Cameroun :} Écosystème ouvert. Dominent : WhatsApp (communication \#1), Facebook (réseau social \#1), YouTube, Instagram, TikTok (version internationale). Mobile Money (\zh{Mobile Money} / \zh{Orange Money} / \zh{MTN MoMo}) = paiement mobile roi, analogue à \zh{微信支付}/\zh{支付宝}.
    \item \textbf{Pont :} \zh{微信} fonctionne au Cameroun (roaming ou Wi-Fi) et permet le lien direct avec correspondants chinois. C'est l'outil privilégié pour les échanges scolaires.
\end{itemize}
\end{savaistu}

\courssub{Structures grammaticales utiles et connecteurs logiques (HSK 2-3)}
\begin{propriete}[Règles syntaxiques obligatoires pour la rédaction]
\begin{itemize}
    \item \textbf{Ordre canonique des compléments :} \textbf{Sujet + Temps + Lieu + Manière/Instrument + Verbe + Objet + Durée/Complément de résultat/degré.}
    \item \textbf{Durée (Complément de durée) :} Placée \textbf{APRÈS} le verbe (et l'objet si verbe VO court). Schéma : S + Temps + V + (O) + Durée.
    \item \textbf{Concession :} \zh{虽然} (\emph{suīrán})…, \zh{但是/可是} (\emph{dànshì/kěshì})… \zh{虽然} en tête de proposition, \zh{但是} en tête de la suivante. Pas de \zh{但是} si \zh{虽然} absent.
    \item \textbf{Cause-Effet :} \zh{因为} (\emph{yīnwèi})…, \zh{所以} (\emph{suǒyǐ})… Structure corrélative. \zh{因为} en tête de cause, \zh{所以} en tête de conséquence.
    \item \textbf{Coordination :}
    \begin{itemize}
        \item \zh{和} (\emph{hé}) : Noms / Pronoms uniquement. \zh{我和他} (Moi et lui). \textbf{Interdit :} *\zh{我上网和看书}.
        \item \zh{也} (\emph{yě}) / \zh{而且} (\emph{érqiě}) : Verbes / Propositions. \zh{我上网，也看书} / \zh{我上网，而且看书}.
    \end{itemize}
    \item \textbf{Aspect :} \zh{了} (\emph{le}) action terminée/accomplie (souvent passé proche ou changement d'état) ; \zh{过} (\emph{guò}) expérience vécue (au moins une fois dans le passé).
\end{itemize}
\end{propriete}

\begin{exemplebox}[Durée]
    \zh{我每天上网两个小时。} (\emph{Wǒ měitiān shàngwǎng liǎng gè xiǎoshí.}) — Je vais sur Internet 2h par jour. \\
    \zh{我用手机上网两个小时。} (\emph{Wǒ yòng shǒujī shàngwǎng liǎng gè xiǎoshí.}) — Instrument \zh{用手机} avant V.
    \end{exemplebox}


\courssub{Modèle de texte commenté : Profil WeChat de Paul (保罗)}
\begin{exoresolu}[Production écrite modèle — Contraintes toutes respectées]
\textbf{Texte source (Caractères + Pinyin) :}
\begin{textebox}[Mon profil WeChat — 保罗 (\emph{Bǎoluó})]
我叫保罗，住在雅温得。\\
我有电脑和手机，每天用电话上网两个小时。\\
我常用微信和抖音聊天、发照片、点赞。\\
虽然网络很方便，但是我不想花太多时间。\\
因为我喜欢足球，所以我关注体育新闻。\\
我想认识中国朋友，一起练习中文。
\end{textebox}

\textbf{Traduction française littéraire :}
Je m'appelle Paul, j'habite à Yaoundé. J'ai un ordinateur et un téléphone portable ; chaque jour, j'utilise le téléphone pour aller sur Internet pendant deux heures. J'utilise souvent WeChat et Douyin pour discuter, poster des photos, liker. Bien que le réseau soit très pratique, je ne veux pas y passer trop de temps. Parce que j'aime le football, je suis l'actualité sportive. Je veux me faire des amis chinois pour pratiquer le chinois ensemble.

\tcblower
\textbf{Commentaire linguistique détaillé (Corrigé type correcteur) :}

\begin{enumerate}
    \item \textbf{Vérification des contraintes obligatoires (Check-list élève) :}
    \begin{itemize}
        \item \textbf{Caractères cibles :} \zh{电} dans \zh{电脑} (ordi), \zh{电话} (tél) ; \zh{话} dans \zh{电话} ; \zh{网} dans \zh{上网}, \zh{网络} ; \zh{信} dans \zh{微信}. \checkmark
        \item \textbf{Vocabulaire Module 5 (>5) :} \zh{微信, 抖音, 上网, 聊天, 发照片, 点赞, 关注, 方便, 网络, 花时间}. \checkmark
        \item \textbf{Connecteurs (>2) :} \zh{和} (noms), \zh{虽然…但是…} (concession), \zh{因为…所以…} (cause-effet). \checkmark
        \item \textbf{Opinion nuancée :} Phrase 4 (\zh{虽然…但是…}). \checkmark
        \item \textbf{Durée :} Phrase 2 (\zh{每天...两个小时}). \checkmark
        \item \textbf{Plateformes/Équipements :} \zh{电脑, 手机, 微信, 抖音}. \checkmark
        \item \textbf{Actions (3 verbes) :} \zh{聊天, 发照片, 点赞}. \checkmark
    \end{itemize}

    \item \textbf{Analyse syntaxique phrase par phrase :}
    \begin{itemize}
        \item \textbf{P1 :} \zh{我叫保罗} (S + V + O) ; \zh{住在雅温得} (V + Loc). \zh{住在} indivisible. Transcription standard \zh{雅温得}.
        \item \textbf{P2 :} \zh{我有电脑和手机} (Possession + coordination noms par \zh{和}). \zh{每天用电话上网两个小时} : \textbf{Ordre parfait} : Temps (\zh{每天}) + Instrument (\zh{用电话}) + Verbe (\zh{上网}) + Durée (\zh{两个小时}). Classificateur \zh{个} obligatoire pour \zh{小时}.
        \item \textbf{P3 :} \zh{常用} (Adv + V) + Objets coordonnés \zh{微信和抖音} + Série verbale \zh{聊天、发照片、点赞}. Ponctuation \zh{、} (dunhao) pour énumération verbes/phrases courtes.
        \item \textbf{P4 :} \zh{虽然网络很方便，但是我不想花太多时间}. Structure concessive standard. \zh{不想} (volition négative) + VO \zh{花太多时间}. \zh{太...了} omis car \zh{太多} (trop de) suffit.
        \item \textbf{P5 :} \zh{因为我喜欢足球，所以我关注体育新闻}. Cause-Effet. \zh{关注} + Objet \zh{体育新闻} (vocabulaire extension naturelle HSK 3).
        \item \textbf{P6 :} \zh{想认识} (Volition + V) + \zh{中国朋友} (O) + \zh{一起练习中文} (But/Manière : \zh{一起} + VO). Très naturel pour bio d'échange.
    \end{itemize}

    \item \textbf{Points de vigilance (Erreurs fréquentes francophones) :}
    \end{enumerate}
\end{exoresolu}

\begin{attention}[Pièges à éviter absolument]
    \begin{itemize}
        \item \textbf{Ordre Durée :} *\zh{我两个小时上网} (Faux, calque français). \rightarrow \textbf{\zh{我上网两个小时}}.
        \item \textbf{\zh{和} vs \zh{也/而且} :} *\zh{我上网和发照片} (Faux, \zh{和} ne lie pas verbes). \rightarrow \textbf{\zh{我上网，也发照片}} ou \textbf{\zh{我又上网又发照片}}.
        \item \textbf{Écriture \zh{电} :} Trait final \zh{乙} (horizontal + crochet bas) ne touche \textbf{pas} le bord droit de \zh{田}. *\zh{由} a un trait vertical traversant.
        \item \textbf{Écriture \zh{信} :} Radical \zh{亻} (2 traits : court puis long) bien distinct de \zh{言}. Ne pas écrire \zh{立日}.
        \item \textbf{Pinyin/Tons :} \zh{微信} = \emph{Wēixìn} (1, 4) ; \zh{抖音} = \emph{Dǒuyīn} (3, 1) ; \zh{雅温得} = \emph{Yǎwēndé} (3, 1, 2) ; \zh{杜阿拉} = \emph{Dù'ālā} (4, 1, 1). Apostrophe pour \zh{Dù'ālā} (séparation syllabes).
    \end{itemize}
    \end{attention}


\courssub{Thème et version guidés (Entraînement à la traduction)}
\begin{methode}[Méthode en 4 étapes : Thème (Français → Chinois)]
\begin{enumerate}
    \item \textbf{Analyse :} Identifier le sujet, le temps, le lieu, le verbe, l'objet, les compléments (durée, manière). Repérer les connecteurs logiques.
    \item \textbf{Transposition lexicale :} Choisir le vocabulaire chinois exact (attention aux collocations : \zh{上网} pas *\zh{去网}, \zh{发照片} pas *\zh{放照片}).
    \item \textbf{Ordre syntaxique :} Appliquer S + T + L + M + V + O + Dur. Placer \zh{用} (instrument) avant V. Placer Durée après V.
    \item \textbf{Vérification :} Caractères (radicaux, traits), pinyin (tons), ponctuation chinoise, mots-outils (\zh{了/过} si aspect).
\end{enumerate}
\end{methode}

\begin{exercice}[Thème guidé : Phrases à traduire en chinois (Écrire caractères + pinyin)]
\begin{enumerate}
    \item J'habite à Douala et j'ai un téléphone portable.
    \item Chaque semaine, j'utilise l'ordinateur pour aller sur Internet pendant trois heures.
    \item J'utilise souvent WeChat pour discuter avec mes amis et poster des vidéos.
    \item Bien que Douyin soit amusant, je ne veux pas passer trop de temps dessus.
    \item Parce que j'aime la cuisine camerounaise, je poste souvent des photos de ndolé.
\end{enumerate}
\end{exercice}

\begin{corrige}[Exercice Thème guidé]
\begin{enumerate}
    \item \zh{我住在杜阿拉，有一个手机。} (\emph{Wǒ zhù zài Dù'ālā, yǒu yí gè shǒujī.}) — \zh{有} ne prend pas \zh{了}. Classificateur \zh{个} pour \zh{手机}.
    \item \zh{我每周用电脑上网三个小时。} (\emph{Wǒ měi zhōu yòng diànnǎo shàngwǎng sān gè xiǎoshí.}) — Ordre : Temps \zh{每周} + Instrument \zh{用电脑} + V \zh{上网} + Durée.
    \item \zh{我常用微信和朋友聊天、发视频。} (\emph{Wǒ cháng yòng Wēixìn hé péngyou liáotiān, fā shìpín.}) — \zh{和} lie noms (\zh{朋友}). \zh{、} pour verbes.
    \item \zh{虽然抖音很有意思，但是我不想花太多时间。} (\emph{Suīrán Dǒuyīn hěn yǒu yìsi, dànshì wǒ bù xiǎng huā tài duō shíjiān.}) — \zh{有意思} (intéressant/amusant). \zh{上面} (dessus) implicite ou \zh{在上面} possible mais lourd.
    \item \zh{因为我喜欢喀麦隆菜，所以我常发恩多莱的照片。} (\emph{Yīnwèi wǒ xǐhuan Kāmàilóng cài, suǒyǐ wǒ cháng fā Ēnduōlái de zhàopiàn.}) — \zh{喀麦隆菜} (cuisine camerounaise). \zh{的} attributif.
\end{enumerate}
\end{corrige}

\begin{methode}[Méthode en 3 étapes : Version (Chinois → Français)]
\begin{enumerate}
    \item \textbf{Segmentation :} Couper en groupes syntaxiques (Sujet / Prédicat / Compléments). Identifier les mots-outils (\zh{了, 过, 的, 得, 地, 和, 虽然, 因为}).
    \item \textbf{Traduction littérale mot à mot :} Écrire le sens brut sous chaque mot pour voir la structure.
    \item \textbf{Refonte française :} Produire une phrase française correcte, naturelle, respectant les temps (passé composé pour \zh{了} action terminée, passé simple/imparfait selon contexte, « a déjà » pour \zh{过}).
\end{enumerate}
\end{methode}

\begin{exercice}[Version guidée : Traduire en français]
\begin{enumerate}
    \item \zh{玛丽每天用手机上网一个小时。}
    \item \zh{她常用微信和同学聊天，也发照片。}
    \item \zh{虽然网络很方便，但是她不想看太久手机。}
    \item \zh{因为她喜欢跳舞，所以她关注了很多舞蹈视频。}
    \item \zh{她想去中国旅游，认识中国朋友。}
\end{enumerate}
\end{exercice}

\begin{corrige}[Exercice Version guidée]
\begin{enumerate}
    \item Marie va sur Internet avec son téléphone une heure par jour. (Structure : Sujet + Temps + Instrument + Verbe + Durée).
    \item Elle discute souvent avec ses camarades de classe sur WeChat, et poste aussi des photos. (\zh{也} = aussi ; \zh{、} = énumération actions).
    \item Bien qu'Internet soit pratique, elle ne veut pas regarder son téléphone trop longtemps. (\zh{看太久} = regarder trop longtemps ; \zh{手机} objet de \zh{看}).
    \item Parce qu'elle aime la danse, elle a suivi beaucoup de vidéos de danse. (\zh{关注了} = action accomplie / état résultant : elle les suit maintenant ; \zh{了} après verbe + objet quantifié \zh{很多}).
    \item Elle veut voyager en Chine et se faire des amis chinois. (\zh{去中国旅游} but ; \zh{认识} but parallèle).
\end{enumerate}
\end{corrige}

\courssub{Entraînement à l'examen (Type Baccalauréat LV2) et Grille d'évaluation}
\begin{experience}[Épreuve écrite type — Durée 30 min]
\textbf{Sujet :} « \zh{我的数字生活} » (Ma vie numérique). Rédigez un texte de 80 à 100 caractères chinois (environ 8-10 phrases) pour le blog de l'échange. Contraintes :
\begin{enumerate}
    \item Présentez-vous (nom, ville, classe).
    \item Décrivez vos équipements et plateformes préférées (min. 3 mots : \zh{手机, 电脑, 微信, 抖音, QQ, 网络}).
    \item Décrivez vos activités (min. 4 verbes : \zh{上网, 聊天, 发照片, 发视频, 点赞, 关注, 看新闻, 听音乐, 直播}).
    \item Indiquez une fréquence + durée (\zh{每天/每周} + \zh{两个小时}).
    \item Exprimez un avis nuancé avec \zh{虽然…但是…} ou \zh{虽然…可是…}.
    \item Justifiez une habitude avec \zh{因为…所以…}.
    \item Utilisez \zh{和} (noms) et \zh{也/而且} (verbes).
    \item Les 4 caractères \zh{电, 话, 网, 信} doivent apparaître dans des mots corrects.
\end{enumerate}
\end{experience}

\begin{aretenir}[Grille d'évaluation critériée (20 pts) — Auto / Pair / Professeur]
\begin{center}
\begin{tabularx}{\linewidth}{|X|X|X|}
\hline
\rowcolor{popblueL} \textbf{Critère} & \textbf{Indicateurs de réussite} & \textbf{Points} \\ \hline
\textbf{1. Caractères \& Écriture} (5 pts) & 4 caractères cibles (\zh{电,话,网,信}) corrects (traits, radicaux, proportion) ; écriture lisible, soignée ; pas de caractères inventés. & /5 \\ \hline
\textbf{2. Grammaire \& Syntaxe} (5 pts) & Ordre S-T-L-M-V-O-D respecté ; \zh{了/over} utilisés à bon escient ; structures \zh{虽然…但是}, \zh{因为…所以} correctes ; pas de \zh{和} liant verbes. & /5 \\ \hline
\textbf{3. Vocabulaire \& Orthographe} (4 pts) & Min. 8 mots du Module 5 employés correctement ; pinyin exact (tons) si demandé ; pas de confusion \zh{发/放}, \zh{上网/上线}. & /4 \\ \hline
\textbf{4. Contenu \& Cohérence} (3 pts) & Contraintes de situation respectées (identité, culture CM, invitation dialogue) ; connecteurs variés (\zh{和, 也, 而且}) ; logique du discours. & /3 \\ \hline
\textbf{5. Expression orale (Présentation)} (3 pts) & Prononciation claire (tons, \zh{zh/ch/sh/r}, \zh{ü}) ; fluidité (pas de lecture mot à mot) ; regard public ; respect temps (1 min). & /3 \\ \hline
\rowcolor{popblueL} \textbf{TOTAL} & & \textbf{/20} \\ \hline
\end{tabularx}
\end{center}
\end{aretenir}

\courssub{Bilan de la leçon et pistes de progression}
\begin{aretenir}[Compétences validées — Niveau visé HSK 3 / Baccalauréat LV2]
\begin{itemize}
    \item \textbf{Écriture (Caractères) :} Maîtrise de la décomposition radicale (sémantique \zh{讠, 纟, 亻, 田} + phonétique) pour \zh{电, 话, 网, 信}. Calligraphie soignée (ordre traits, proportion) validée par affichage public.
    \item \textbf{Grammaire/Syntaxe :} Ancrage des structures complexes HSK 3 (\zh{虽然…但是}, \zh{因为…所以}, compléments de durée post-verbaux, ordre canonique) dans une production personnelle authentique (pas d'exercice mécanique).
    \item \textbf{Lexique :} Appropriation du vocabulaire numérique chinois (\zh{微信, 抖音, 点赞, 关注, 直播}) vs camerounais (WhatsApp, Facebook, Mobile Money). Conscience sociolinguistique.
    \item \textbf{Traduction :} Automatisation des pièges classiques thème/version (ordre mots, \zh{和} vs \zh{也}, \zh{了/过}, classificateurs).
    \item \textbf{Expression orale :} Passage à l'oral sans notes (prosodie, tons, confiance) préparé par l'écrit.
    \item \textbf{Évaluation formative :} Grille pair-à-pair internalisée → autonomie, regard critique, métalangage grammatical.
\end{itemize}
\end{aretenir}

\begin{methode}[Pour aller plus loin : Devoirs / Approfondissement / Différenciation]
\begin{enumerate}
    \item \textbf{Enrichissement culturel (Niveau +) :} Ajoutez 2 phrases : \zh{我喜欢吃恩多莱，也喜欢看喀麦隆队踢足球。} (\emph{Wǒ xǐhuan chī Ēnduōlái, yě xǐhuan kàn Kāmàilóng duì tī zúqiú.}) — Vocabulaire : \zh{踢足球} (jouer au foot), \zh{看...踢...} (regarder... jouer...).
    \item \textbf{Nuance temporelle :} Remplacez \zh{两个小时} par \zh{一到两个小时} (\emph{yī dào liǎng gè xiǎoshí}, 1 à 2 heures) ou \zh{大概两个小时} (\emph{dàgài liǎng gè xiǎoshí}, environ 2h).
    \item \textbf{Observation caractères traditionnels :} Écrivez et observez les formes traditionnelles : \zh{電} (\emph{diàn}), \zh{話} (\emph{huà}), \zh{網} (\emph{wǎng}), \zh{信} (\emph{xìn} - inchangé). Repérez la clé \zh{雨} dans \zh{電}, \zh{言} dans \zh{話}, \zh{糸} (\emph{mì}, soie) dans \zh{網}. \textbf{Pas d'apprentissage par cœur}, observation comparative seulement.
    \item \textbf{Production orale enregistrée :} Enregistrez votre profil final (audio WeChat/WhatsApp/Vocaroo) et déposez-le sur l'espace classe OnBuch+ pour correction individualisée de la prononciation (tons, \zh{ü}, rétroflexes \zh{zh/ch/sh/r}, nasales \zh{an/ang/en/eng}).
    \item \textbf{Défi « Correspondant » :} Rédigez une question en chinois à poser à votre correspondant pékinois (ex: \zh{你喜欢喀麦隆的足球吗？} / \zh{你们用什么软件听音乐？}).
\end{enumerate}
\end{methode}

\newpage

\coursec{Méthodes et exercices résolus}

\coursec{Leçon 38 — Expression écrite : origine du caractère 电 et radicaux ; caractères des réseaux sociaux}

\begin{center}
\begin{tabularx}{\linewidth}{|X|X|X|X|}
\hline
\rowcolor{popblueL} \textbf{Caractères} & \textbf{Pinyin} & \textbf{Nature} & \textbf{Traduction} \\
\hline
手机 & shǒujī & n. & téléphone portable \\
网络 & wǎngluò & n. & Internet, réseau \\
上网 & shàngwǎng & v. & aller sur Internet, se connecter \\
微信 & Wēixìn & n. & WeChat \\
发消息 & fā xiāoxi & v. & envoyer un message \\
朋友圈 & péngyouquān & n. & fil d'actualité (Moments) \\
照片 & zhàopiàn & n. & photo, photographie \\
视频 & shìpín & n. & vidéo \\
电话 & diànhuà & n. & téléphone (appel) \\
电视 & diànshì & n. & télévision \\
电脑 & diànnǎo & n. & ordinateur \\
电影 & diànyǐng & n. & cinéma, film \\
电 & diàn & n. & électricité, éclair \\
田 & tián & n. & champ (rizière) \\
由 & yóu & prep./v. & par, à partir de ; provenir \\
申 & shēn & v. & déclarer, expliquer \\
信 & xìn & n. & lettre, message,信息 \\
聊天 & liáotiān & v. & discuter, bavarder \\
用 & yòng & v. & utiliser, se servir de \\
因为……所以…… & yīnwèi… suǒyǐ… & conj. & parce que… donc… \\
但是 & dànshì & conj. & mais, cependant \\
觉得 & juéde & v. & penser, trouver (avis) \\
不好 & bù hǎo & adj. & pas bien, mauvais \\
有用 & yǒuyòng & adj. & utile \\
不能 & bù néng & modal & ne pas pouvoir, ne pas devoir \\
班 & bān & m. & classe (classeurs) \\
部 & bù & m. & classificateur : appareils, téléphones \\
条 & tiáo & m. & classificateur : messages, nouvelles, choses longues \\
张 & zhāng & m. & classificateur : photos, papiers, tables \\
个 & gè & m. & classificateur général \\
\hline
\end{tabularx}
\end{center}

\courssub{Méthode 1 — Produire un texte simple sur les réseaux sociaux}

\begin{methode}[Rédiger un court texte sur les réseaux sociaux]
Pour produire un texte simple et correct sur ce thème, suis toujours ces étapes :
\begin{enumerate}
\item \textbf{Comprendre la consigne} : identifier le sujet, la longueur demandée (souvent 5 à 8 phrases), le destinataire (ami, journal du lycée, blog).
\item \textbf{Lister le vocabulaire utile} avant d'écrire : \zh{手机} shǒujī (téléphone portable), \zh{网络} wǎngluò (Internet), \zh{微信} Wēixìn (WeChat), \zh{发消息} fā xiāoxi (envoyer un message), \zh{朋友圈} péngyouquān (fil d'actualité / « moments »), \zh{照片} zhàopiàn (photo), \zh{视频} shìpín (vidéo), \zh{聊天} liáotiān (discuter).
\item \textbf{Construire un plan en trois temps} : introduction (situation : \zh{我有一部手机}), développement (actions : \zh{用手机发消息、看视频} ; avantages / inconvénients avec \zh{因为……所以……} / \zh{但是}), conclusion (avis personnel avec \zh{我觉得} wǒ juéde).
\item \textbf{Utiliser les connecteurs logiques} : \zh{首先} shǒuxiān (d'abord), \zh{然后} ránhòu (ensuite), \zh{但是} dànshì (mais), \zh{因为……所以……} yīnwèi … suǒyǐ … (parce que … donc …), \zh{最后} zuìhòu (enfin), \zh{还有} hái yǒu (de plus).
\item \textbf{Relire et corriger} : vérifier l'ordre \textbf{Sujet + Temps + Lieu + Outil (\zh{用}) + Verbe + Complément}, la place du complément de temps (avant le verbe), les tons du pinyin, les traits des caractères et la ponctuation chinoise pleine chasse (， 。).
\end{enumerate}
\end{methode}

\begin{exoresolu}[Rédaction guidée : « Mon téléphone portable et moi »]
Énoncé : Rédige un texte de 5 à 6 phrases en chinois sur l'emploi de ton téléphone portable et des réseaux sociaux. Utilise au moins deux connecteurs et donne ton avis.
\tcblower
Solution détaillée :
\begin{enumerate}
\item \textbf{Plan} : (a) Possession : \zh{我有一部新手机} ; (b) Usage quotidien : \zh{每天用手机发消息、看视频} ; (c) Avantage (cause) : \zh{因为可以用微信和杜阿拉的朋友聊天} ; (d) Inconvénient (concession) : \zh{但是玩手机太多对学习不好} ; (e) Avis nuancé : \zh{我觉得手机很有用，可是我们不能玩太多}。
\item \textbf{Structures clés} : \zh{用} yòng (outil) + V ; \zh{因为……所以……} (cause-conséquence) ; \zh{对……不好} (mauvais pour…) ; \zh{可是} kěshì (mais, plus formel que \zh{但是} à l'oral).
\end{enumerate}
Texte rédigé :

\zh{我有一部新手机。我每天用手机发消息、看视频。因为我可以用微信和杜阿拉的朋友聊天，所以我很喜欢它。但是，玩手机太多对学习不好。我觉得手机很有用，可是我们不能玩太多。}

\emph{Wǒ yǒu yí bù xīn shǒujī. Wǒ měitiān yòng shǒujī fā xiāoxi, kàn shìpín. Yīnwèi wǒ kěyǐ yòng Wēixìn hé Dù'ālā de péngyou liáotiān, suǒyǐ wǒ hěn xǐhuan tā. Dànshì, wán shǒujī tài duō duì xuéxí bù hǎo. Wǒ juéde shǒujī hěn yǒuyòng, kěshì wǒmen bù néng wán tài duō.}

« J'ai un nouveau téléphone portable. Chaque jour, je m'en sers pour envoyer des messages et regarder des vidéos. Comme je peux discuter avec mes amis de Douala sur WeChat, je l'aime beaucoup. Mais trop jouer avec le téléphone n'est pas bon pour les études. Je pense que le téléphone est très utile, mais nous ne devons pas en abuser. »

\textbf{Points de grammaire vérifiés} : classificateur \zh{部} bù pour téléphone ; complément de temps \zh{每天} placé avant le verbe ; structure \zh{对……不好} ; connecteurs \zh{因为……所以……} et \zh{但是} / \zh{可是} bien employés ; \zh{不能} pour l'interdiction / le conseil.
\end{exoresolu}

\courssub{Méthode 2 — Écrire correctement les caractères du thème : 电 et familles}

\begin{methode}[Maîtriser l'origine, les traits et les radicaux de 电]
\begin{enumerate}
\item \textbf{Origine et évolution} : \zh{电} diàn (électricité) est la forme \textbf{simplifiée} du caractère traditionnel \zh{電} diàn. L'ancien caractère \zh{電} est un pictogramme composé : \zh{雨} yǔ (pluie/ciel) en haut + \zh{田} tián (champ) + \zh{乚} (éclair jaillissant) en bas → « éclair dans le ciel orageux ». La forme simplifiée \zh{电} conserve \zh{田} (champ) traversé par un trait unique en crochet \zh{乚} (éclair) : mnémotechnique \emph{« un champ traversé par un éclair »}.
\item \textbf{Ordre des 5 traits de \zh{电}} (règle : haut → bas, gauche → droite, horizontal avant vertical, cadre avant contenu, fermeture finale) :
      \begin{enumerate}
      \item \zh{丨} (shù) : trait vertical central, haut → bas.
      \item \zh{(trait brisé)} (héng zhé) : horizontal + crochet vers le bas (coin sup. gauche).
      \item \zh{一} (héng) : trait horizontal médian.
      \item \zh{丨} (shù) : trait vertical central bas.
      \item \zh{乚} (shù wān gōu) : crochet vertical arrondi final, partant du centre vers la droite-bas. \textbf{Dernier trait = crochet qui « sort » du carré.}
      \end{enumerate}
\item \textbf{Radicaux (clés) des mots composés} :
      \begin{itemize}
      \item \zh{电话} diànhuà : \zh{电} (électricité) + \zh{话} huà (parole, radical \zh{讠} yán « parole » + \zh{舌} shé « langue » phonétique). Litt. « paroles électriques ».
      \item \zh{电视} diànshì : \zh{电} + \zh{视} shì (regarder, radical \zh{目} mù « œil » + \zh{示} shì phonétique). Litt. « vision électrique ».
      \item \zh{电脑} diànnǎo : \zh{电} + \zh{脑} nǎo (cerveau, radical \zh{月} ròu « chair/corps » + \zh{乃} nǎi phonétique). Litt. « cerveau électrique ».
      \item \zh{电影} diànyǐng : \zh{电} + \zh{影} yǐng (ombre/image, radical \zh{彡} shān « ornements/poils » + \zh{景} jǐng phonétique). Litt. « ombres électriques ».
      \end{itemize}
\item \textbf{Caractères proches (pièges graphiques)} :
      \begin{center}
      \begin{tabularx}{\linewidth}{|X|X|X|}
      \hline
      \rowcolor{poporangeL} \textbf{Caractère} & \textbf{Pinyin / Sens} & \textbf{Différence graphique clé} \\
      \hline
      \zh{电} & diàn, électricité & Trait central qui \textbf{dépasse en bas} avec crochet arrondi \zh{乚}. \\
      \zh{田} & tián, champ & Carré \textbf{fermé}, symétrique, 5 traits, aucun trait ne sort. \\
      \zh{由} & yóu, par / depuis & Haut : \zh{田} ; Bas : \zh{口} kǒu (bouche) → 5+3=8 traits. \\
      \zh{申} & shēn, déclarer / 9e branche terrestre & Haut : \zh{田} ; Bas : \zh{丨} (vertical) + \zh{一} (horizontal) → 5 traits, pas de crochet final. \\
      \hline
      \end{tabularx}
      \end{center}
\item \textbf{Apprentissage par familles lexicales} (méthode des radicaux phonétiques/sémantiques) :
      \begin{itemize}
      \item Famille \zh{网} wǎng (filet, radical \zh{纟} sī « soie ») → \zh{网络} wǎngluò, \zh{上网} shàngwǎng, \zh{网友} wǎngyǒu (internaute/ami en ligne).
      \item Famille \zh{信} xìn (message, radical \zh{亻} rén « homme » + \zh{言} yán « parole ») → \zh{微信} Wēixìn, \zh{信息} xìnxī (information), \zh{信件} xìnjiàn (lettre).
      \item Famille \zh{视} shì (regarder, radical \zh{目} mù) → \zh{电视} diànshì, \zh{视频} shìpín, \zh{视力} shìlì (acuité visuelle).
      \end{itemize}
\end{enumerate}
\end{methode}

\begin{exoresolu}[Analyse de caractères et formation de mots]
Énoncé : (1) Donne le nombre de traits et le nom du dernier trait de \zh{电}. (2) Quel est le radical de \zh{话} ? Donne son sens et deux autres caractères le contenant. (3) Recopie la paire \zh{电 / 田} et explique la différence graphique et sémantique. (4) Forme deux mots avec \zh{电} (autres que ceux vus en classe) et traduis-les.
\tcblower
Solution détaillée :
\begin{enumerate}
\item \zh{电} s'écrit en \textbf{5 traits} ; le dernier trait est le \textbf{crochet vertical arrondi} \zh{乚} (shù wān gōu).
\item Le radical de \zh{话} est \zh{讠} (yán, clé de la \textbf{parole}). Sens : relatif à la parole, au langage. Autres caractères : \zh{说} shuō (dire), \zh{语} yǔ (langue), \zh{谢} xiè (remercier), \zh{读} dú (lire).
\item \zh{电} diàn : le 4e trait (vertical central) se prolonge par un 5e trait \zh{乚} qui \textbf{sort du cadre} vers la droite. \zh{田} tián : le 4e trait (vertical central) s'arrête au bas du carré, \textbf{fermé}. Sens : « électricité/éclair » vs « champ/rizière ».
\item \zh{电梯} diàntī « ascenseur » (litt. « échelle électrique ») ; \zh{电池} diànchí « pile, batterie » (litt. « étang électrique »). Autres acceptables : \zh{电子} diànzǐ (électronique), \zh{电力} diànlì (énergie électrique).
\end{enumerate}
Conclusion : connaître l'origine pictographique (éclair dans le ciel → champ traversé) et les radicaux sémantiques (\zh{讠}, \zh{目}, \zh{月}) permet de mémoriser durablement, de deviner le sens de nouveaux mots et d'éviter les confusions graphiques coûteuses en examen.
\end{exoresolu}

\courssub{Méthode 3 — Thème : traduire du français vers le chinois}

\begin{methode}[Réussir le thème (français → chinois) : ordre des mots et classificateurs]
\begin{enumerate}
\item \textbf{Sujet en tête} : l'ordre de base est \textbf{Sujet + (Temps / Lieu) + Outil (\zh{用}) + Verbe + Complément}.
\item \textbf{Compléments de temps et de lieu AVANT le verbe} : « Je vais sur Internet le soir » → \zh{我晚上上网。} (Sujet + Temps + Verbe). Interdit : *\zh{我上网晚上。}
\item \textbf{Outil / Instrument avec \zh{用} yòng} placé \textbf{avant le verbe} : « Je prends une photo avec le téléphone » → \zh{我用手机拍照。}
\item \textbf{Destinataire avec \zh{给} gěi} placé \textbf{avant le verbe} (structure \zh{给} + personne + V) : « J'envoie un message à mon ami » → \zh{我给朋友发消息。}
\item \textbf{Classificateurs obligatoires} pour dénombrer : \zh{一部手机} (téléphone), \zh{一条消息} (message), \zh{一张照片} (photo), \zh{两个视频} (vidéos), \zh{一个电话} (appel). Jamais de nom nu après un chiffre.
\item \textbf{Aspect en cours avec \zh{在} zài} devant le verbe : « Il est en train d'écrire » → \zh{他在写。}
\item \textbf{Vérification finale} : tons du pinyin, ponctuation chinoise (， 。), traits des caractères.
\end{enumerate}
\end{methode}

\begin{exoresolu}[Thème guidé : phrases types Bac]
Énoncé : Traduis en chinois :
(1) « Ma sœur envoie des photos à ses amis avec son téléphone. »
(2) « Le soir, nous regardons la télévision à la maison. »
(3) « Il est en train d'écrire un message. »
(4) « Hier, j'ai passé un appel vidéo à mon père. » (vidéo = \zh{视频} shìpín, appel = \zh{电话} diànhuà / \zh{视频电话})
\tcblower
Solution détaillée :
\begin{enumerate}
\item Sujet \zh{我姐姐} wǒ jiějie ; Outil \zh{用手机} ; Destinataire \zh{给朋友} ; Verbe + Complément \zh{发照片} (classif. \zh{张} implicite ou \zh{发几张照片}).
      \zh{我姐姐用手机给朋友发照片。} \emph{Wǒ jiějie yòng shǒujī gěi péngyou fā zhàopiàn.}
\item Temps \zh{晚上} wǎnshang ; Sujet \zh{我们} ; Lieu \zh{在家} zài jiā ; Verbe \zh{看} kàn ; Objet \zh{电视} diànshì. Ordre : Temps + Sujet + Lieu + V + O.
      \zh{晚上我们在家看电视。} \emph{Wǎnshang wǒmen zài jiā kàn diànshì.}
\item Sujet \zh{他} ; Aspect \zh{在} zài ; Verbe \zh{写} xiě ; Objet \zh{一条消息} (classif. \zh{条} tiáo).
      \zh{他在写一条消息。} \emph{Tā zài xiě yì tiáo xiāoxi.}
\item Temps \zh{昨天} zuótiān ; Sujet \zh{我} ; Verbe \zh{打} dǎ (passer un coup de fil) ou \zh{做} zuò ; Objet \zh{一个视频电话} yí gè shìpín diànhuà ; Complément \zh{给爸爸} gěi bàba. Aspect accompli \zh{了} le après le verbe.
      \zh{昨天我给爸爸打了一个视频电话。} \emph{Zuótiān wǒ gěi bàba dǎ le yí gè shìpín diànhuà.}
\end{enumerate}
Conclusion : au thème, la place des compléments (temps, lieu, outil \zh{用}, destinataire \zh{给}) \textbf{avant le verbe} est le critère le plus sanctionné ; le choix du classificateur (\zh{部, 条, 张, 个}) est obligatoire.
\end{exoresolu}

\courssub{Méthode 4 — Version : traduire du chinois vers le français}

\begin{methode}[Réussir la version (chinois → français) : aspect et mots-outils]
\begin{enumerate}
\item \textbf{Lecture globale} : lire le texte entier une fois sans traduire pour saisir le contexte (qui ? quand ? quoi ?).
\item \textbf{Découpage en groupes syntaxiques} : isoler [Sujet] [Temps] [Lieu] [Verbe + Aspect] [Complément]. Souligner les mots-outils : \zh{了} (accompli), \zh{过} guò (expérience), \zh{着} zhe (duratif), \zh{的} de (attributif), \zh{在} zài (progressif/locatif).
\item \textbf{Traduction des mots-outils} :
      \begin{itemize}
      \item \zh{了} le (post-verbal) = action accomplie, terminée → \textbf{passé composé} (ou imparfait selon contexte) en français. Ex: \zh{看了} kàn le → « j'ai regardé / je regardai ».
      \item \zh{过} guò = expérience vécue au moins une fois → \textbf{« avoir déjà »} / passé composé avec « déjà ». Ex: \zh{去过} qù guò → « être déjà allé ».
      \item \zh{的} de = lien entre modificateur et nom → relatif « que/qui » ou adjectif. Ex: \zh{我看的书} wǒ kàn de shū → « le livre que je lis / lu ».
      \end{itemize}
\item \textbf{Réordonnancement en français naturel} : remettre le temps au début ou à la fin, le lieu après le verbe, l'outil avec « avec », le destinataire avec « à ».
\item \textbf{Relecture} : accords, temps des verbes, articles, prépositions, style idiomatique (pas de calque).
\end{enumerate}
\end{methode}

\begin{exoresolu}[Version guidée : texte court]
Énoncé : Traduis en français :
\zh{昨天我在网络上看了两个很有意思的视频。我还给朋友发了一条消息。}

\emph{Zuótiān wǒ zài wǎngluò shàng kàn le liǎng gè hěn yǒu yìsi de shìpín. Wǒ hái gěi péngyou fā le yì tiáo xiāoxi.}
\tcblower
Solution détaillée :
\begin{enumerate}
\item \textbf{Phrase 1} : \zh{昨天} (Hier, temps) — \zh{我} (sujet) — \zh{在网络上} (sur Internet, lieu) — \zh{看了} (verbe + \zh{le} accompli) — \zh{两个} (deux, classif. \zh{个}) — \zh{很有意思的} (très intéressants, \zh{的} attributif) — \zh{视频} (vidéos).
      \zh{了} marque l'accompli → \textbf{passé composé}.
\item \textbf{Phrase 2} : \zh{我} (sujet) — \zh{还} hái (aussi / en plus) — \zh{给朋友} (à un ami) — \zh{发了} (envoyé, accompli) — \zh{一条消息} (un message, classif. \zh{条}).
\item \textbf{Restitution} : « Hier, j'ai regardé sur Internet deux vidéos très intéressantes. J'ai aussi envoyé un message à un ami. »
\end{enumerate}
Conclusion : le piège principal est le calque mot à mot (\zh{在网络上} → « sur Internet » placé après le verbe en français) ; il faut restituer un français idiomatique tout en respectant les marques aspectuelles (\zh{了} = accompli).
\end{exoresolu}

\courssub{Méthode 5 — Comprendre et exploiter un texte sur les médias}

\begin{methode}[Répondre aux questions de compréhension écrite]
\begin{enumerate}
\item \textbf{Lire les questions avant le texte} pour cibler la recherche d'information.
\item \textbf{Repérer les mots-clés} de la question dans le texte (noms propres, verbes, chiffres, mots-outils).
\item \textbf{Répondre par une phrase complète en chinois}, en reprenant le sujet et le verbe de la question :
      \begin{itemize}
      \item Question \zh{谁} shéi (qui) → Réponse avec la personne.
      \item Question \zh{什么} shénme (quoi) → Réponse avec l'objet/l'action.
      \item Question \zh{为什么} wèishénme (pourquoi) → Réponse \textbf{obligatoire} commençant par \zh{因为} yīnwèi.
      \item Question \zh{怎么} zěnme (comment) → Réponse avec \zh{用} yòng (outil) ou manière.
      \end{itemize}
\item \textbf{Ne pas recopier tout le texte} : extraire uniquement l'information demandée, adapter les pronoms (\zh{他} au lieu de \zh{小明}).
\item \textbf{Questions d'opinion} (\zh{你觉得……？} nǐ juéde…?) : donner un avis simple justifié : \zh{我觉得……，因为……。}
\end{enumerate}
\end{methode}

\begin{exoresolu}[Compréhension de texte : « Xiao Ming et son téléphone »]
Énoncé : Lis le texte puis réponds en chinois par des phrases complètes.

\begin{textebox}[Texte support]
\zh{小明是雅温得一所中学的学生。他有一部手机，每天都用微信和朋友聊天。他也很喜欢在网络上看足球视频。但是他妈妈觉得玩手机太多不好，所以他晚上九点以后不玩手机。}
\emph{Xiǎomíng shì Yǎwēndé yí suǒ zhōngxué de xuésheng. Tā yǒu yí bù shǒujī, měitiān dōu yòng Wēixìn hé péngyou liáotiān. Tā yě hěn xǐhuan zài wǎngluò shàng kàn zúqiú shìpín. Dànshì tā māma juéde wán shǒujī tài duō bù hǎo, suǒyǐ tā wǎnshang jiǔ diǎn yǐhòu bù wán shǒujī.}
\end{textebox}

Questions :
(1) 小明用什么和朋友聊天？(2) 他喜欢在网上看什么？(3) 为什么他晚上九点以后不玩手机？
\tcblower
Solution détaillée :
\begin{enumerate}
\item Question sur l'outil (\zh{什么}) → structure \zh{用}… Réponse : \zh{他用微信和朋友聊天。} \emph{Tā yòng Wēixìn hé péngyou liáotiān.} « Il discute avec ses amis sur WeChat. »
\item Question sur l'objet de \zh{看} (regarder) → \zh{足球视频}. Réponse : \zh{他喜欢在网上看足球视频。} \emph{Tā xǐhuan zài wǎngshang kàn zúqiú shìpín.} « Il aime regarder des vidéos de football en ligne. »
\item Question \zh{为什么} → réponse avec \zh{因为}. Réponse : \zh{因为他妈妈觉得玩手机太多不好。} \emph{Yīnwèi tā māma juéde wán shǒujī tài duō bù hǎo.} « Parce que sa maman pense que trop jouer avec le téléphone n'est pas bien. »
\end{enumerate}
Conclusion : une réponse complète reprend le sujet (\zh{他}) et le verbe de la question ; la cause s'exprime systématiquement par \zh{因为}… ; \zh{所以} n'est pas répété dans la réponse courte.
\end{exoresolu}

\courssub{Entraînement à l'examen : Expression écrite et Traduction}

\begin{exercice}[Sujet type Bac — Expression écrite]
\textbf{Consigne :} Tu t'appelles Marie (马丽 Mǎlì). Écris un message de 6 à 8 phrases à ton ami chinois Liu Wei (刘伟 Liú Wěi) pour lui parler de ton utilisation des réseaux sociaux. Tu dois inclure :
\begin{itemize}
\item Le type de téléphone que tu as (neuf / ancien) et son usage principal.
\item L'application que tu utilises le plus (WeChat / TikTok / QQ) et avec qui tu parles.
\item Un avantage (ex: nouvelles, cours, amis loin) et un inconvénient (ex: yeux, temps, sommeil).
\item Ton avis personnel avec \zh{我觉得} ou \zh{我认为}.
\item Au moins \textbf{trois connecteurs} parmi : \zh{首先、然后、但是、因为……所以……、最后、还有}.
\end{itemize}
\textbf{Indications :} Soigne l'ordre des mots (Temps/Lieu/Outil avant V), les classificateurs (\zh{部, 条, 个}), la ponctuation chinoise et les tons.
\end{exercice}

\begin{corrige}[Exercice — Expression écrite]
\textbf{Proposition de rédaction (niveau HSK 3 / Première A) :}

\zh{刘伟，你好！首先，我有一部新手机，是生日礼物。我每天用微信给同学发消息、发语音。因为我在喀麦隆，爸爸在北京工作，所以我常用视频电话和他说话，很方便。还有，我喜欢在网上看中文新闻和足球视频。但是，玩手机太久对眼睛不好，也影响睡眠。最后，我觉得手机很有用，可是我们要控制时间，不能玩太多。祝好！马丽}

\emph{Liú Wěi, nǐ hǎo! Shǒuxiān, wǒ yǒu yí bù xīn shǒujī, shì shēngrì lǐwù. Wǒ měitiān yòng Wēixìn gěi tóngxué fā xiāoxi, fā yǔyīn. Yīnwèi wǒ zài Kāmàilóng, bàba zài Běijīng gōngzuò, suǒyǐ wǒ cháng yòng shìpín diànhuà hé tā shuōhuà, hěn fāngbiàn. Hái yǒu, wǒ xǐhuan zài wǎngshang kàn Zhōngwén xīnwén hé zúqiú shìpín. Dànshì, wán shǒujī tài jiǔ duì yǎnjing bù hǎo, yě yǐngxiǎng shuìmián. Zuìhòu, wǒ juéde shǒujī hěn yǒuyòng, kěshì wǒmen yào kòngzhì shíjiān, bù néng wán tài duō. Zhù hǎo! Mǎlì.}

\textbf{Traduction :} « Liu Wei, bonjour ! D'abord, j'ai un nouveau téléphone, c'est un cadeau d'anniversaire. Chaque jour, j'utilise WeChat pour envoyer des messages et des vocaux à mes camarades. Comme je suis au Cameroun et que mon père travaille à Pékin, j'utilise souvent l'appel vidéo pour lui parler, c'est très pratique. De plus, j'aime regarder les infos chinoises et des vidéos de football en ligne. Mais jouer trop longtemps au téléphone est mauvais pour les yeux et affecte le sommeil. Enfin, je trouve le téléphone très utile, mais nous devons contrôler le temps, ne pas trop jouer. Amicalement, Marie. »

\textbf{Grille d'évaluation (sur 10 pts) :}
\begin{center}
\begin{tabularx}{\linewidth}{|X|X|}
\hline
\rowcolor{popgreenL} \textbf{Critère} & \textbf{Points} \\
\hline
Contenu : 5 éléments requis présents (téléphone, app, avantage, inconvénient, avis) & 3 \\
Grammaire : Ordre des mots (Temps/Lieu/Outil avant V), \zh{用}/\zh{给} corrects, \zh{因为……所以……} & 3 \\
Vocabulaire : Richesse (mots du module), classificateurs corrects (\zh{部, 条, 个}) & 2 \\
Écriture : Caractères corrects, traits, ponctuation chinoise, pinyin (si demandé) & 1 \\
Connecteurs : Au moins 3 utilisés correctement & 1 \\
\hline
\textbf{Total} & \textbf{10} \\
\hline
\end{tabularx}
\end{center}
\end{corrige}

\begin{exercice}[Sujet type Bac — Thème / Version]
\textbf{A. Thème (Français → Chinois)} Traduis :
(1) Mon frère aîné utilise l'ordinateur pour faire ses devoirs le soir.
(2) Nous avons regardé un film chinois hier soir à la télévision.
(3) Est-ce que tu as déjà envoyé un message à ton professeur ? (indice : \zh{过} guò)
(4) Cette photo est très belle, je veux l'envoyer à ma mère tout de suite. (\zh{马上} mǎshàng)

\textbf{B. Version (Chinois → Français)} Traduis :
\zh{现在的学生都有智能手机。他们上课不用手机，下课以后常常上网看视频、玩游戏。老师觉得这样不好，因为影响学习。所以学校有规定：上课不能带手机进教室。}
\emph{Xiànzài de xuésheng dōu yǒu zhìnéng shǒujī. Tāmen shàngkè bù yòng shǒujī, xiàkè yǐhòu chángcháng shàngwǎng kàn shìpín, wán yóuxì. Lǎoshī juéde zhèyàng bù hǎo, yīnwèi yǐngxiǎng xuéxí. Suǒyǐ xuéxiào yǒu guīdìng: shàngkè bù néng dài shǒujī jìn jiàoshì.}
\end{exercice}

\begin{corrige}[Exercice — Thème / Version]
\textbf{A. Thème}
\begin{enumerate}
\item \zh{我哥哥晚上用电脑做作业。} \emph{Wǒ gēge wǎnshang yòng diànnǎo zuò zuòyè.} (Sujet + Temps + Outil + V + Objet)
\item \zh{昨天晚上我们在电视上看了一个中国电影。} \emph{Zuótiān wǎnshang wǒmen zài diànshì shàng kàn le yí gè Zhōngguó diànyǐng.} (Temps + Sujet + Lieu + V + \zh{le} + Classif + Adj + N). Note : \zh{看电影} kàn diànyǐng.
\item \zh{你给老师发过消息吗？} \emph{Nǐ gěi lǎoshī fā guò xiāoxi ma?} (Structure \zh{给} + Pers + V + \zh{过} + Objet + \zh{吗}).
\item \zh{这张照片很漂亮，我马上要发给我妈妈。} \emph{Zhè zhāng zhàopiàn hěn piàoliang, wǒ mǎshàng yào fā gěi wǒ māma.} (Démonstratif + Classif \zh{张} + Adj ; \zh{马上} + \zh{要} + V + \zh{给} + Pers).
\end{enumerate}

\textbf{B. Version}
« Les étudiants d'aujourd'hui ont tous des smartphones. En cours, ils n'utilisent pas leur téléphone ; après les cours, ils vont souvent sur Internet regarder des vidéos, jouer à des jeux. Les professeurs trouvent cela pas bien, car cela influence les études. Donc l'école a un règlement : en cours, on ne peut pas apporter de téléphone dans la salle de classe. »

\textbf{Points clés version} :
\begin{itemize}
\item \zh{现在的学生} = les étudiants actuels / d'aujourd'hui.
\item \zh{智能手机} zhìnéng shǒujī = smartphone (téléphones intelligents).
\item \zh{上课} / \zh{下课} = en cours / après les cours (fin de cours).
\item \zh{常常} chángcháng = souvent.
\item \zh{影响} yǐngxiǎng = influencer, affecter (verbe transitif).
\item \zh{规定} guīdìng = règlement, règle.
\item \zh{带...进...} dài ... jìn ... = apporter ... dans ...
\end{itemize}
\end{corrige}

\begin{savaistu}[Cultures numériques : Cameroun — Chine]
\begin{itemize}
\item \textbf{Applications reines} : Au Cameroun, \zh{WhatsApp} et \zh{Facebook} (Meta) dominent ; en Chine, c'est l'écosystème \zh{微信} Wēixìn (WeChat) : messagerie, paiement (\zh{微信支付} Wēixìn zhīfù), réseaux sociaux (\zh{朋友圈}), services publics, taxis, livraison… « Tout dans WeChat ».
\item \textbf{Paiement mobile} : En Chine, le liquide a presque disparu ; on paie même les petits commerçants de rue via QR code (\zh{扫码} sǎomǎ « scanner le code »). Au Cameroun, \zh{Mobile Money} (MTN MoMo, Orange Money) remplit ce rôle via USSD/SMS, sans nécessiter de smartphone haut de gamme.
\item \textbf{Censure / Grande Muraille} : L'Internet chinois (\zh{中国互联网}) est un intranet national filtré (\zh{防火长城} Fánghuǒ Chángchéng « Grande Muraille pare-feu ») : pas d'accès direct à Google, YouTube, Facebook, WhatsApp, Instagram sans VPN (\zh{翻墙} fānqiáng « franchir le mur »). Les équivalents locaux : \zh{百度} Bǎidù (recherche), \zh{优酷} Yòukù / \zh{哔哩哔哩} Bīlībīlī (vidéo), \zh{微博} Wēibó (microblog type Twitter).
\item \textbf{Éducation} : \zh{钉钉} Dīngdīng (Alibaba) et \zh{腾讯会议} Téngxùn Huìyì (Tencent Meeting) sont les standards pour les cours en ligne (équivalents Zoom/Teams), massivement utilisés pendant le Covid.
\item \textbf{Jeux vidéo} : \zh{王者荣耀} Wángzhě Róngyào (Honor of Kings / Arena of Valor) est un phénomène de société en Chine (MOBA mobile), comme \zh{Free Fire} ou \zh{PUBG Mobile} au Cameroun.
\end{itemize}
\end{savaistu}

\begin{attention}[Les 5 erreurs qui coûtent des points au Bac — Récapitulatif]
\begin{enumerate}
\item \textbf{Confondre les caractères proches} : \zh{电} diàn (électricité, crochet final qui sort) ≠ \zh{田} tián (champ, carré fermé) ≠ \zh{由} yóu (bouche en bas) ≠ \zh{申} shēn (vertical + horizontal en bas). \zh{信} xìn (homme + parole) ≠ \zh{言} yán (parole seule, radical).
\item \textbf{Mauvais ordre des mots (calque français)} : Complément de temps / lieu / outil / destinataire \textbf{TOUJOURS AVANT LE VERBE}. \checkmark \zh{我晚上在家用手机给妈妈发消息。} ✗ *\zh{我发消息给妈妈用手机在家晚上。}
\item \textbf{Oublier ou mal choisir le classificateur} : \zh{一部手机} \checkmark vs *\zh{一手机} ✗ ; \zh{一条消息} \checkmark vs *\zh{一个消息} (accepté à l'oral mais \zh{条} est standard pour message) ; \zh{一张照片} \checkmark.
\item \textbf{Mal placer l'outil \zh{用} et le destinataire \zh{给}} : \zh{用} Outil + V ; \zh{给} Pers + V. Jamais après le verbe. \checkmark \zh{我用电脑写作业。} ✗ *\zh{我写作业用电脑。}
\item \textbf{Négliger les tons, l'aspect et la ponctuation} : \zh{发} fā (1er ton, envoyer) ≠ \zh{伐} fá (abattre) ; \zh{了} le (accompli) vs \zh{过} guò (expérience) vs \zh{着} zhe (duratif) ; \zh{在} zài (progressif/locatif, 4e ton) ≠ \zh{再} zài (de nouveau, 4e ton) — homophones, caractères différents ; Ponctuation : \zh{，} \zh{。} \zh{？} \zh{！} (pleine chasse, pas d'espace).
\end{enumerate}
\end{attention}

\newpage

\coursec{Exercices}

\courssub{Je vérifie mes connaissances}

\begin{aretenir}[Rappel : l'origine du caractère 电]
Le caractère \zh{电} (diàn, « électricité ») vient de l'ancien caractère \zh{申}, qui représentait un éclair zigzaguant dans le ciel. Plus tard, on a ajouté la clé de la pluie \zh{雨} pour écrire « éclair » (forme traditionnelle \zh{電}) ; la forme simplifiée moderne est \zh{电}. Aujourd'hui, \zh{电} sert à écrire : \zh{电话} (diànhuà, téléphone), \zh{电视} (diànshì, télévision), \zh{电影} (diànyǐng, cinéma), \zh{电脑} (diànnǎo, ordinateur).
\end{aretenir}

\begin{exercice}[QCM : la bonne traduction]
Entoure la bonne réponse. Une seule réponse correcte par question.

1. « Liker une photo » se dit en chinois :
\begin{itemize}
\item A. \zh{点赞照片} (diǎn zàn zhàopiàn)
\item B. \zh{点照片赞} (diǎn zhàopiàn zàn)
\item C. \zh{照片点赞} (zhàopiàn diǎn zàn)
\end{itemize}

2. « Envoyer un message » se dit :
\begin{itemize}
\item A. \zh{收信息} (shōu xìnxī)
\item B. \zh{发信息} (fā xìnxī)
\item C. \zh{看信息} (kàn xìnxī)
\end{itemize}

3. « Se connecter à Internet » se dit :
\begin{itemize}
\item A. \zh{下网} (xià wǎng)
\item B. \zh{上网} (shàng wǎng)
\item C. \zh{打网} (dǎ wǎng)
\end{itemize}

4. Le caractère \zh{电} (diàn, « électricité ») contient à l'origine l'idée de :
\begin{itemize}
\item A. l'éclair dans le ciel (puis associé à la pluie)
\item B. le soleil
\item C. la montagne
\end{itemize}
\end{exercice}

\begin{exercice}[Vrai ou faux justifié]
Dis si chaque affirmation est vraie ou fausse, puis justifie ta réponse en une phrase en français.

1. Le radical de \zh{话} (huà, « parole ») est \zh{讠}, la clé de la parole.
2. \zh{电} et \zh{申} s'écrivent exactement de la même façon.
3. \zh{网友} (wǎngyǒu) signifie « ami rencontré sur Internet ».
4. Dans la phrase \zh{我每天上网。}, le complément de temps \zh{每天} se place après le verbe.
\end{exercice}

\begin{exercice}[Vocabulaire : complète le tableau]
Recopie et complète le tableau suivant avec la nature du mot (nom, verbe) et sa traduction française.

\begin{center}\begin{tabularx}{\linewidth}{|X|X|X|X|}\hline
\rowcolor{popblueL} Caractères & Pinyin & Nature & Traduction \\ \hline
\zh{电话} & diànhuà & \ldots & \ldots \\ \hline
\zh{网络} & wǎngluò & \ldots & \ldots \\ \hline
\zh{发} & fā & \ldots & \ldots \\ \hline
\zh{短信} & duǎnxìn & \ldots & \ldots \\ \hline
\zh{聊天} & liáotiān & \ldots & \ldots \\ \hline
\end{tabularx}\end{center}
\end{exercice}

\begin{exercice}[Associe le pinyin au caractère]
Relie chaque caractère de la colonne de gauche au bon pinyin et à la bonne traduction (les colonnes sont mélangées).

\begin{center}\begin{tabularx}{\linewidth}{|X|X|X|}\hline
\rowcolor{popblueL} Caractère & Pinyin & Traduction \\ \hline
\zh{电} & xìn & composer, frapper \\ \hline
\zh{网} & diàn & lettre, message \\ \hline
\zh{信} & dǎ & électricité \\ \hline
\zh{打} & wǎng & filet, réseau \\ \hline
\zh{话} & huà & parole, langue \\ \hline
\end{tabularx}\end{center}
\end{exercice}

\courssub{Je m'entraîne}

\begin{attention}[Caractères proches : ne pas confondre]
\zh{话} (huà, parole ; clé de la parole \zh{讠}) et \zh{活} (huó, vivre ; clé de l'eau \zh{氵}) se ressemblent mais n'ont ni le même radical ni le même sens. De même, \zh{电} (diàn) se termine par un crochet recourbé, alors que \zh{申} (shēn, dans \zh{申请}, shēnqǐng, « faire une demande ») se termine par une verticale droite. Observe bien le dernier trait avant d'écrire.
\end{attention}

\begin{exercice}[Caractères proches : \zh{话} ou \zh{活} ? \zh{电} ou \zh{申} ?]
Complète chaque phrase avec le caractère qui convient, puis traduis-la en français.

1. 打电\zh{＿＿}。(dǎ diàn \ldots)
2. 他说中国\zh{＿＿}说得很好。(Tā shuō Zhōngguó \ldots shuō de hěn hǎo.)
3. 他在雅温得生\zh{＿＿}了十年。(Tā zài Yǎwēndé shēng \ldots le shí nián.)
4. 请把你的\zh{＿＿}请表交给老师。(Qǐng bǎ nǐ de \ldots qǐng biǎo jiāo gěi lǎoshī.)
5. 这个\zh{＿＿}影很有意思。(Zhège \ldots yǐng hěn yǒu yìsi.)
\end{exercice}

\begin{exercice}[Texte à trous : mots-outils et vocabulaire]
Complète le texte suivant avec les mots de la liste : \zh{的} — \zh{了} — \zh{因为} — \zh{但是} — \zh{每天} — \zh{信息}。Puis traduis le texte complet en français.

\begin{textebox}[Texte à compléter]
\zh{＿＿} 网络很方便，我 \zh{＿＿} 都用手机给朋友发 \zh{＿＿}。昨天我收到 \zh{＿＿} 一个网友 \zh{＿＿} 消息，\zh{＿＿} 我还没有回复。\\
(\ldots wǎngluò hěn fāngbiàn, wǒ \ldots dōu yòng shǒujī gěi péngyou fā \ldots. Zuótiān wǒ shōudào \ldots yí ge wǎngyǒu \ldots xiāoxi, \ldots wǒ hái méiyǒu huífù.)\\
« \ldots »
\end{textebox}
\end{exercice}

\begin{exercice}[Remets les mots dans l'ordre]
Remets les mots dans le bon ordre pour former une phrase correcte, écris-la en caractères, puis traduis-la.

1. 上网 / 我 / 每天 / 小时 / 两个 (wǒ / shàng wǎng / měi tiān / liǎng ge xiǎoshí)
2. 因为 / 所以 / 很方便 / 我喜欢 / 微信 / 它 (yīnwèi / suǒyǐ / hěn fāngbiàn / wǒ xǐhuan / Wēixìn / tā)
3. 给 / 发 / 我 / 短信 / 妈妈 / 一条 (gěi / fā / wǒ / duǎnxìn / māma / yì tiáo)
4. 网友 / 我的 / 在 / 住 / 杜阿拉 (wǎngyǒu / wǒ de / zài / zhù / Dù'ālā)
\end{exercice}

\begin{exercice}[Écriture des caractères : traits et radicaux]
\begin{enumerate}
\item Écris les caractères \zh{电}, \zh{话}, \zh{网}, \zh{信}, \zh{发} en numérotant chaque trait (1, 2, 3\ldots) dans l'ordre correct d'écriture.
\item Pour chaque caractère, indique son radical (clé) et donne le nom de cette clé en français (par exemple : clé de la parole, clé du filet\ldots).
\item Trouve un autre caractère que tu connais et qui possède le même radical que \zh{话}.
\end{enumerate}
\end{exercice}

\courssub{Je me prépare au Bac}

\begin{methode}[Réussir l'expression écrite sur les réseaux sociaux]
\begin{enumerate}
\item \textbf{Lis le sujet} et souligne les consignes obligatoires (connecteurs, nombre de phrases, mots imposés).
\item \textbf{Prépare un mini-plan} : situation (1 phrase) → avantages (2 phrases avec \zh{因为}\ldots\zh{所以}\ldots) → inconvénient (1 phrase avec \zh{但是}) → opinion personnelle (\zh{我觉得}\ldots).
\item \textbf{Rédige avec des phrases courtes} : Sujet + (complément de temps) + verbe + objet. N'écris que des structures que tu maîtrises.
\item \textbf{Vérifie} : ordre des mots, place de \zh{了}, classificateurs (\zh{一条短信}, \zh{一个网友}), dernier trait des caractères \zh{电}/\zh{申} et \zh{话}/\zh{活}.
\end{enumerate}
\end{methode}

\begin{exercice}[Compréhension de l'écrit et questions de langue]
Lis le texte suivant, puis réponds aux questions.

\begin{textebox}[Les réseaux sociaux au Cameroun]
现在，很多喀麦隆年轻人每天都用社交媒体。他们用手机上网、发信息、看视频，也喜欢给朋友点赞。网络很方便，大家可以很快地知道新闻。但是，有些年轻人每天花太多时间上网，这影响他们的学习。我觉得，我们应该好好利用网络，但是不应该忘记学习和运动。\\
(Xiànzài, hěn duō Kāmàilóng niánqīngrén měi tiān dōu yòng shèjiāo méitǐ. Tāmen yòng shǒujī shàng wǎng, fā xìnxī, kàn shìpín, yě xǐhuan gěi péngyou diǎn zàn. Wǎngluò hěn fāngbiàn, dàjiā kěyǐ hěn kuài de zhīdào xīnwén. Dànshì, yǒuxiē niánqīngrén měi tiān huā tài duō shíjiān shàng wǎng, zhè yǐngxiǎng tāmen de xuéxí. Wǒ juéde, wǒmen yīnggāi hǎohǎo lìyòng wǎngluò, dànshì bù yīnggāi wàngjì xuéxí hé yùndòng.)\\
« Aujourd'hui, beaucoup de jeunes Camerounais utilisent les réseaux sociaux chaque jour. Avec leur téléphone, ils se connectent à Internet, envoient des messages, regardent des vidéos et aiment aussi liker leurs amis. Internet est très pratique : tout le monde peut connaître les nouvelles très vite. Mais certains jeunes passent trop de temps en ligne chaque jour, et cela nuit à leurs études. Je pense que nous devons bien utiliser Internet, mais que nous ne devons pas oublier les études et le sport. »
\end{textebox}

\textbf{Questions de compréhension} (réponds en français) :
\begin{enumerate}
\item Que font les jeunes Camerounais sur les réseaux sociaux ? Cite trois activités.
\item Quel est l'avantage d'Internet mentionné dans le texte ?
\item Quel est le danger signalé par l'auteur ?
\item Quelle est l'opinion personnelle de l'auteur ? Relevez la phrase chinoise qui l'exprime.
\end{enumerate}

\textbf{Questions de langue} :
\begin{enumerate}
\item Relevez dans le texte deux connecteurs et expliquez leur fonction.
\item Quel est le radical du caractère \zh{影} dans \zh{影响} ? Et celui de \zh{信} dans \zh{信息} ?
\item Traduisez en français la phrase : \zh{有些年轻人每天花太多时间上网。}
\end{enumerate}
\end{exercice}

\begin{exercice}[Thème et version]
\textbf{A. Thème} — Traduis en chinois (caractères obligatoires, pinyin en plus si tu veux) :

1. « Parce qu'Internet est très pratique, je discute chaque jour avec mes amis, mais je ne perds pas de temps. »
2. « Ma sœur envoie un message à son amie de Douala. »
3. « J'aime beaucoup liker les photos de mes camarades de classe. »

\textbf{B. Version} — Traduis en français :

1. \zh{我的网友每天花三个小时上网，我觉得这影响学习。}\\
(Wǒ de wǎngyǒu měi tiān huā sān ge xiǎoshí shàng wǎng, wǒ juéde zhè yǐngxiǎng xuéxí.)
2. \zh{昨天我给老师发了一条短信。}\\
(Zuótiān wǒ gěi lǎoshī fā le yì tiáo duǎnxìn.)
3. \zh{因为手机很方便，所以大家都喜欢用手机聊天。}\\
(Yīnwèi shǒujī hěn fāngbiàn, suǒyǐ dàjiā dōu xǐhuan yòng shǒujī liáotiān.)
\end{exercice}

\begin{exercice}[Expression écrite : sujet de type Baccalauréat]
\textbf{Sujet :} \zh{我和社交媒体} (Wǒ hé shèjiāo méitǐ) — « Les réseaux sociaux et moi ».

Rédige un texte en chinois de \textbf{6 phrases minimum} (environ 60 à 80 caractères) en respectant les consignes suivantes :
\begin{itemize}
\item utilise obligatoirement les trois connecteurs \zh{因为} (yīnwèi), \zh{所以} (suǒyǐ) et \zh{但是} (dànshì) ;
\item emploie au moins quatre mots du vocabulaire du thème : \zh{上网}, \zh{发信息}, \zh{聊天}, \zh{点赞}, \zh{网友}, \zh{手机}, \zh{网络} ;
\item présente au moins un avantage et un inconvénient des réseaux sociaux ;
\item termine par ton opinion personnelle avec \zh{我觉得}\ldots
\end{itemize}

\textbf{Critères d'évaluation :} correction des caractères (4 pts), grammaire et ordre des mots (4 pts), respect des consignes et connecteurs (4 pts), richesse du vocabulaire (4 pts), cohérence du texte (4 pts). Total : 20 pts.
\end{exercice}

\begin{exoresolu}[Modèle de production écrite commenté]
Rédige le sujet \zh{我和社交媒体} en suivant les consignes de l'exercice précédent.
\tcblower
\textbf{Modèle de réponse :}\\
\zh{我每天都用手机上网。因为网络很方便，所以我可以很快地知道新闻。我也喜欢给朋友发信息、聊天和点赞。但是，有些年轻人花太多时间上网，这影响学习。我有一个网友，他住在杜阿拉，我们常常聊天。我觉得，社交媒体很有用，但是我们不应该忘记学习和运动。}\\
(Wǒ měi tiān dōu yòng shǒujī shàng wǎng. Yīnwèi wǎngluò hěn fāngbiàn, suǒyǐ wǒ kěyǐ hěn kuài de zhīdào xīnwén. Wǒ yě xǐhuan gěi péngyou fā xìnxī, liáotiān hé diǎn zàn. Dànshì, yǒuxiē niánqīngrén huā tài duō shíjiān shàng wǎng, zhè yǐngxiǎng xuéxí. Wǒ yǒu yí ge wǎngyǒu, tā zhù zài Dù'ālā, wǒmen chángcháng liáotiān. Wǒ juéde, shèjiāo méitǐ hěn yǒuyòng, dànshì wǒmen bù yīnggāi wàngjì xuéxí hé yùndòng.)\\
« Chaque jour, je me connecte à Internet avec mon téléphone. Parce qu'Internet est très pratique, je peux connaître les nouvelles très vite. J'aime aussi envoyer des messages à mes amis, discuter et liker. Mais certains jeunes passent trop de temps en ligne, et cela nuit aux études. J'ai un ami d'Internet qui habite à Douala ; nous discutons souvent. Je pense que les réseaux sociaux sont très utiles, mais que nous ne devons pas oublier les études et le sport. »\\
\textbf{Commentaire :} le texte compte 6 phrases, utilise \zh{因为}\ldots\zh{所以}\ldots, \zh{但是} (deux fois), cinq mots du thème (\zh{上网}, \zh{发信息}, \zh{聊天}, \zh{点赞}, \zh{网友}, \zh{手机}, \zh{网络}), présente un avantage (connaître vite les nouvelles) et un inconvénient (perte de temps, études), et se termine par \zh{我觉得}.
\end{exoresolu}

\newpage

\coursec{Corrigés des exercices}

\coursec{Corrigés des exercices — Leçon 38 : Expression écrite — Origine du caractère 电 et radicaux ; caractères des réseaux sociaux}

\begin{corrige}[Exercice 1 — QCM : la bonne traduction]
1. \textbf{A} — \zh{点赞照片} (diǎn zàn zhàopiàn). Le verbe \zh{点赞} (« liker ») précède son complément d'objet \zh{照片} : ordre Verbe + Objet, comme en français. Les propositions B et C violent cet ordre.\\
2. \textbf{B} — \zh{发信息} (fā xìnxī). \zh{发} signifie « envoyer, émettre » ; \zh{收} (shōu) = « recevoir » (sens inverse) ; \zh{看} (kàn) = « regarder, lire ».\\
3. \textbf{B} — \zh{上网} (shàng wǎng), littéralement « monter sur le filet ». \zh{下网} (xià wǎng) signifie au contraire « se déconnecter » ; \zh{打网} n'existe pas dans ce sens.\\
4. \textbf{A} — \zh{电} représentait à l'origine un éclair zigzaguant ; la forme ancienne complète \zh{電} comportait la clé de la pluie \zh{雨}.\\
\textbf{Point de vigilance :} au Bac, ce type de QCM teste l'ordre Verbe + Objet et les verbes fréquents du thème (\zh{发}, \zh{收}, \zh{上}, \zh{点}). Apprends les couples verbe + nom par cœur : \zh{发信息}, \zh{上网}, \zh{打电话}, \zh{点赞}.
\end{corrige}

\begin{corrige}[Exercice 2 — Vrai ou faux justifié]
1. \textbf{Vrai.} \zh{话} se compose de la clé de la parole \zh{讠} (à gauche, 2 traits) et de la partie phonétique \zh{舌} (shé, « langue ») qui indique la prononciation.\\
2. \textbf{Faux.} \zh{申} se termine par une verticale droite qui traverse la case de part en part ; \zh{电} se termine par un crochet vertical recourbé vers la gauche (乚). C'est le dernier trait qui les différencie.\\
3. \textbf{Vrai.} \zh{网} (wǎng, « réseau ») + \zh{友} (yǒu, « ami ») = « ami du réseau », c'est-à-dire un ami connu sur Internet.\\
4. \textbf{Faux.} En chinois, le complément de temps se place \emph{avant} le verbe, jamais après : \zh{我每天上网。} (Wǒ měi tiān shàng wǎng. — « Je me connecte chaque jour. »). La phrase *\zh{我上网每天} est incorrecte.\\
\textbf{Point de vigilance :} la place du complément de temps (avant le verbe) est l'erreur la plus fréquente des francophones ; elle est systématiquement sanctionnée à l'examen.
\end{corrige}

\begin{corrige}[Exercice 3 — Vocabulaire : complète le tableau]
\begin{center}\begin{tabularx}{\linewidth}{|X|X|X|X|}\hline
\rowcolor{popblueL} Caractères & Pinyin & Nature & Traduction \\ \hline
\zh{电话} & diànhuà & nom & téléphone \\ \hline
\zh{网络} & wǎngluò & nom & réseau, Internet \\ \hline
\zh{发} & fā & verbe & envoyer, émettre \\ \hline
\zh{短信} & duǎnxìn & nom & SMS, texto (litt. « message court ») \\ \hline
\zh{聊天} & liáotiān & verbe & bavarder, discuter \\ \hline
\end{tabularx}\end{center}
\textbf{Remarque :} \zh{聊天} est un verbe séparable : on peut dire \zh{聊了一个小时天} (liáole yí ge xiǎoshí tiān — « discuter une heure »), la durée s'insérant entre les deux syllabes.
\end{corrige}

\begin{corrige}[Exercice 4 — Associe le pinyin au caractère]
\zh{电} — diàn — électricité ; \zh{网} — wǎng — filet, réseau ; \zh{信} — xìn — lettre, message ; \zh{打} — dǎ — composer, frapper ; \zh{话} — huà — parole, langue.\\
\textbf{Point de vigilance :} attention aux tons : \zh{打} est au 3\textsuperscript{e} ton (dǎ), \zh{话} au 4\textsuperscript{e} ton (huà). Une erreur de ton change le sens : \zh{dǎ} (frapper) / \zh{dà} (grand).
\end{corrige}

\begin{corrige}[Exercice 5 — Caractères proches : \zh{话} ou \zh{活} ? \zh{电} ou \zh{申} ?]
1. \zh{打电话。} (Dǎ diànhuà.) « Téléphoner. / Passer un coup de fil. » — \zh{电} : crochet final recourbé.\\
2. \zh{他说中国话说得很好。} (Tā shuō Zhōngguóhuà shuō de hěn hǎo.) « Il parle très bien le chinois. » — \zh{话} : clé de la parole, car il s'agit de langue. Remarque la structure : verbe + objet + verbe répété + \zh{得} + adjectif.\\
3. \zh{他在雅温得生活了十年。} (Tā zài Yǎwēndé shēnghuó le shí nián.) « Il a vécu dix ans à Yaoundé. » — \zh{活} : clé de l'eau \zh{氵}, verbe \zh{生活} (« vivre »).\\
4. \zh{请把你的申请表交给老师。} (Qǐng bǎ nǐ de shēnqǐngbiǎo jiāo gěi lǎoshī.) « Merci de remettre ton formulaire de demande au professeur. » — \zh{申} : verticale droite finale, dans \zh{申请} (« faire une demande »).\\
5. \zh{这个电影很有意思。} (Zhège diànyǐng hěn yǒu yìsi.) « Ce film est très intéressant. » — \zh{电} dans \zh{电影} (« cinéma », litt. « ombres électriques »).\\
\textbf{Point de vigilance :} avant d'écrire, identifie le sens : parole → \zh{讠} ; vie/eau → \zh{氵} ; électricité → crochet recourbé ; demande → verticale droite.
\end{corrige}

\begin{corrige}[Exercice 6 — Texte à trous : mots-outils et vocabulaire]
Texte complété :\\
\zh{因为网络很方便，我每天都用手机给朋友发信息。昨天我收到了一个网友的消息，但是我还没有回复。}\\
(Yīnwèi wǎngluò hěn fāngbiàn, wǒ měi tiān dōu yòng shǒujī gěi péngyou fā xìnxī. Zuótiān wǒ shōu dào le yí gè wǎngyǒu de xiāoxi, dànshì wǒ hái méiyǒu huífù.)\\
« Parce qu'Internet est très pratique, j'utilise chaque jour mon téléphone pour envoyer des messages à mes amis. Hier, j'ai reçu le message d'un ami d'Internet, mais je n'ai pas encore répondu. »\\
\textbf{Justifications :} \zh{因为} ouvre la proposition de cause ; \zh{每天} (complément de temps) se place avant le verbe \zh{用} ; \zh{信息} est le complément d'objet de \zh{发} ; \zh{了} après \zh{收到} marque l'action accomplie (hier) ; \zh{的} relie le possesseur \zh{网友} au nom \zh{消息} ; \zh{但是} introduit l'opposition finale.
\end{corrige}

\begin{corrige}[Exercice 7 — Remets les mots dans l'ordre]
1. \zh{我每天上两个小时网。} (Wǒ měi tiān shàng liǎng gè xiǎoshí wǎng.) « Je me connecte deux heures chaque jour. » — \zh{上网} est un verbe séparable : la durée \zh{两个小时} s'insère entre \zh{上} et \zh{网}.\\
2. \zh{因为微信很方便，所以我喜欢它。} (Yīnwèi Wēixìn hěn fāngbiàn, suǒyǐ wǒ xǐhuan tā.) « Parce que WeChat est très pratique, je l'aime beaucoup. » — en chinois, \zh{因为} et \zh{所以} s'emploient ensemble dans la même phrase (contrairement au français).\\
3. \zh{我给妈妈发了一条短信。} (Wǒ gěi māma fā le yì tiáo duǎnxìn.) « J'ai envoyé un SMS à maman. » — structure \zh{给} + personne + verbe ; classificateur des messages : \zh{条}.\\
4. \zh{我的网友住在杜阿拉。} (Wǒ de wǎngyǒu zhù zài Dù'ālā.) « Mon ami d'Internet habite à Douala. » — le complément de lieu introduit par \zh{在} se place après le verbe \zh{住}.\\
\textbf{Point de vigilance :} schéma de base à revérifier à chaque phrase : Sujet + (temps) + (\zh{给} + personne) + verbe + (了) + (classificateur) + objet.
\end{corrige}

\begin{corrige}[Exercice 8 — Écriture des caractères : traits et radicaux]
1. Ordre des traits (règles : de haut en bas, de gauche à droite, l'horizontal avant le vertical, l'extérieur avant l'intérieur, on ferme en dernier) :
\begin{itemize}
\item \zh{电} (5 traits) : 1. trait vertical (丨) ; 2. pli horizontal-crochet ((trait brisé)) formant le haut et le côté droit du cadre ; 3. trait horizontal (一) au milieu ; 4. trait horizontal (一) en bas ; 5. grand crochet vertical recourbé (乚, shù wān gōu) qui traverse et ferme.
\item \zh{话} (8 traits) : d'abord \zh{讠} à gauche (1. point ; 2. pli horizontal-crochet), puis \zh{舌} à droite (3. jeté ; 4. horizontale ; 5. verticale ; 6--8. la bouche \zh{口} : verticale, horizontale-crochet, horizontale de fermeture).
\item \zh{网} (6 traits) : 1. verticale gauche du cadre ; 2. horizontale-crochet ; 3--4. première croix intérieure (jeté, point) ; 5--6. deuxième croix (jeté, point). Le bas reste ouvert.
\item \zh{信} (9 traits) : \zh{亻} à gauche (1. jeté ; 2. verticale), puis \zh{言} à droite (3. point ; 4--6. trois horizontales ; 7--9. la bouche \zh{口}).
\item \zh{发} (5 traits) : 1. pli horizontal ; 2. jeté ; 3. pli horizontal-crochet ; 4. jeté long ; 5. point final.
\end{itemize}
2. Radicaux : \zh{电} : clé du champ \zh{田} (la forme ancienne \zh{電} se classait sous la clé de la pluie \zh{雨}) ; \zh{话} : clé de la parole \zh{讠} ; \zh{网} : clé du filet \zh{网} (variante \zh{罒}) ; \zh{信} : clé de l'homme \zh{亻} ; \zh{发} : clé de la main droite \zh{又}.\\
3. Même radical que \zh{话} (clé de la parole \zh{讠}) : \zh{说} (shuō, parler), \zh{语} (yǔ, langue), \zh{请} (qǐng, s'il vous plaît), \zh{谢} (xiè, remercier).\\
\textbf{Point de vigilance :} à l'examen, un trait mal placé (crochet de \zh{电} tracé droit → \zh{申}) rend le caractère faux : soigne le dernier trait.
\end{corrige}

\begin{corrige}[Exercice 9 — Compréhension de l'écrit et questions de langue]
\textbf{Questions de compréhension :}\\
1. Trois activités parmi : se connecter à Internet avec le téléphone (\zh{上网}), envoyer des messages (\zh{发信息}), regarder des vidéos (\zh{看视频}), liker leurs amis (\zh{点赞}).\\
2. L'avantage : Internet est pratique, tout le monde peut connaître les nouvelles très vite (\zh{网络很方便，大家可以很快地知道新闻}).\\
3. Le danger : certains jeunes passent trop de temps en ligne chaque jour, ce qui nuit à leurs études.\\
4. L'opinion personnelle : \zh{我觉得，我们应该好好利用网络，但是不应该忘记学习和运动。} (« Je pense que nous devons bien utiliser Internet, mais que nous ne devons pas oublier les études et le sport. »)\\
\textbf{Questions de langue :}\\
1. \zh{但是} (dànshì, « mais ») : introduit une opposition entre avantage et inconvénient ; \zh{也} (yě, « aussi ») : ajoute une activité à la liste ; on accepte aussi \zh{这} (zhè, « cela ») qui reprend la phrase précédente comme sujet.\\
2. \zh{影} : clé \zh{彡} (clé des poils / traits décoratifs) ; \zh{信} : clé de l'homme \zh{亻}.\\
3. \zh{有些年轻人每天花太多时间上网。} : « Certains jeunes passent trop de temps chaque jour à se connecter à Internet. » (\zh{花} + temps = « passer / dépenser du temps »).\\
\textbf{Barème indicatif (10 pts) :} compréhension 4 × 1,5 pt = 6 pts ; langue 1 + 1 + 2 = 4 pts.
\end{corrige}

\begin{corrige}[Exercice 10 — Thème et version]
\textbf{A. Thème :}\\
1. \zh{因为网络很方便，所以我每天和朋友聊天，但是我不浪费时间。} (Yīnwèi wǎngluò hěn fāngbiàn, suǒyǐ wǒ měi tiān hé péngyou liáotiān, dànshì wǒ bù làngfèi shíjiān.)\\
2. \zh{我姐姐给她杜阿拉的朋友发信息。} (Wǒ jiějie gěi tā Dù'ālā de péngyou fā xìnxī.) — \zh{给} + personne avant le verbe ; \zh{的} entre le complément de lieu et \zh{朋友}.\\
3. \zh{我很喜欢给同学的照片点赞。} (Wǒ hěn xǐhuan gěi tóngxué de zhàopiàn diǎn zàn.)\\
\textbf{B. Version :}\\
1. « Mon ami d'Internet passe trois heures par jour en ligne ; je pense que cela nuit aux études. »\\
2. « Hier, j'ai envoyé un SMS au professeur. »\\
3. « Parce que le téléphone portable est très pratique, tout le monde aime discuter avec. »\\
\textbf{Points de vigilance :} en thème, ne jamais traduire mot à mot : « perdre du temps » = \zh{浪费时间} ; « liker les photos de quelqu'un » = \zh{给} + personne + \zh{的照片点赞}. En version, \zh{花} + durée se traduit par « passer (du temps) », non par « fleur ».
\end{corrige}

\begin{corrige}[Exercice 11 — Expression écrite : sujet de type Baccalauréat]
\textbf{Proposition de rédaction modèle :}\\
\zh{我每天都用手机上网。因为网络很方便，所以我可以很快地知道新闻。我也喜欢给朋友发信息、聊天和点赞。但是，有些年轻人花太多时间上网，这影响学习。我有一个网友，他住在杜阿拉，我们常常聊天。我觉得，社交媒体很有用，但是我们不应该忘记学习和运动。}\\
(Wǒ měi tiān dōu yòng shǒujī shàng wǎng. Yīnwèi wǎngluò hěn fāngbiàn, suǒyǐ wǒ kěyǐ hěn kuài de zhīdào xīnwén. Wǒ yě xǐhuan gěi péngyou fā xìnxī, liáotiān hé diǎn zàn. Dànshì, yǒuxiē niánqīngrén huā tài duō shíjiān shàng wǎng, zhè yǐngxiǎng xuéxí. Wǒ yǒu yí gè wǎngyǒu, tā zhù zài Dù'ālā, wǒmen chángcháng liáotiān. Wǒ juéde, shèjiāo méitǐ hěn yǒuyòng, dànshì wǒmen bù yīnggāi wàngjì xuéxí hé yùndòng.)\\
« Chaque jour, je me connecte à Internet avec mon téléphone. Parce qu'Internet est très pratique, je peux connaître les nouvelles très vite. J'aime aussi envoyer des messages à mes amis, discuter et liker. Mais certains jeunes passent trop de temps en ligne, et cela nuit aux études. J'ai un ami d'Internet qui habite à Douala ; nous discutons souvent. Je pense que les réseaux sociaux sont très utiles, mais que nous ne devons pas oublier les études et le sport. »\\
\textbf{Barème indicatif (20 pts) :} correction des caractères (4 pts : un caractère faux = $-$0,5 pt) ; grammaire et ordre des mots (4 pts : place du temps, de \zh{了}, classificateurs) ; respect des consignes et connecteurs (4 pts : \zh{因为}/\zh{所以}/\zh{但是} présents et correctement employés) ; richesse du vocabulaire (4 pts : au moins 4 mots du thème) ; cohérence et plan (4 pts : situation → avantage → inconvénient → opinion).\\
\textbf{Points de vigilance :} vérifier le dernier trait de \zh{电} (crochet, pas verticale droite) ; ne pas confondre \zh{话}/\zh{活} ; placer \zh{每天} avant le verbe ; terminer impérativement par \zh{我觉得}\ldots ; compter ses phrases (6 minimum) avant de rendre la copie.
\end{corrige}

\newpage

\coursec{Fiche bilan}

\begin{popfigure}
\begin{center}\tcbox[colback=popink,colframe=popink,coltext=white,arc=9pt,boxsep=4pt,left=6pt,right=6pt]{\bfseries\small Leçon 38 Caractères \& Réseaux Sociaux}\end{center}
\vspace{-2pt}
\begin{tcbraster}[raster columns=3,raster equal height=rows,raster column skip=6pt,raster row skip=6pt,enhanced,blankest]
\begin{tcolorbox}[enhanced,colback=poporange,colframe=poporange,coltext=white,boxrule=0.9pt,arc=3pt,left=4pt,right=4pt,top=3pt,bottom=3pt,boxsep=1pt,halign=left]\bfseries\scriptsize Origine de \zh{电} éclair $\to$ électricité\end{tcolorbox}
\begin{tcolorbox}[enhanced,colback=popgreen,colframe=popgreen,coltext=white,boxrule=0.9pt,arc=3pt,left=4pt,right=4pt,top=3pt,bottom=3pt,boxsep=1pt,halign=left]\bfseries\scriptsize Radicaux \& Écriture 田 雨 糸 讠 扌\end{tcolorbox}
\begin{tcolorbox}[enhanced,colback=poppurple,colframe=poppurple,coltext=white,boxrule=0.9pt,arc=3pt,left=4pt,right=4pt,top=3pt,bottom=3pt,boxsep=1pt,halign=left]\bfseries\scriptsize Vocab. Réseaux Sociaux 微信 微博 发 评论 点赞\end{tcolorbox}
\begin{tcolorbox}[enhanced,colback=poppink,colframe=poppink,coltext=white,boxrule=0.9pt,arc=3pt,left=4pt,right=4pt,top=3pt,bottom=3pt,boxsep=1pt,halign=left]\bfseries\scriptsize Connecteurs Logiques 首先 而且 但是 因为\ldots 所以\end{tcolorbox}
\begin{tcolorbox}[enhanced,colback=popgold,colframe=popgold,coltext=white,boxrule=0.9pt,arc=3pt,left=4pt,right=4pt,top=3pt,bottom=3pt,boxsep=1pt,halign=left]\bfseries\scriptsize Modèle de Texte Sujet + Temps + Lieu + Verbe\end{tcolorbox}
\begin{tcolorbox}[enhanced,colback=popink,colframe=popink,coltext=white,boxrule=0.9pt,arc=3pt,left=4pt,right=4pt,top=3pt,bottom=3pt,boxsep=1pt,halign=left]\bfseries\scriptsize Thème / Version Ordre mots + Particules\end{tcolorbox}
\begin{tcolorbox}[enhanced,colback=popdark,colframe=popdark,coltext=white,boxrule=0.9pt,arc=3pt,left=4pt,right=4pt,top=3pt,bottom=3pt,boxsep=1pt,halign=left]\bfseries\scriptsize Prépa Examen Caractères Grammaire Cohérence\end{tcolorbox}
\end{tcbraster}

\legende{Carte mentale de la leçon 38 : points clés à maîtriser}
\end{popfigure}

\begin{bilan}

\end{bilan}

\courssub{Lexique clé — Réseaux sociaux \& Caractères apparentés}

\begin{center}
\begin{tabularx}{\linewidth}{|X|X|X|X|}
\hline
\rowcolor{popblueL} \textbf{Caractères} & \textbf{Pinyin} & \textbf{Nature} & \textbf{Traduction} \\ \hline
\zh{电} & diàn & nom / verbe & électricité ; télégraphique, électronique \\ \hline
\zh{电脑} & diànnǎo & nom & ordinateur \\ \hline
\zh{电影} & diànyǐng & nom & film, cinéma \\ \hline
\zh{电视} & diànshì & nom & télévision \\ \hline
\zh{手机} & shǒujī & nom & téléphone portable \\ \hline
\zh{网络} & wǎngluò & nom & réseau, Internet \\ \hline
\zh{微信} & Wēixìn & nom propre & WeChat \\ \hline
\zh{微博} & Wēibó & nom propre & Weibo (Twitter chinois) \\ \hline
\zh{发} & fā & verbe & envoyer, publier (un message) \\ \hline
\zh{评论} & pínglùn & verbe / nom & commenter / commentaire \\ \hline
\zh{点赞} & diǎnzàn & verbe & liker, mettre un « j'aime » \\ \hline
\zh{转发} & zhuǎnfā & verbe & partager, retweeter \\ \hline
\zh{关注} & guānzhù & verbe & suivre (un compte) \\ \hline
\zh{视频} & shìpín & nom & vidéo \\ \hline
\zh{直播} & zhíbō & verbe / nom & diffuser en direct / direct \\ \hline
\zh{网红} & wǎnghóng & nom & influenceur, star du web \\ \hline
\zh{流量} & liúliàng & nom & trafic, audience (data) \\ \hline
\end{tabularx}
\end{center}

\courssub{Règles de grammaire \& structures à connaître par cœur}

\begin{center}
\begin{tabularx}{\linewidth}{|X|X|X|}
\hline
\rowcolor{popgreenL} \textbf{Structure} & \textbf{Exemple (Caractères + Pinyin)} & \textbf{Traduction} \\ \hline
\zh{在 + Lieu/Plateforme + 动词} & \zh{我在微信上发照片。}\\\zh{(Wǒ zài Wēixìn shàng fā zhàopiàn.)} & Je publie des photos sur WeChat. \\ \hline
\zh{给 + Destinataire + 动词} & \zh{我给朋友发信息。}\\\zh{(Wǒ gěi péngyou fā xìnxī.)} & J'envoie un message à mon ami. \\ \hline
\zh{把 + Objet + 动词 + 结果/方向} & \zh{他把视频转发了。}\\\zh{(Tā bǎ shìpín zhuǎnfā le.)} & Il a partagé la vidéo. \\ \hline
\zh{因为\ldots 所以\ldots} & \zh{因为网络快，所以我看直播。}\\\zh{(Yīnwèi wǎngluò kuài, suǒyǐ wǒ kàn zhíbō.)} & Comme le net est rapide, je regarde le direct. \\ \hline
\zh{不但\ldots 而且\ldots} & \zh{这位网红不但唱歌好，而且跳舞也行。}\\\zh{(Zhè wèi wǎnghóng búdàn chànggē hǎo, érqiě tiàowǔ yě xíng.)} & Cet influenceur non seulement chante bien, mais danse aussi. \\ \hline
\zh{动词 + 了 (action achevée)} & \zh{我已经关注了他。}\\\zh{(Wǒ yǐjīng guānzhù le tā.)} & Je l'ai déjà suivi. \\ \hline
\zh{正在 + 动词 + 呢 (action en cours)} & \zh{他正在看评论呢。}\\\zh{(Tā zhèngzài kàn pínglùn ne.)} & Il est en train de lire les commentaires. \\ \hline
\end{tabularx}
\end{center}

\courssub{Méthodes express}

\begin{itemize}
  \item \textbf{Écriture caractères :} Respecter l'ordre des traits (haut→bas, gauche→droite, horizontal avant vertical, fermer en dernier). Identifier le radical (clé) pour classer/mémoriser : \zh{田} (champ) pour \zh{男}, \zh{雨} (pluie) pour \zh{电} (forme traditionnelle \zh{電}), \zh{糸} (fil) pour \zh{络}, \zh{讠} (parole) pour \zh{评}, \zh{扌} (main) pour \zh{发} \zh{点} \zh{转}.
  \item \textbf{Origine de \zh{电} :} Forme traditionnelle \zh{電} = \zh{雨} (nuage/pluie) + \zh{田} (champ) + \zh{ㄅ} (éclair jaillissant) $\to$ « éclair sous la pluie ». Forme simplifiée \zh{电} : radical \zh{田} + trait d'éclair.
  \item \textbf{Thème (Fr $\to$ Ch) :} 1. Identifier le sujet/temps/lieu. 2. Choisir la structure (把/被/在/给). 3. Placer les compléments (Temps $\to$ Lieu $\to$ Manière $\to$ Verbe). 4. Vérifier particules (\zh{了},\zh{过},\zh{着}) et classificateurs.
  \item \textbf{Version (Ch $\to$ Fr) :} 1. Segmenter la phrase (mots). 2. Identifier le noyau verbal. 3. Reconstruire l'ordre français (Suj + Verbe + Compléments). 4. Adapter le temps (passé/présent/futur) selon \zh{了}/\zh{过}/\zh{将}/\zh{正在}.
  \item \textbf{Rédaction :} Utiliser connecteurs (\zh{首先},\zh{其次},\zh{最后},\zh{但是},\zh{因此}) pour structurer : Introduction (sujet) $\to$ Développement (arguments/exemples) $\to$ Conclusion (avis personnel).
\end{itemize}



\begin{aretenir}[Checklist avant le Bac — Leçon 38]
$\square$ Écrire correctement \zh{电} (5 traits) et expliquer son origine (éclair + pluie) ainsi que la différence avec \zh{電}.
$\square$ Identifier et écrire les radicaux clés du thème : \zh{田} \zh{雨} \zh{糸} \zh{讠} \zh{扌} \zh{目} \zh{手}.
$\square$ Différencier \zh{发} (fā, envoyer) et \zh{发} (fà, cheveux) / \zh{发} (fā, partir) — homophones.
$\square$ Maîtriser le vocabulaire réseaux sociaux : \zh{微信} \zh{微博} \zh{点赞} \zh{评论} \zh{转发} \zh{关注} \zh{直播} \zh{网红} \zh{流量}.
$\square$ Construire une phrase avec \zh{在 + Plateforme + 上 + Verbe} (ex: \zh{在微博上发视频}).
$\square$ Utiliser la structure \zh{把} pour action sur objet défini (ex: \zh{把手机关了}).
$\square$ Relier deux idées avec \zh{因为\ldots 所以\ldots} et \zh{不但\ldots 而且\ldots}.
$\square$ Traduire un court paragraphe Fr $\to$ Ch (Thème) en respectant l'ordre T-L-M-V et les particules d'aspect.
$\square$ Traduire un court paragraphe Ch $\to$ Fr (Version) en segmentant correctement et rendant les aspects verbaux.
$\square$ Rédiger un message type « publication + hashtags » ou « commentaire » en chinois (50-80 caractères) avec connecteurs logiques.
$\square$ Reconnaître les caractères proches / pièges : \zh{电/田/由/甲} ; \zh{网/罗/罔} ; \zh{评/平} ; \zh{论/伦} ; \zh{赞/赃}.
$\square$ Connaître 2 différences socioculturelles : écosystème « Super App » (WeChat) vs apps spécialisées FR ; paiement mobile QR code omniprésent en Chine.
\end{aretenir}

\findecours

\end{document}
